\documentclass[11pt]{article}
\usepackage[T1]{fontenc}
\usepackage{lmodern,microtype,amsmath,amssymb,amsthm,mathtools}
\usepackage[a4paper,margin=25mm]{geometry}
\usepackage{booktabs,longtable,tabularx,array,enumitem,xurl}
\usepackage[hidelinks]{hyperref}
\hypersetup{pdftitle={Nonequilibrium Calibration and Faithful Records in Autonomous Effective Measurement Models},pdfauthor={Jeremy Rodgers}}
\setlength{\emergencystretch}{3em}
\newtheorem{theorem}{Theorem}[section]
\newtheorem{lemma}[theorem]{Lemma}
\newtheorem{proposition}[theorem]{Proposition}
\newtheorem{corollary}[theorem]{Corollary}
\theoremstyle{definition}\newtheorem{definition}[theorem]{Definition}
\newtheorem{assumption}[theorem]{Assumption}
\theoremstyle{remark}\newtheorem{remark}[theorem]{Remark}
\newcommand{\pr}{}\newcommand{\cert}{}
\newcommand{\norm}[1]{\left\lVert#1\right\rVert}
\newcommand{\Prob}{\mathbb P}\newcommand{\E}{\mathbb E}
\newcommand{\TV}{\operatorname{TV}}\newcommand{\Var}{\operatorname{Var}}
\newcommand{\one}{\mathbf1}\newcommand{\R}{\mathbb R}\newcommand{\T}{\mathbb T}
\newcommand{\dd}{\,\mathrm d}\newcommand{\Imv}{\operatorname{Im}}\newcommand{\ii}{\mathrm i}
\title{Nonequilibrium Calibration and Faithful Records\\in Autonomous Effective Measurement Models}
\author{Jeremy Rodgers\thanks{Website: \href{https://everythingequation.com}{everythingequation.com}. Manuscript DOI: \url{https://doi.org/10.5281/zenodo.23131081}}\\\small Independent Researcher}
\date{4 October 2026}
\begin{document}\maketitle
\begin{abstract}
A calibrated final pointer distribution does not by itself establish that a
record is faithful to an earlier event. We consolidate two finite autonomous
Schr\"odinger--guidance constructions in which both questions have quantitative
answers under explicitly restricted nonequilibrium initial laws. The first uses
periodic focusing, a finite recoiling reference rotor and a delayed compensated
oscillator. Its record error is at most $0.00443484008607784720002304$ and its
failure to copy and hold its actual earlier label is at most
$0.00232725479707784720002304$. The second starts with a small radial Gaussian,
includes autonomous preparation and activation, and contains its writer from
the original entrance. Its joint earlier-label and entire-record law has total
variation distance below $0.006086770113$ from a constant Born-labelled path;
its actual-label copy-and-hold failure is below $0.000552421956$.
The proofs combine bounded-variation calibration, exact recoil reduction,
conditional-rank currents, finite-clock derivative estimates and absolute
holding flux. The two Hamiltonians and admissible law classes remain separate.
Neither construction assumes full initial equilibrium, but both retain
substantive readiness and statistical restrictions. Neither establishes a
finite local source of every prescribed interaction or a realization in
established matter. The appendices supply the specialized derivative, current and numerical
estimates supporting these bounds. External specialist verification remains
an open step.
\end{abstract}
\clearpage
\begingroup\small
\makeatletter\renewcommand{\l@section}{\@dottedtocline{1}{0em}{2.3em}}\makeatother
\tableofcontents
\endgroup\clearpage

\section{The problem and the contribution}
An account of a measurement should explain not only the frequency of its final
display but also what that display records. A pointer may have a suitable
one-time distribution while subsequently changing sides, or while failing to
copy the system's earlier position. We therefore specify an earlier
configuration event and a whole holding interval. The statistical conclusion
concerns that joint observable.

The two constructions below separate three mathematical requirements:
calibration of the earlier event, physical copying under the same dynamics,
and retention at every time of a finite interval. Their error estimates are
uniform over stated input and initial-law classes. They are not claims of
universal relaxation from an arbitrary configuration law.

Quantum equilibrium, subsystem statistics and Bohmian measurement theory
have a substantial established literature \cite{DGZ,Operators}. The method of
arbitrary functions has also been applied to quantum probabilities: Bonds
and collaborators study a random perturbation of a double-well model in
long-time and small-$\hbar$ limits \cite{Bonds}. Our finite apparatus statements
instead concern restricted configuration ensembles and explicit recorded
histories. Inverse engineering of harmonic transport is established
\cite{Tor}; it is used here with an autonomous quantum clock and a separate
historical-current estimate. Coarse-grained relaxation studies
\cite{ValentiniWestman} address a different statistical question. We make no
claim that any of these general ideas originates in this work.

The contribution being consolidated is the particular assumption-to-record
chain, including its finite parameters and its current-sensitive error
accounting. It consolidates the author's earlier research; the numerical constructions
are not presented as new discoveries in this revision. Its relationship to the author's
equilibrium and hybrid pilot-medium papers is comparative: those papers use
different initial laws or dynamics and do not supply missing premises here
\cite{RodgersMonograph,RodgersMassive,RodgersPilot}.

\subsection{What is and is not a physical claim}
Throughout, the wave obeys a specified Schr\"odinger equation and actual
configurations obey its guidance law. These are the constitutive assumptions.
The potentials and coherent entrance stock are supplied mathematically.
``Autonomous'' means the complete displayed Hamiltonian has no externally
switched laboratory-time parameter. It does not mean that its coefficients,
supports and stock have been manufactured with established interactions.
The intrinsic-circle model has an explicitly finite recoil coordinate.
The radial model is an exact two-body $s$-wave reduction, but its clock and
pointer remain effective scalar coordinates. Adding a field or a material
support changes the model and requires a new same-model history proof.

\section{Probability and complete record observables}
For a spinor $\Psi$, put $\rho=\Psi^\dagger\Psi$ and
$j_k=(\hbar/m_k)\Im(\Psi^\dagger\partial_k\Psi)$ in physical coordinates.
Internal projectors label wave components; they are not independently sampled
configuration bits. The wave measure and the actual initial law $\nu_0$ are
distinct. A bound $\nu_0\le C\rho_0\dd q$ transports under the same conservative
equivariant flow, and implies $\E_{\nu_0}F\le C\E_{\rho_0}F$ for every
nonnegative path functional $F$. It does not establish equilibrium or
domination by an analytical comparison wave.

We use $\TV(P,Q)=\sup_A|P(A)-Q(A)|$. Fix an earlier label
$L\in\{0,1,\dagger\}$, a compact holding interval $I$, and disjoint closed
pointer regions $R_0,R_1$. For a continuous pointer trajectory define
$K_i=\{t\in I:Z_t\in R_i\}$, allowing the empty compact set, and
\[
 Y=(L,K_0,K_1),\qquad
 \iota(0)=(0,I,\varnothing),\quad
 \iota(1)=(1,\varnothing,I).
\]
The random compact sets give a measurable representation of the entire
symbolic record; no unwarranted c\`adl\`ag assumption on a thresholded path
is needed. For example, avoidance of a compact time set $J$ is the condition
$\min_{t\in J}\operatorname{dist}(Z_t,R_i)>0$, a measurable condition on
continuous paths. Countable rational interval approximations generate the
corresponding compact-set sigma-field. A visit outside both record regions,
however brief, counts as failure. Undefined $L$ also counts as failure.

\begin{lemma}[Direct ideal-record composition]\label{lem:joint}
Let $L_*$ be a binary reference on the same coupling. Suppose $E,A,B$ are
events such that outside $E\cup A$, $L=L_*$, and outside $E\cup B$ the
entire actual record is the constant record of the defined label $L$.
If $\TV(\mathcal L(L_*),\operatorname{Bernoulli}(p))\le e$, then
\[
 \TV(\mathcal L(Y),(\iota)_*\operatorname{Bernoulli}(p))
 \le \Prob(E)+\Prob(A)+\Prob(B)+e.
\]
\end{lemma}
\begin{proof}
On the complement of the three events, $Y=\iota(L_*)$. The coupling
inequality bounds the distance to $\mathcal L(\iota(L_*))$ by the union
probability; pushforward contracts the remaining distance. This proves the
claim, including undefined outcomes. Shared events are paid once because
their identity is established, not by subtracting unrelated upper bounds.
\end{proof}

A comparison to an imperfect reference's entire path is a different
statement: its own history failure generally must be paid as well.
Likewise, conditioning on a rare successful event may amplify error.
If $\TV(P,Q)\le\epsilon$ and $P(E)>\epsilon$, then the elementary
normalization bound is $\TV(P(\cdot\mid E),Q(\cdot\mid E))
\le\min(1,2\epsilon/P(E))$. No conditional conclusion is asserted for
zero-probability branches.

\subsection{Why absolute current is indispensable}
For real even normalized functions $g,h$, the wave
$\Psi(x,y)=g(x)h(y)e^{ikxy}$ has
$j_x(0,y)=(\hbar k/m)y g(0)^2h(y)^2$. Its signed hidden-coordinate
integral vanishes, whereas $\int |j_x(0,y)|\dd y>0$ when $k\ne0$.
Thus averaging before taking the modulus can erase the quantity that
bounds crossings. The record estimates below take the full-configuration
absolute value first. This example diagnoses an invalid inference, not
a no-go theorem for either apparatus.

\part{A periodic detector with finite recoil}
The first construction uses units $\hbar=1$. The active rotor/pointer/clock
coordinates and their law are specified afresh in this part. They are not
identified with the physical radius and clock of Part II.
% Integrated section: periodic_model
\section{All moving coordinates and finite parameters}
Let $x,R\in\T_2$ be a detector rotor and a new reference rotor, and set $X=x-R\pmod2$. Let $Z,S,C\in\R$, and retain the original 255 differential source coordinates $u\in\one^\perp\subset\R^{256}$. The unknown normalized qubit/reference vector is $\psi$, with exactly conserved projectors $P_0+P_1=I$.
Set $p=\norm{P_1\psi}^2$. The target projector $P_1^{\rm tar}$ is an
orthogonal projector with $\norm{P_1-P_1^{\rm tar}}\le10^{-4}$,
and $p_{\rm tar}=\norm{P_1^{\rm tar}\psi}^2$. Both act as the identity
on any admitted inaccessible reference factor. That reference is an internal
Hilbert-space factor and introduces no additional guidance coordinates. Thus $|p-p_{\rm tar}|\le10^{-4}$.
The periodic classifier uses the centred cell coordinate
$u=((X-o)/h)\bmod1\in[-1/2,1/2)$: it is zero for $|u|\le1/4$
and one otherwise, with a fixed boundary convention.


Set
\[
 N=256,\quad h=1/128,\quad K=2^{38},\quad
 \eta_t=2^{-16},\quad\eta_p=2^{-17}.
\]
Let $m= m_xm_R/(m_x+m_R)=\alpha K/h^2$. Independently choosing
\[
 m_x,m_R\in[1.998,2.002]K/h^2
\]
ensures $\alpha\in[.999,1.001]$. Nominally $m=2^{52}$ and $m_x=m_R=2^{53}$. The other masses are
\[
 M_s=10^{18},\quad M\in[.99,1.01]10^{18},\quad
 \mu\in[.00999,.01001],\quad M_C=10^8.
\]
The retained source has $\kappa_s=2.5\,10^{43}$ and frequencies $\Omega_k=10^{13}\sin(\pi k/256)$. It is a spectator Hamiltonian, not a field generator in this redesign.

Let $S_5(z)=10z^3-15z^4+6z^5$. The even period-one $U_{\eta_p}$ equals $z^2$ for $|z|\le1/2-\eta_p$ and has positive-cap derivative
\[
 U'_{\eta_p}(z)=2z\left[1-S_5((z-1/2+\eta_p)/\eta_p)\right].
\]
It is $C^3$, nonnegative, and bounded by $1/4$. Define
\[
 V_i(X)=\frac{g_iK}{2}U_{\eta_p}\left(\frac{X-o}{h}-\frac i2-d_i\right),
 \quad g_i\in[.999,1.001],\quad |d_i|\le.01.
\]
The common origin $o$ is arbitrary. The relative sector offsets and gains can vary independently in these intervals; arbitrary nonperiodic site defects are not admitted.

For $B_9(a)=126a^5-420a^6+540a^7-315a^8+70a^9$ on $[0,1]$, with constant extensions, put
\[
 b(a)=L B_9((a-a_0)/w_b),\quad y=Z-z_0,
\]
\[
 W_0(a,y)=\frac{y^2}{2\mu},\qquad
 W_1(a,y)=\frac{(y-b(a))^2}{2\mu}-\mu b''(a)(y-b(a)/2).
\]
The finite writer family is
\[
 \begin{gathered}
 L\in[23.99,24.01],\quad |z_0|\le.001,\quad v\in[127,129],\quad
 \lambda\in[.99999,1.00001],\\
 a_0\in[1.2299,1.2301],\qquad w_b\in[.03999,.04001].
 \end{gathered}
\]
Its launch momentum is $Mv\lambda$. Spring coefficient $1/\mu$ and the scalar compensation are matched to the selected $\mu$. General symmetry-breaking or independently mismatched compensation errors are outside this theorem.

\section{The ordinary autonomous Hamiltonian}
\begin{equation}\label{pm:eq:H}
\boxed{\begin{aligned}
 H_{\mathrm{per}}={}&\frac{p_x^2}{2m_x}+\frac{p_R^2}{2m_R}
 +\frac{p_Z^2}{2\mu}+\frac{p_S^2}{2M}
 +H_s(u)+\frac{p_C^2}{2M_C}+50000\\
 &+\sum_{i=0}^1P_i[V_i(x-R)+W_i(S/v,Z-z_0)],\\
 H_s={}&\frac{p_u^2}{2M_s}+\frac{\kappa_s}{2}\sum_j(u_{j+1}-u_j)^2.
\end{aligned}}
\end{equation}
No kinetic coefficient is switched off. No external time occurs in this expression. Every microscopic kinetic term shown is positive and quadratic; the relative reduced mass is a derived constant.

\begin{proposition}[Operator realization]
The finite operator \eqref{pm:eq:H} has a self-adjoint semibounded form realization, with ordinary configuration currents $j_a=m_a^{-1}\Imv(\Psi^\dagger\partial_a\Psi)$. The shift makes it nonnegative. The added reference body is finite, not an infinitely massive frame.
\end{proposition}
\begin{proof}
Use the free-rotor/clock and positive oscillator form. The angular potentials are bounded. Complete the pointer square:
\[
 W_1=\frac{[y-b-\mu^2b'']^2}{2\mu}
 -\frac{\mu bb''}{2}-\frac{\mu^3(b'')^2}{2}.
\]
The stated rectangle gives $|b|<25$, $|b''|<169000$, and a negative part below $43000$. The affine terms are infinitesimally oscillator-form bounded. The form construction and ordinary Green identity give the assertion. Finite-horizon conservative currents are assumed on the usual differentiated domains; the initial active-law domination transfers their wave-null exceptions. The passive-coordinate issue is handled exactly in the comparison companion, not by an arbitrary-clock almost-sure argument.
\end{proof}

\section{Simple initial waves; independent actual-law data}
Let $A_c(s)=1$ on $|s|\le.3$, $A_c(s)=1-B_9(5(|s|-.3))$ until $.5$, and zero outside. Set $Z_c=176818/230945$ and $\varphi=A_c/\sqrt{Z_c}$. The supplied wave is
\[
 \Psi_0=\frac{\psi}{2}\,\zeta(u)\gamma(Z-z_0)
 e^{iMv\lambda S}\varphi(S)\,
 e^{iM_CC}\varphi(C+1),
 \quad \gamma(z)=\pi^{-1/4}e^{-z^2/2}.
\]
The factor $1/2$ is the product of two flat rotor waves. No shaped rotor wave is supplied. $\zeta$ is the retained harmonic source ground wave.

In relative coordinates, the actual law is
\[
 f_0(X,Z,S)g(u\mid X,Z,S)\,
 \kappa_R(dR\mid X,Z,S,u)\,
 \kappa_C(dC\mid X,Z,S,u,R).
\]
Here $f_0$ is a nonnegative density with integral one. Define
$V_j=|D_j f_0|(\R^3)$ as the distributional total variation of its
zero extension, for $j=X,Z,S$, including boundary jumps, and
$B=V_X+V_Z+V_S$. The admitted same-support total-variation neighbourhood
retains absolute continuity of the active configuration marginal,
as required for the almost-sure flow; allowed singular passive
conditionals retain their separately stated supports.

The original class is retained: support $[-1,1]^2\times[-1/4,1/4]$, $V_X,V_Z\le2$, $V_S\le8$, $B\le12$, $f_0\le2$, and $0\le g\le2|\zeta|^2$ normalized. $\kappa_R$ is any normalized law on the reference circle; $\kappa_C$ is any normalized law on $[-1.25,-.75]$. Neither must be Born, independent, or absolutely continuous. Taking $R_0=0$ embeds the original detector positions without even a common translation. An independent actual orientation law can also be appended to an original bare-position ensemble. Its translated mixture preserves the density cap, $V_Z,V_S$, circular rotor-marginal variation at most two, and rotor-marginal density at most one; its mixed source conditional preserves the source domination. Zero-extension rotor variation and the raw total $B$ need not remain unchanged. The independent-orientation extension is a separate corollary because the active proof uses precisely the circular variation, marginal bound and density cap. This is not permission to correlate $R$ arbitrarily with bare $x$ while keeping the relative law unchanged. For example, $R=x$ would collapse the relative seed and is outside the declared relative-density class. The same-support $10^{-4}$ actual-law neighbourhood is charged once.

% Integrated section: periodic_auxiliary
\section{Exact relative-coordinate reduction with finite recoil}
On the subspace $\Psi(x,R,\ldots)=2^{-1/2}\Phi(x-R,\ldots)$,
\[
 \left(\frac{p_x^2}{2m_x}+\frac{p_R^2}{2m_R}\right)\Psi
 =-\frac1{2m}\partial_X^2\Psi,
 \qquad\frac1m=\frac1{m_x}+\frac1{m_R}.
\]
Both the Hamiltonian and the product of flat initial rotor waves preserve this subspace. No physical mixed kinetic term has been introduced: these are the restrictions of two constant ordinary kinetic operators.

\begin{theorem}[Finite recoil and spectator independence]\label{pa:thm:recoil}
The complete wave has the exact form
\[
 \Psi_t=\frac{\zeta(u)e^{-i(E_s+50000)t}\chi_{C,t}(C)}{\sqrt2}
 \sum_i(P_i\psi)r_i(X,t)\Xi_i(S,Z,t).
\]
Its active relative velocity, pointer velocity, and writer-clock velocity are independent of the instantaneous $R,u,C$. Moreover
\[
 \dot x=\frac m{m_x}v_X,\qquad
 \dot R=-\frac m{m_R}v_X,\qquad \dot x-\dot R=v_X.
\]
The old differential source positions are stationary under their nodeless ground wave, with their Hamiltonian and phase retained. Their arbitrary normalized actual conditional law is not used as an equilibrium reference.
\end{theorem}
\begin{proof}
The differential identity gives an invariant wave subspace. Every $u$ and $C$ term commutes with the active Hamiltonian. Differentiating the exact factorization gives the displayed ordinary currents; division cancels spectator densities wherever defined. The source wave is positive and has a coordinate-independent phase. The free $C$ tube proved below is positive at every allowed actual clock path. Active exceptions are wave-null and, by $f_0\le16q_0$, actual-null. These exceptions depend only on active initial data, so singular conditional $R,C$ laws cannot select them at otherwise good active positions.
\end{proof}

In particular all active event probabilities equal those calculated from $f_0(X,Z,S)$ alone, after integrating the normalized $g,\kappa_R,\kappa_C$. This is stronger than assigning a Gaussian actual law, but weaker physically than generating the field from $u$: the latter coupling is absent.

\section{The writer's exact driven comparison}
Remove the common Galilean phase and set $\xi=S-v\lambda t$, $a=\lambda t+\xi/v$. For sector one let
\[
 q''+\Omega^2q=\Omega^2b(a)+b''(a),\quad q(0)=q'(0)=0,
 \qquad\Omega=1/\mu.
\]
The normalized comparison is
\[
 G_1=\varphi(\xi)e^{-it/(2\mu)}
 \exp\{i[\mu q'(y-q/2)+\vartheta]\}\gamma(y-q),
\]
\[
 \vartheta'=\frac\mu2[\Omega^2b+b''](q-b).
\]
Sector zero has $q=\vartheta=0$. Direct substitution proves the fibre Schr\"odinger equation when the clock's residual kinetic term is omitted. For $\lambda=1$, $q=b$ and $\vartheta=0$ exactly.

The difference $e=q-b(a)$ solves
\[
 e''+\Omega^2e=(1-\lambda^2)b''(a).
\]
Sine/cosine integration gives $|e|<.001$, $|e'|<.07$ on the complete finite rectangle. Derivatives obey
\[
 |q|<25,\ |q'|<1630,\ |q''|<170000,\ |q'''|<36000000.
\]
The clock-envelope derivative norms are below $6,112,4000,190000$. Appendix~\ref{app:periodic} gives the polynomial derivative and envelope
integrals explicitly. Keeping the Gaussian polynomial before taking norms gives
\[
 \boxed{\norm{\partial_\xi^2G_i}<110000,\qquad
 \norm{p_Z\partial_\xi^2G_i},\ 
 \norm{(Z-z_0)\partial_\xi^2G_i}<3\times10^6.}
\]
Appendix~\ref{app:periodic} derives these jet estimates from the two
coherent-Gaussian derivative polynomials, with explicit coefficient bounds. Harmonic transport inverse engineering is established prior art \cite{Tor}.

\section{Finite interacting clock, not an infinite-mass limit}
The exact moving-frame Hamiltonian retains $-\partial_\xi^2/(2M)$. Consequently the comparison residual is $\partial_\xi^2G_i/(2M)$. Its propagator is the full interacting $S,Z$ propagator.

\begin{lemma}[Full-time auxiliary errors]\label{pa:lem:aux}
For $0\le t\le3$ and $M_-=.99\,10^{18}$,
\[
 \norm{\Xi_i-G_i}\le\frac{3(110000)}{2M_-}<2\times10^{-13},
\]
\[
 \norm{p_Z(\Xi_i-G_i)}
 \le\frac{3[6\times10^6+15000(110000)]}{2M_-}
 <3\times10^{-9}.
\]
\end{lemma}
\begin{proof}
The norm estimate is Duhamel with zero initial defect. For the derivative estimate solve the oscillator operator equations. The remaining $Z$ force has norm below $5000$, so
\[
 \norm{p_ZU(t,s)r}\le\norm{p_Zr}+\norm{(Z-z_0)r}
 +(t-s)5000\norm r.
\]
Apply this to the residual and integrate. Clock recoil is in $U$, not prescribed away. No exponential in $\Omega t$ is paid, and no source phase enters these $S,Z$ derivatives.
\end{proof}

\section{An earlier actual label is protected separately}
For the full initial offset support,
\[
 \lambda_+t_a+.5/127<1.2299,
 \qquad
 \lambda_-t_b-.5/127>1.2301+.04001.
\]
The compact comparison is therefore still in the common potential before $t_a$ and has completed writing at $t_b$. These guards concern a comparison support; they are not a claim that the exact clock has compact support at later times.

Before $t_a$, use
\[
 F_1(t,\xi,Z)=e^{-it/(2\mu)}
 \left[\varphi(\xi)+\frac{it}{2M}\varphi''(\xi)\right]\gamma(Z-z_0).
\]
Its residual norm is at most $t\norm{\varphi''''}/(4M^2)$. Both the exact common auxiliary evolution and each exact driven sector differ from $F_1$ by at most $t_a^2\norm{\varphi''''}/(8M^2)$. Therefore
\[
 \epsilon_{\rm pre}\le\frac{t_a^2(400000)}{4M_-^2}.
\]
This expansion has not discarded any exact clock tails.

Let $\rho_r,j_r$ be the isolated spinor rotor current. Its initial density is $1/2$, and circle Sobolev/energy estimates give
\[
 \norm{\rho_r}_\infty\le3K/h,\quad
 \norm{j_r}_\infty\le8K,\quad \norm{r_i'}_2\le K/h.
\]
Exact sector factorization bounds the pretrigger wave and its $X$ derivative errors by $\epsilon_{\rm pre}$ and $(K/h)\epsilon_{\rm pre}$. Thus
\[
 \norm{\delta j_X}_1\le4h\epsilon_{\rm pre},\qquad
 \norm{\delta\rho}_1\le3\epsilon_{\rm pre}.
\]
The isolated lifted cumulative rank along the coupled ordinary path obeys
\[
 \rho\dot{\mathcal K}=\rho_rj_X-j_r\rho,
\]
including its circulation. Hence
\[
 \boxed{\E_Q\Var_{[0,t_a]}\mathcal K
 \le120K\epsilon_{\rm pre}<6\times10^{-18}.}
\]
This is the physical clock-to-label error used in the periodic-current section. A global wave norm alone is not used as its replacement.

\section{Every admitted preparation-clock position is harmless for a reason}
The old preparation clock $C$ is retained but no longer controls the active programme. It is an ordinary free particle with $M_C=10^8$. In its co-moving coordinate, the free propagator and one-dimensional Sobolev bound give, through time three,
\[
 e_0\le\frac{3(112+4000)}{2M_C},\quad
 e_1\le\frac{3(4000+190000)}{2M_C},\quad
 |\dot C-1|\le\frac{e_1}{M_C(1-e_0)}<3\times10^{-11}.
\]
The displacement error is below $9\times10^{-11}$, less than the $.05$ plateau clearance. A first-exit argument proves the tube for every $C_0\in[-1.25,-.75]$. No averaging over $|\chi_C|^2$ is used. Its correlations are irrelevant to active histories by Theorem~\ref{pa:thm:recoil}.

The active writer clock is $S$. Its actual distribution is controlled by the original joint cap $f_0\le2$, not by a new Born-clock premise. A pointwise trajectory tube for every $S_0$ is neither claimed nor needed by the active-law current proof.

% Integrated section: periodic_history
\section{Reduced active law and periodic averaging}
Use the finite Hamiltonian and initial data of the periodic-model section. The recoil and auxiliary calculation above proves the exact relative dynamics and spectator reduction. The active initial reference density is
\[
 q_0(X,Z,S)=\tfrac12|\gamma(Z-z_0)|^2|\varphi(S)|^2.
\]
On the readiness box, $q_0>1/8$. Indeed $\exp((1.001)^2)<11/4$, $\sqrt\pi<9/5$, $Z_c<4/5$, and $2(11/4)(9/5)(4/5)<8$. Thus every active actual law is dominated by $16q_0$. This reference is only a proof measure.

\begin{lemma}[Periodic event averaging]\label{ph:lem:avg}
Let $A$ be any measurable $h$-periodic subset of $\T_2$, with volume fraction $\theta$. For a normalized actual marginal $f_X$ of zero-extension variation at most $V_X$,
\[
 |P_{f_X}(A)-\theta|\le\frac{hV_X}{8}=\frac{V_X}{4N}.
\]
\end{lemma}
\begin{proof}
The mean-zero function $\one_A-\theta$ has an $h$-periodic primitive of oscillation at most $h\theta(1-\theta)$. Subtract its midrange value; its sup norm is at most $h\theta(1-\theta)/2\le h/8$. Integration by parts against the circular BV derivative proves the estimate. Its variation is bounded by the given zero-extension variation, including boundary jumps. No smoothness of the event boundary and no equilibrium actual law is needed.
\end{proof}

\section{Rotor wave and velocity-weighted error}
In one sector write $z=(X-o)/h-i/2-d_i$, $\omega_i=\sqrt{g_i/\alpha}$, and $b_i=\cos(\omega_i t)$. Let $a_{\eta_t}$ be the unnormalized quintic taper: one on $|z|\le1/2-2\eta_t$, zero on $|z|\ge1/2-\eta_t$, and a $S_5$ transition between. The harmonic compact comparison carries amplitude $a_{\eta_t}(z/b_i)/\sqrt{2b_i}$ and phase $-\alpha K\omega_i\tan(\omega_it)z^2/2$ in each cell.

The physical initial wave is $1/\sqrt2$, so its initial error is not zero. Exact integrals give
\[
 d_0^2=\frac{643\eta_t}{231},\quad
 A_1^2=\frac{20}{7\eta_t},\quad
 A_2^2=\frac{240}{7\eta_t^3},\quad
 A_3^2=\frac{1440}{\eta_t^5}.
\]
The comparison lies in the quadratic potential core whenever it is used. Its residual norm is $A_2/(2\alpha Kb_i^2)$.

\begin{lemma}[Finite rotor comparison]\label{ph:lem:rotor}
Through capture time $t_b=1.28$, let $r_i$ be the exact rigid-sector wave. Then
\[
 \norm{r_i-\Phi_i}\le d_0+\epsilon_R=:d,
 \qquad \epsilon_R=\frac{2A_2}{.999K},
\]
\[
 \frac{\norm{p_X(r_i-\Phi_i)}}{mh}
 \le \frac{A_1}{.999K}+.501d_0+\frac{4.02A_3}{K^2}
                +1.002\epsilon_R=:q.
\]
For the selected row, $d<D=.007233$ and $q<Q=.003983$.
\end{lemma}
\begin{proof}
The residual norm integral uses $\omega_it<1.3$ and $\tan1.3<4$. The static nonnegative rotor Hamiltonian propagates its own square-root norm. Multiplying that estimate by $\sqrt{2/(\alpha K)}$ gives the initial derivative contribution $A_1/(\alpha K)+\omega_i d_0/2$. Differentiating the residual introduces $A_3/(2\alpha^2K^2b_i^3)$ and its harmonic phase derivative; the potential part adds the same $\omega_i/2$ residual bound. Integrate with $\int_0^{1.3}\sec^3\theta\,d\theta<8$, $\omega_i<1.002$, and $4/(\alpha^2\omega_i)<4.02$. Weak $H^3$ taper derivatives suffice; no fourth derivative of the sharp taper is invoked. The outward rational square-root calculation in Appendix~\ref{app:periodic}
proves the final decimals.
\end{proof}

At $t_a=1.22$, each $\Phi_i$ is supported in its own outcome region even with $|d_i|\le.01$. The wrong-side wave probability is at most $D^2$. For the isolated spinor rotor, the guidance map commutes with $X\mapsto X+h$. The initial preimage of the classifier is therefore periodic, including its circle circulation. Lemma~\ref{ph:lem:avg} and the wrong-side estimate yield the actual isolated-label error
\[
 |P(\ell(X^{\rm iso}_{t_a})=1)-p|\le\frac1{2N}+D^2.
\]

\section{Pretrigger disturbance belongs to the same ordinary model}
The recoil and auxiliary calculation above supplies the exact ordinary auxiliary estimates, including all clock tails and backreaction. In particular,
\[
 \epsilon_{\rm pre}\le\frac{t_a^2(400000)}{4M_-^2},\qquad
 I_{\rm pre}\le120K\epsilon_{\rm pre}<6\times10^{-18}.
\]
Here $I_{\rm pre}$ is expected variation of the isolated rotor's lifted rank along the exact active trajectory. It is not the old source-inclusive logical-current quantity with a substituted constant.

The initial rotor marginal has density at most one because its variation is at most two. Its initial cumulative-rank density is at most two. The $2N$ pulled-back classifier boundaries therefore have actual strip mass at most $8Nr$. With $r=10^{-9}$,
\[
 \epsilon_{\rm prelabel}
 \le8Nr+\frac{16(6\times10^{-18})}{r}
 =.000002144.
\]
The first term uses the actual marginal directly; the second uses the derived active domination factor16.

\section{Copying the actual earlier label}
Let $B_X$ be the union of the intervals at relative distances $h/5$ to $3h/10$ between adjacent zero and one wells. Let $H_X\in[0,1]$ be the quintic soft classifier, equal to the binary label outside these bands, with $|H_X'|\le20/h$.

During $[t_a,t_b]$, both $\Phi_i$ and $\partial_X\Phi_i$ vanish in these bands. Exact sector factorization gives
\[
 Q(X_{t_a}\in B_X)\le D^2,
 \qquad
 \E_Q\Var_{[t_a,t_b]}H_X(X_t)\le\frac65DQ.
\]
For clarity $Q$ in the last product denotes the fixed number $.003983$; $\E_Q$ denotes wave-reference expectation.

Use $I_0=[-8,8]$, $I_1=[16,32]$ and define
\[
 G_{\rm mis}(x,z)=1-(1-H_X(x))\one_{I_0}(z)-H_X(x)\one_{I_1}(z).
\]
For any continuous paths,
\[
 \one_{Z_{t_b}\notin I_{\ell(X_{t_a})}}
 \le\one_{B_X}(X_{t_a})+
 \Var H_X(X_t)+G_{\rm mis}(X_{t_b},Z_{t_b}).
\]
The comparison writer has completed its drive for every retained offset at $t_b$ and its wrong-inner-region amplitude is below $10^{-12}$. With the exact auxiliary error $e_A<2\times10^{-13}$,
\[
 \E_QG_{\rm mis}\le(D+e_A+10^{-12})^2.
\]
\begin{theorem}[Actual historical copy]\label{ph:thm:copy}
Uniformly over the original actual class and all qubit inputs,
\[
 \boxed{\epsilon_{\rm copy}
 =16\left[D^2+\frac65DQ+(D+2\times10^{-13}+10^{-12})^2\right]
 =.00222725479707774720002304.}
\]
This bounds failure to copy the actual coupled relative rotor label at $1.22$ into the pointer at $1.28$.
\end{theorem}
The internal sector in this proof is an orthogonal wave component, not a sampled actual spin. On the full wave the pointer velocity depends on the actual relative rotor coordinate through its branch weights.

\section{Whole-interval retention, including later rotor evolution}
The outer regions are $\mathcal R_0=[-10,10]$ and $\mathcal R_1=[14,34]$. Four width-two bands separate them from the corresponding inner regions. After capture, the auxiliary comparison center stays within $.012$ of its intended value and has speed below $.07$. Its band wave and momentum norms are below $10^{-14}$ and $10^{-13}$.

The exact rotor waves are used during retention; no harmonic comparison is continued through a caustic. Their unit norms suffice. With $e_A=2\times10^{-13}$, $p_A=3\times10^{-9}$, $\mu_-=.00999$, the total actual probability of leaving after starting in an inner region is bounded by
\[
 16\cdot4\cdot\frac{15}{16}\cdot\frac{43}{25}
 \left[.07\,10^{-28}
 +\frac{e_A10^{-13}+10^{-14}p_A+e_Ap_A}{\mu_-}\right]
 <10^{-16}.
\]
This uses band-integrated absolute current and is simultaneous over $[1.28,3]$;
Appendix~\ref{app:periodic} supplies the cutoff and Gaussian-tail calculation. Later rotor revival or source phase evolution is not frozen. The same-support actual-law neighbourhood adds $10^{-4}$ once to the complete history event.

% Integrated section: periodic_result
\section{Finite complete theorem}
Let $X=x-R$. Define $t_a=1.22$, $t_b=1.28$, $T_H=3$ and
\[
 I_0=[-8,8],\ I_1=[16,32],\qquad
 \mathcal R_0=[-10,10],\ \mathcal R_1=[14,34].
\]
The record symbol is $0$ or $1$ in the corresponding outer region and $\perp$ elsewhere. The relative rotor classifier is halfway between the nominal sector centers. No observation is performed at $t_a$; it identifies the earlier actual outcome which the pointer records.

\begin{theorem}[Complete direct ordinary record]\label{pr:thm:final}
Assume the stated conservative current flow on the ordinary model, the exact conserved internal sectors, the specified initial product waves and actual-law family, and the prescribed-profile realization of the periodic-model section. Uniformly over every normalized qubit/reference input, every admitted actual source and clock conditional law, and every $t\in[1.28,3]$,
\[
 \boxed{
 \TV\!\left(\mathcal L(\mathfrak r(Z_t)),
 \operatorname{Bernoulli}_{\{0,1,\perp\}}(p_{\rm tar})\right)
 \le\varepsilon_{\mathrm{per}}=.00443484008607784720002304<.01.}
\]
Moreover,
\[
 \boxed{
 P\{Z_t\in\mathcal R_{\ell(X_{1.22})}\ \forall t\in[1.28,3]\}
 \ge1-.00232725479707784720002304.}
\]
\end{theorem}
\begin{proof}
Exact relative-coordinate reduction and spectator independence leave the active initial law $f_0(X,Z,S)$. Its density is bounded by16 times the initial active wave reference. The isolated rotor's periodic-flow event is calibrated by $1/(2N)+D^2$. The current-sensitive quiet-prefix calculation changes its actual label by at most $8Nr+16I_{\rm pre}/r$. The source-free sector factorization is used explicitly; it is not inferred from a small source error.

The soft-classifier proof copies the actual coupled earlier label with error $C_{16}$, and the pointer-band current gives the whole holding error $H_{16}$. Add the same-support actual-law neighbourhood and target-projector error once. The proof compares the ordinary record to its own actual earlier label and then to the Born target. It does not need a small difference from squared-controller's individual trajectories.

For the simultaneous history claim, only copying, holding and the actual-law neighbourhood are required. All bounds concern one Hamiltonian and one parameter family. Source and preparation-clock conditionals integrate exactly because the event is independent of their initial values; they are not averaged under an assumed physical Born law.
\end{proof}

\begin{center}\small
\begin{tabular}{@{}p{.56\linewidth}p{.36\linewidth}@{}}\toprule
Contribution & Certified bound\\\midrule
Cell averaging & $1/512=.001953125$\\
Wrong-side rotor probability & $D^2=.000052316289$\\
Clock-prefix label change & $.000002144$\\
Actual historical copy $C_{16}$ & $.00222725479707774720002304$\\
Whole-interval holding $H_{16}$ & $10^{-16}$\\
Actual-law neighbourhood & $.0001$\\
Target-projector calibration & $.0001$\\\midrule
Complete sum & $.00443484008607784720002304$\\\bottomrule
\end{tabular}
\end{center}
Here $D=.007233$, $Q=.003983$, $e_A=2\times10^{-13}$, $r=10^{-9}$, and $I_{\rm pre}\le6\times10^{-18}$. The exact record fraction is
\[
 \frac{1732359408624159062509}{390625000000000000000000}.
\]
The unused direct probability margin is $.00556515991392215279997696$. It could be spent on a separately proved microscopic field-realization history error; it is not such an error certificate.


\begin{corollary}[Joint periodic-record law]\label{cor:periodicjoint}
With the same assumptions and parameters, the law of the actual earlier
label and its entire symbolic holding path has distance at most
$0.00443484008607784720002304$ from the ideal constant path with the
specified calibrated Bernoulli label.
\end{corollary}
\begin{proof}
For a base law, calibration and prefix transfer cost
$1/512+D^2+0.000002144$. Copying and holding cost
$C_{16}+10^{-16}$. Apply Lemma~\ref{lem:joint} with the isolated label
and these actual-history events. The actual-law neighbourhood costs
$10^{-4}$ once on the full path law because the same flow is used.
The target-projector allowance costs another $10^{-4}$. Their sum is
the displayed bound. This strengthens the presentation to a joint
observable using the already proved common-event argument, without
inferring path stability from norm error alone.
\end{proof}

\subsection{A genuinely nonequilibrium admissible stock}
On the readiness box take
$f_0=(1+0.8X)\one_{[-1,1]^2\times[0,1/4]}$.
Its mass is one, maximum is $1.8$, directional variations are
$(1.8,1,8)$ and rotor half-probability is $0.7$. Its clock is confined
to one half of a symmetric wave packet, so the initial full-law TV
distance from the wave measure is at least $1/2$. An admitted source
conditional is $(1+\tfrac12\operatorname{sgn}u_1)|\zeta|^2$.
These are examples, not new assumptions imposed on the whole class.
Invertibility of the full flow preserves fine-grained TV. Calibration
of a coarse record therefore does not prove that the full ensemble
equilibrates or that previously retained information is destroyed.

\part{A radial apparatus with preparation and a writer from entrance}
% Integrated section: radial_definition
\section{The fixed radial preparation operator}\label{sec:radial_definition}
We now use physical units. Put $\ell_0=10^{-4}\,\mathrm m$,
$T_0=0.03\,\mathrm s$, $m=2^{52}\hbar T_0/\ell_0^2$,
$M=10\,\mathrm{kg}$ and $v=1\,\mathrm{m/s}$. The relative radial
Hamiltonian has Dirichlet boundary at zero. The two ordinary Cartesian
bodies each have mass $2m$. Their supplied normalized free centre Gaussian
and constant angular factor give an exact $s$-wave reduction;
the centre and constant angular factor separate from the active dynamics.

All smoothing choices are fixed. Let
\[
 \kappa_\epsilon(x)=Z_\epsilon^{-1}
 e^{-1/(1-(x/\epsilon)^2)}\one_{|x|<\epsilon},\qquad
 \Theta(x)=\begin{cases}1&x\le0,\\
 [1+e^{-1/x+1/(1-x)}]^{-1}&0<x<1,\\0&x\ge1.
 \end{cases}
\]
Here $Z_\epsilon$ normalizes the mollifier. Let $B_9$ be the clamped
polynomial already defined, and let $b$ be the convolution of
$1+99B_9(s/0.1\,\mathrm s)$ with $\kappa_{10^{-5}\,\mathrm s}$.
Its constant extensions are one and one hundred. Set $a=b'/b$;
derivatives of these preparation functions are in clock time $s=S/v$.
The distinct pointer displacement introduced below will always carry
its own argument and derivative convention.

Take $L_i=10^{-4}\,\mathrm m$, $L_f=100L_i$. The radial comparison
amplitude is
\[
 \chi(s,r)=\frac{2r}{\pi^{1/4}[L_i b(s)]^{3/2}}
             e^{-r^2/(2L_i^2b(s)^2)}.
\]
The initial wave uses $\chi$ on the $b=1$ plateau, the specified
centre Gaussian, an arbitrary normalized qubit, and a clock envelope
centred at $S_c=-0.002\,\mathrm m$ with carrier $Mv$. The envelope
is the normalized convolution of
$\sqrt{64/(35\sigma)}\cos^4(\pi\xi/(2\sigma))
\one_{|\xi|\le\sigma}$, $\sigma=10^{-3}\,\mathrm m$,
with $\kappa_{10^{-10}\,\mathrm m}$. Its actual support halfwidth is
$0.0010000001\,\mathrm m$; it is not replaced by an unsmoothed packet.

Define an odd tent by $f_0(r)=r$ on $[0,R]$, $f_0(r)=2R-r$ on
$[R,2R]$ and zero beyond, with $R=0.201\,\mathrm m$. Set
$f=f_0*\kappa_{0.001\,\mathrm m}$ and $F(r)=2\int_0^r f(u)\dd u$.
Thus $f=r,F=r^2$ through $0.2\,\mathrm m$, $f=0$ from
$0.403\,\mathrm m$, $\norm f_\infty\le0.201\,\mathrm m$,
$\norm{f'}_\infty\le1$ and $F_\infty\le0.080802\,\mathrm m^2$.
Put
\begin{align}
 \omega_i&=\hbar/(mL_i^2),\quad E_0=3\hbar\omega_i/2,\\
 Q(s,r)&=m\omega_i^2r^2/(2b(s)^4)-E_0/b(s)^2,\\
 Q_c(S,r)&=\Gamma(S)\Theta((r-0.2\,\mathrm m)/(0.001\,\mathrm m))Q(S/v,r),\\
 V_c(S,r)&=-\frac m2a'(S/v)F(r)-\frac m2a(S/v)^2 f(r)^2
       -\frac{m^2a'(S/v)^2F(r)^2}{8Mv^2}.\label{eq:Vc}
\end{align}
The clock cutoff $\Gamma$ is one on $[-0.004,0.105]\,\mathrm m$,
zero outside $[-0.005,0.106]\,\mathrm m$, with left factor
$1-\Theta((S+0.005\,\mathrm m)/(0.001\,\mathrm m))$ and right factor
$\Theta((S-0.105\,\mathrm m)/(0.001\,\mathrm m))$. In particular,
the clock-only far-radial part of $V_c$ is retained. It is not silently
subtracted as a constant.

In dimensionless radius $x=r/\ell_0$, let $h=1/128$, $K=2^{38}$,
$\eta=2^{-17}$ and let $U_\eta$ be the capped quadratic of Part I.
The physical measured potentials are
\[
 W_i^{\rm phys}(r)=\frac\hbar{T_0}\Theta(x-1000)
                         \frac K2 U_\eta(x/h-i/2).
\]
The outer cutoff ends at $0.1001\,\mathrm m$. Its cap transition
width is $5.9604644775390625$ picometres. The finite radial classifier uses the centred representative
$(r/\ell_0)\bmod h\in[-h/2,h/2)$. It is zero when the absolute
value of that representative is at most $h/4$, one otherwise, and
undefined unless $0\le r/\ell_0<1000$. The boundary convention is
the frozen one; ordinary nonnegative modulo is not substituted.
There are $256000$ internal cuts plus an exterior ray. Define
$\gamma(s)=B_9((s-0.104\,\mathrm s)/(0.0001\,\mathrm s))$.
The preparation-plus-activation operator is exactly
\begin{equation}\label{eq:Hrad}
 H_{\rm rad}=\frac{p_S^2}{2M}+\frac{p_r^2}{2m}+V_c
       +(1-\gamma(S/v))Q_c
       +\gamma(S/v)\sum_i\Pi_iW_i^{\rm phys}.
\end{equation}
It is this operator, not a subsequently fitted effective wave, that is
enlarged by the pointer. Laboratory entrance is $-0.106\,\mathrm s$,
handoff is $-0.002\,\mathrm s$, witness is $0.0366\,\mathrm s$,
and holding is $[0.0384,0.09]\,\mathrm s$.

\subsection{The full initial-law contract}\label{sec:radial_law}
The new radius is coupled to a periodic-reference initial coordinate
by $X=(100r_0/\ell_0)\bmod2$. Retain one complete reference law
$\nu_{\rm ref}(\dd X,\dd\zeta_{\rm ref})$, where
$\zeta_{\rm ref}=(Z_{\rm ref},S_{\rm ref},u_{\rm ref},R_{\rm ref},C_{\rm ref})$.
The comparison rotor is fixed to the nominal Part I member:
$o=d_0=d_1=0$, $g_0=g_1=1$, and reduced mass $2^{52}$ in its
reference units. The radial theorem does not inherit the entire
independent gain, offset and mass rectangle of Part I.
It satisfies the periodic readiness and BV restrictions and, for this
radial construction, the additional full reference cap $\nu_{\rm ref}\le32\mu_{\rm ref}$,
where $\mu_{\rm ref}$ is the complete original reference entrance
wave measure. Its
marginal obeys $f_X\le1$. The earlier admitted arbitrary singular
passive laws are not all in this narrowed class.

Let $w_{\rm new}(\dd\mathrm{aux}\mid X)$ be the normalized
disintegration of the specified new entrance \emph{wave measure}.
The winding weights are $\chi_{100}(X+2k)^2/W_{100}(X)$ for positive
radius, together with the centre, angular, clock and pointer wave
factors. The dimensionless amplitude is
$\chi_{100}(y)=2y e^{-y^2/(2\cdot100^2)}/(\pi^{1/4}100^{3/2})$
for $y>0$, and zero otherwise. Here
\[
 W_{100}(X)=\sum_{k:X+2k>0}\chi_{100}(X+2k)^2
       \ge\frac{67499}{135000}=W_{\min}.
\]
Choose one normalized input-independent conditional measure
\[
 0\le g_{\rm new}(\dd\mathrm{aux}\mid X,\zeta_{\rm ref})
             \le2w_{\rm new}(\dd\mathrm{aux}\mid X).
\]
The complete actual coupling is $\nu_{\rm ref}g_{\rm new}$. Correlations
among new auxiliaries and with every retained reference variable are
allowed. The new physical clock and pointer are not identified with
the reference clock and pointer. Integration gives
\[
 C_{\rm micro}=2/W_{\min}=270000/67499,\qquad
 C_{\rm joint}=32/W_{\min}=4320000/67499.
\]
For the joint comparison measure use
\[
 \mu_{\rm joint}=2W_{100}(X)\,
       \mu_{\rm ref}(\dd X,\dd\zeta_{\rm ref})
       w_{\rm new}(\dd\mathrm{aux}\mid X).
\]
The reference wave has uniform $X$ marginal $\dd X/2$, so this is
normalized. The two likelihood bounds then give
$\nu_{\rm joint}\le(32/W_{\min})\mu_{\rm joint}$.
These are consequences of the complete coupling, not substitutes for it.
All original reference neighbourhood and projector-calibration allowances
remain charged in the final comparison.

For a concrete nonequilibrium example, take the biased reference law of
Part I and reference passives whose combined full likelihood ratio is
below $23.04<32$. Choose $g_{\rm new}=2\one_{S>S_c}w_{\rm new}$.
Evenness of the entrance clock makes this conditional normalized for
each $X$. It is input independent, satisfies the single cap and has
actual-wave TV at least $1/2$. Thus the permitted class is genuinely
nonequilibrium while remaining strongly restricted.

\subsection{The preparation gauge and why it is legitimate}
After the common Galilean carrier is removed, use the analytical unitary
$e^{-i\theta}$ with $\theta=ma(S/v)F(r)/(2\hbar)$. Expanding the
kinetic squares and \eqref{eq:Vc} gives exactly
\[
 H_g=\frac{p_r^2}{2m}+\frac{p_S^2}{2M}+vp_S
   +\frac12\{af,p_r\}+\frac12\{\delta,p_S\}+Q_c+V_A,
 \quad \delta=\frac{ma'F}{2Mv},\quad V_A=\gamma(W-Q_c).
\]
The clock drift $\delta$ remains. During preparation write $\tau=t_{\rm lab}+0.106\,\mathrm s$
for elapsed time, so $\tau=0$ at the entrance. The real helper
$G_g=\phi(S-S_c-v\tau)\chi(S/v,r)$ has residual
\begin{align*}
 R_{\rm core}&=\frac{\hbar^2}{2M}
 [(\phi\chi)_{SS}+2i\theta_S(\phi\chi)_S+i\theta_{SS}\phi\chi],\\
 R_{\rm tail}&=i\hbar a\phi[(f-r)\chi_r+(f'-1)\chi/2]
                       +(Q-Q_c)\phi\chi.
\end{align*}
The tails begin at twenty maximum packet widths and are bounded, never
deleted. Since the helper has no remote activation support during
preparation, activation acts on the error and is controlled by moving
clock collars. Those collars are proof multipliers, not restrictions of
the physical initial law.

For error momenta $P_r,P_S$ and a right-collar amplitude $n$, the
mass-weighted norm $Z_E=(P_r^2/m+P_S^2/M)^{1/2}$ satisfies a positive
system $Z_E'\le\gamma_0Z_E+e_0+c_0dt+L_A n$,
$n'\le g_0+b_0Z_E$. The fixed weight
$10^{-6}\sqrt{\mathrm J}$ makes both row growth coefficients below
$141\,\mathrm s^{-1}$. Gronwall and Gaussian residual integrals then
bound the exact gauged error. The stronger anisotropic and localized
Hessian estimates used by the writer are derived in Appendix~\ref{prep:main}. A scalar norm alone is never used to
infer those derivatives.

% Integrated section: radial_result
\section{Model, stock, and scope}

The active configuration is $(S,r,Z)\in\mathbb R\times\mathbb R_+\times
\mathbb R$, with a two-component spinor. The relative radius $r$ is the
exact $s$-wave reduction of the frozen two-body scalar model. The free
centre coordinate and constant spherical harmonic factor from the active
Hamiltonian. The clock $S$ and pointer $Z$ are effective scalar coordinates;
no three-dimensional lifting error for those two coordinates is asserted.

Write $\Pi_0,\Pi_1$ for the conserved qubit projectors. Retain the exact
radial preparation baseline bounded-phase preparation, smooth activation, finite radial
potentials, and all clock self terms. Denote that operator by
$H_{\mathrm{rad}}$. Its physical clock mass is $M=10\,\mathrm{kg}$, its carrier speed
is $v=1\,\mathrm{m/s}$, and its reduced radial mass is
\[
 m_r=2^{52}\hbar T_0/\ell_0^2,\qquad
 T_0=0.03\,\mathrm{s},\quad \ell_0=10^{-4}\,\mathrm m.
\]
The scalar preparation operator is defined in Section~\ref{sec:radial_definition}.
Its fixed cutoffs and admitted laws are part of the theorem. The appendices derive the specialized coefficient and current estimates
used below.

The enlarged autonomous Hamiltonian is
\begin{equation}\label{rad:eq:H}
 H=H_{\mathrm{rad}}+\frac{\hbar}{T_0}
 \left[\frac{-\partial_Z^2+Z^2-1}{2\mu_Z}
              +\Pi_1\{-F(s)Z+C_b(s)\}\right],
\end{equation}
where
\begin{gather*}
 \mu_Z=0.01,\quad s=\frac{S-S_{\rm on}}{vT_0},\quad
 S_{\rm on}=0.14145\,\mathrm m,\quad w=0.003,\\
 b(s)=24 B_9(s/w),\qquad
 F=b/\mu_Z+\mu_Z b'',\qquad
 C_b=b^2/(2\mu_Z)+\mu_Z b b''/2.
\end{gather*}
Here $B_9(x)=126x^5-420x^6+540x^7-315x^8+70x^9$ on $[0,1]$,
extended by $0$ and $1$. Primes denote $s$ derivatives. Thus the pulse
width is $90\,\mu\mathrm s$ in the characteristic clock, and its last
spatial point is $0.14154\,\mathrm m$. The physical pointer mass is
$10^{-19}\,\mathrm{kg}$; its length unit is
$\sqrt{\hbar\mu_Z T_0/10^{-19}\mathrm{kg}}\simeq0.562469\,\mathrm{nm}$.
Every displayed compensation term remains in~\eqref{rad:eq:H}.

At laboratory time $t_0=-0.106\,\mathrm s$, the state is precisely the
original full small-Gaussian radial and smooth compact clock stock,
with its Galilean carrier and arbitrary normalized qubit, tensored with
\[
 g_0(Z)=\pi^{-1/4}\exp(-Z^2/2).
\]
There is no wave truncation or later preparation reset. The original
radial width is $10^{-4}\,\mathrm m$; the clock is the positive
$10^{-10}\,\mathrm m$ mollification of its normalized $\cos^4$ packet
of half-width $10^{-3}\,\mathrm m$. Its initial centre is
$S=-0.002\,\mathrm m$. The analytical handoff is
$t_h=-0.002\,\mathrm s$, with clock centre $0.102\,\mathrm m$.

The narrowed nominal reference-law class of Section~\ref{sec:radial_law} is retained exactly. The pointer is
included in the same normalized, input-independent joint auxiliary
conditional rule $g_{\rm new}\leq2w_{\rm new}$. This is one rule for
the whole auxiliary configuration, allowing the declared correlations;
it is not a separate cap or independent Born assignment for each added
coordinate. Its disintegration gives
\begin{equation}\label{rad:eq:cap}
 \nu_{\rm micro,0}\leq C|\Psi_0|^2\dd q,
 \qquad C=\frac{270000}{67499}.
\end{equation}
The retained periodic reference marginal, its allowed neighbourhood, and
its target-calibration allowance remain part of the hypotheses. The cap
alone is not substituted for that complete reference-coupling contract.

\section{Observable and theorem}

Let $L$ be the frozen radial detector classifier evaluated at the
\emph{actual} radius at $t_a=0.0366\,\mathrm s$, with the inherited
cemetery outcome for a failed earlier event. Set
\[
 R_0=[-10,10],\qquad R_1=[14,34],\qquad
 R(Z)=\begin{cases}0&Z\in R_0,\\1&Z\in R_1,\\
 \dagger&\text{otherwise.}\end{cases}
\]
The complete holding interval is $I=[0.0384,0.09]\,\mathrm s$ and the
observable is
\[
 Y=\bigl(L,(R(Z_t))_{t\in I}\bigr).
\]
An incorrect or undefined record at any holding time is a failure.
An undefined earlier label is a failure even if a cemetery record could
be assigned the same symbol. No time grid replaces the interval.

\begin{theorem}[Effective joint-path theorem]\label{rad:thm:main}
For Hamiltonian~\eqref{rad:eq:H}, the full entrance stock above, every
normalized qubit $\alpha_0|0\rangle+\alpha_1|1\rangle$, and every admitted
original reference/joint auxiliary law, there is the unchanged compatible
reference coupling for which
\begin{align}
 \Prob(L\neq L_{\mathrm{per}})&<0.003331233001,\label{rad:eq:label}\\
 \Prob\bigl(L\text{ fails or }\exists t\in I:R(Z_t)\neq L\bigr)
      &<0.000552421956.\label{rad:eq:copy}
\end{align}
Writing $p=|\alpha_1|^2$ and $B\sim\mathrm{Bernoulli}(p)$ only as an
ideal comparison variable,
\begin{equation}\label{rad:eq:TV}
 \TV\!\left(\mathcal L(Y),
        \mathcal L\bigl(B,(B)_{t\in I}\bigr)\right)
       <0.006086770113<0.01.
\end{equation}
The microscopic dynamics contains no sampled sector variable $B$.
\end{theorem}

Theorem~\ref{rad:thm:main} is about one effective scalar Hamiltonian. It is
not a claim that the finite source candidate realizes this Hamiltonian.
The difference to $0.01$ is an unachieved sufficient allocation for a
future compatible microscopic comparison, not a measured or proved error
of added hardware.

\section{A conservative flow for the effective model}

For current-based existence background see \cite{TT}. The three active kinetic operators have positive coefficients. The
pointer-linear term has bounded clock coefficients and is infinitesimally
bounded relative to the free kinetic plus oscillator operator. Square
completion gives a finite lower bound on the potential. Consequently
the full operator is self-adjoint on that reference operator's domain.

After odd extension in $r$, the original stock belongs to $D(H^2)$.
The writer coefficients vanish on its clock support. Functional calculus
therefore gives $\Psi\in C_tD(H^2)$ and $\partial_t\Psi\in C_tD(H)$.
The writer potential has the locally bounded weak second derivatives
needed for local elliptic regularity: $B_9$ is $C^4$, while its occurrence
through $b''$ gives only the required $C^2$ statement. Thus
$D(H^2)\subset H^4_{\rm loc}$ and $D(H)\subset H^2_{\rm loc}$.
In three active dimensions these imply the local continuity and local
Lipschitz velocity needed for a unique nonnodal guiding flow. No unjustified
$H^6$ claim for the writer is used.

Energy conservation and $\|H\Psi\|<\infty$ give, on any finite interval,
\begin{align*}
 \int\!\dd t\int |J_j|\dd q
 &\leq\frac{\hbar}{m_j}\int\!\dd t\,\|\Psi\|\|\partial_j\Psi\|<\infty,\\
 \int\!\dd t\int \rho\,|D_t\log\rho|\dd q
 &\leq\int\!\dd t\left(2\|\Psi\|\|\partial_t\Psi\|
       +2\sum_j\frac{\hbar}{m_j}\|\partial_j\Psi\|^2\right)<\infty.
\end{align*}
Stopped partial equivariance and exhaustion exclude finite-time escape
and node hitting almost surely. The origin is an odd-extension node.
The free centre flow and constant angular factors then restore the
two-body Cartesian description. Domination~\eqref{rad:eq:cap} transfers
the null exceptional set and transports the same cap along this one flow.
The much higher-dimensional interacting finite-field model requires its
own existence proof.

\section{Preparation and prewitness comparison}

Appendix~\ref{prep:main} proves the preparation estimates used here.
The writer's presence during preparation is paid by a remote-support
Duhamel estimate and weighted oscillator hierarchy, from the original
stock. It gives the handoff norm change below $2.768\times10^{-30}$,
first moving-annihilator norm below $5.050\times10^{-24}$, and second
annihilator norm below $1.229\times10^{-17}$. The strengthened handoff
clock-error derivative is
\[
 \|\partial_S^2(\Psi_{\rm new}-\Psi_{\rm old}g_0)\|
       <1.829\times10^{19}\,\mathrm m^{-2}.
\]
The proof uses ordinary fourth spatial derivatives of the smooth
preparation error, with nonzero Gaussian tails and cutoff derivatives
retained. It never assumes the capped reference entrance belongs to a
fourth radial Hamiltonian graph domain.

The preparation history itself is rederived with the original explicit
Gaussian material rank $K(S,r)$. Its $Z$ derivative is zero, but its
clock and radial currents are those of the new exact wave. Both phase
restoration corrections are included. For $N=256001$ rank cuts, including
the exterior ray, the preparation event and handoff conditional-CDF
bridge are respectively below
\[
 3.084950\times10^{-6},\qquad 6.1578014\times10^{-5}.
\]

For the later prefix, use the constant Galilean frame for both waves:
$\Psi_b=e^{-\ii MvS/\hbar+\ii Mv^2t/(2\hbar)}\Psi$ and
$H_b=H-\ii\hbar v\partial_S$, with the identical scalar terms.
Let $\psi_{\rm old}$ denote the exact old wave in that frame, with
$\psi_{{\rm old},i}=\Pi_i\psi_{\rm old}$ including its qubit coefficient.
Define the normalized Weyl
comparison $G_*=V(S)(\psi_{\rm old}g_0)$, with sector-one factor
\[
 \gamma_1(s,Z)=\pi^{-1/4}
 \exp\!\left[-(Z-b)^2/2+\ii\mu_Z b'(Z-b/2)\right]
\]
and sector-zero factor $g_0$. Direct differentiation gives the exact
residual
\[
 (H_b-\ii\hbar\partial_t)G_*
 =-\frac{\hbar^2}{2M}\sum_i
       \left(2\psi_{{\rm old},i,S}\gamma_{i,S}
                       +\psi_{{\rm old},i}\gamma_{i,SS}\right).
\]
It is supported only on the pulse. The displaced stationary plateau
has no residual source. Keeping the entire old clock tail and the
leading edge of its characteristic stock gives
\[
\sup_{t_h\leq t\leq t_a}\|\Psi_{
 \rm new,b}-G_*\|<5.056\times10^{-22}.
\]
This norm estimate is not used as a trajectory-agreement assertion.

For $\delta=\Psi_{\rm new,b}-G_*$, two fixed left localizations and the
oscillator moment hierarchy give
\[
 \sup\|\delta_S\|\leq0.1\,\mathrm m^{-1},\qquad
 \sup\|\delta_{SS}\|<1.118\times10^{21}\,\mathrm m^{-2}.
\]
The differentiated residual includes the coefficients $1,5/2,2,1/2$
on $\gamma_S\psi_{SSS}$, $\gamma_{SS}\psi_{SS}$,
$\gamma_{SSS}\psi_S$, and $\gamma_{SSSS}\psi$. Its old-wave input is
\[
 \int_{t_h}^{t_a}\|\psi_{{\rm old},SSS}\|\dd t
                <9.747\times10^{46}\,\mathrm{s/m^3}.
\]
On the remote preparation-force support the characteristic helper is
identically zero, so propagation of this third clock derivative uses
the already bounded old \emph{error} derivatives there. All activation
terms elsewhere remain.

\section{Rank current and actual earlier label}

Appendices~\ref{old:main} and~\ref{prefix:main} provide the quantitative
inputs for this step. Use the new wave's own radial conditional CDF given the live pair
$(S,Z)$. Its material derivative has clock and pointer contributions;
the radial current cancels pointwise in the conditional-rank identity.
The hidden-coordinate absolute values precede all integrations. The
hybrid comparison yields
\[
 I_{\rm post}=I_{\rm old}+\Delta I_S+I_Z
       <3.695139\times10^{-15},\quad
 \Delta I_S<4.547787\times10^{-16},\quad
 I_Z<2.574815\times10^{-15}.
\]
The rank collar is optimized once on this sum, giving
$2C\sqrt{2NI_{\rm post}}$. Independently optimized component costs
are not silently substituted for that expression.

The common good-clock event requires
$|S_h-0.102|\leq0.000839\,\mathrm m$ and whole-path clock deviation
at most $10^{-5}\,\mathrm m$ through $0.09\,\mathrm s$.
On that event $S_t<S_{\rm on}$ for every prewitness time, with a strict
$1\,\mu\mathrm m$ margin at $t_a$. Hence $G_*=\psi_{\rm old}g_0$
there. The same event places the witness conditional age in
$[355/300,377/300]$, completes the writer by $0.038389\,\mathrm s$,
and keeps the clock beyond the pulse throughout the hold.

The earlier-label composition retains the new preparation event, both
distinct exact-CDF bridges, the original conditional age-union and
outer-ray events, the common good-clock complement, the actual radial
exit event, and the unchanged conditional-reference endpoint comparison.
The new absolute barrier estimate gives radial-exit probability below
$9.127218\times10^{-8}$ through readout. These terms prove
\eqref{rad:eq:label} on the original reference coupling. Conditional rank
uniformity is used only under its own wave law, then transferred with
the same cap $C$; no conditional Born law is assigned to actual states.

\section{Copy, whole hold, and the direct path bound}

Appendix~\ref{hold:main} proves the complete copy and holding inequalities.
The exact moving annihilator is
\[
 A=\partial_Z+Z-\Pi_1\bigl(b+\ii\mu_Zb'\bigr).
\]
With $c=b+\ii\mu_Zb'$ and the characteristic variables, its forced
equation has the sole finite-clock source
\[
 (\ii\partial_\tau-\widetilde H-1/\mu_Z)A\Psi
 =-\epsilon\Pi_1(c'\partial_s+c''/2)\Psi,
 \qquad \epsilon=\frac{\hbar}{Mv^2T_0}.
\]
In the completed-clock region the pointer current is controlled by
$\operatorname{Im}(\Psi^\dagger A\Psi)/\mu_Z$. Positive inner pointer
regions $[-8,8]$ and $[16,32]$, followed by soft barriers to $R_0,R_1$,
give an entire-interval estimate; recrossings are charged by absolute
variation. A soft radial classifier on the same exact trajectories
connects the conserved-sector proxy to the actual earlier position label.
It does not replace that label by a sampled logical bit.

The copy proof retains its baseline radial band current, finite-clock
radial correction, positive pointer-inner event, entire hold current,
and the shared clock-core, clock-deviation and radial-exit exceptions.
The hold-current term is below $3.272\times10^{-9}$; the full sum is
\eqref{rad:eq:copy}. Every term concerns the new Hamiltonian and original
entrance law.

In the union of failed earlier-label transfer and failed record history,
the same three exceptional events need be charged only once. Their
identity is explicit in both proofs; this is a union of events, not
subtraction of unrelated numerical upper bounds. The outward arithmetic in Section~\ref{sec:events} gives
\[
 p_{\rm union}<0.003879184824.
\]
The preserved reference earlier-label estimate, including both original
$0.0001$ allowances, is
$e_{\rm label,26}=0.002207585289$. Coupling first to the correct constant
record of $L_{\mathrm{per}}$ and then to its target Bernoulli label proves
\[
 \TV(\mathcal L(Y),\mathcal L(B,(B)_{t\in I}))
 \leq p_{\rm union}+e_{\rm label,26}<0.006086770113.
\]
This establishes the whole observable of Theorem~\ref{rad:thm:main}.

The former comparison to the entire imperfect periodic reference apparatus has the
different sufficient upper expression
\[
 p_{\rm union}+h_{\mathrm{per}}<0.006206439621,\qquad
 h_{\mathrm{per}}=0.00232725479707784720002304.
\]
That expression does not meet the old target
$0.00556515991392215279997696$. Its failure is not an actual error lower
bound. The theorem uses the expressly permitted direct ideal joint-path
alternative.

% Integrated section: technical_interfaces
\section{Conditional-rank stability with zero fibres included}\label{sec:rankproof}
This section supplies the nonlinear step behind the radial proof. It is
useful precisely because the actual clock/pointer marginal need not have a
positive lower bound. Let $q=\psi(S,\cdot)$ be a radial spinor slice,
$u=\psi_S$, $z=\psi_{SS}$, $a=\norm q$, $\eta=q/a$, and let
$P_r$ multiply by $\one_{[0,r]}$. Set $K=\langle\eta,P_r\eta\rangle$
and $Z_\eta=(P_r-K)\eta$. The constant clock carrier is removed from
all derivatives, but its actual drift is retained in continuity.

\begin{lemma}[Exact weighted rank current]\label{lem:rankcurrent}
On nonzero fibres, with $\kappa=\hbar/M$,
\begin{equation}\label{eq:rankF}
 \rho\frac{\dd K}{\dd t}=\mathcal F[\psi]
 =\kappa\left\{2\Im(\eta(r)^\dagger u(r))
          \Re\langle Z_\eta,u\rangle
       -a|\eta(r)|^2\Im\langle Z_\eta,z\rangle\right\}.
\end{equation}
For two live conditioning coordinates $S,Z$, the right side is the sum
of the corresponding two functionals, with coefficients $\hbar/M$ and
$1/(\mu_ZT_0)$ respectively.
\end{lemma}
\begin{proof}
Write $w=\int\rho\dd r$, $A_r=\int_0^r\rho\dd y$, and
$K=A_r/w$. Let $\delta B$ and $\delta j$ be the cumulative and full
clock-current integrals after subtracting the constant carrier. Continuity
and zero origin flux give
\[
 \frac{\dd K}{\dd t}
  =\frac{\delta J}{\rho}K_S
       -\frac{\partial_S\delta B-K\partial_S\delta j}{w}.
\]
The radial current has cancelled pointwise against $K_r\dot r$.
Now $K_S=2\Re\langle q,(P_r-K)u\rangle/a^2$ and
$\partial_S\delta B-K\partial_S\delta j
=\kappa\Im\langle q,(P_r-K)z\rangle$; the $u^\dagger u$ term
is real. Substitution proves \eqref{eq:rankF}. Applying continuity with
both transverse currents gives the two summands, not an additional
mixed-derivative term. Each local product sums the original spinor
components. No component has been sampled as an actual sector.
\end{proof}

\begin{lemma}[Stability without an exact-marginal denominator]\label{lem:rankstable}
For normalized $\psi,G$, write $E=\psi-G$, and suppose the quantities
\[
 D=\norm E,\quad e_1=\norm{E_S},\quad e_2=\norm{E_{SS}},\quad
 g_j=\norm{\partial_S^jG},\quad
 Q_G^2=\int\frac{\norm{G_S(S,\cdot)}^4}{\norm{G(S,\cdot)}^2}\dd S
\]
are finite. The ratio is zero where $G$ and its required derivatives
vanish. Then
\begin{equation}\label{eq:rankstable}
 \norm{\mathcal F[\psi]-\mathcal F[G]}_1
 \le\frac\hbar M
 (4g_1e_1+2e_1^2+14DQ_G+9Dg_2+e_2).
\end{equation}
Here $\mathcal F[G]$ is algebraic evaluation, including whatever
residual the comparison has under the actual equation.
\end{lemma}
\begin{proof}
At a fixed slice put $g=G(S,\cdot)$, $b=\norm g$, $\xi=g/b$,
$U=G_S$, $V=G_{SS}$, and $d=\norm{q-g}$,
$d_1=\norm{u-U}$, $d_2=\norm{z-V}$. The elementary estimates
\[
 \norm{\eta-\xi}\le2d/b,\quad \norm{Z_\eta}\le1/2,
 \quad\norm{Z_\eta-Z_\xi}\le3\norm{\eta-\xi},
 \quad\norm{|\eta|^2-|\xi|^2}_1\le2\norm{\eta-\xi}
\]
follow by adding and subtracting $q/b$, by projection variance, and by
$|K_\eta-K_\xi|\le2\norm{\eta-\xi}$. Split the first term of
\eqref{eq:rankF} first in $u-U$ and then in normalized directions.
Its derivative part is at most $4\norm U d_1+2d_1^2$; its direction
part is at most $7\norm{\eta-\xi}\norm U^2\le14d\norm U^2/b$.
For the second term, the $z-V$ part is at most $ad_2$. Splitting the
remaining amplitude, density and projection yields
$(d/2+4b\norm{\eta-\xi})\norm V\le9d\norm V$.
Integrate these inequalities and use Cauchy--Schwarz and
$\int a^2\dd S=1$. If $a=0$, direct bounding of $\mathcal F[G]$
works with $d=b$; if $b=0$ and its derivatives vanish, the bound
$2d_1^2+ad_2$ suffices. Thus zero fibres are included, not discarded.
\end{proof}

For $G=\phi(S-S_c-vt)u(S,r)$ with real $\phi$ and normalized $u$,
\[
 Q_G\le\norm{\phi'^2/\phi}_2+\sup_S\norm{u_S}^2,
 \qquad
 \norm{\phi_0'^2/\phi_0}_2
  =\frac{4\pi^2}{\sigma^2}\sqrt{\frac3{35}}.
\]
The last equality is a direct trigonometric integral for the unsmoothed
cosine packet. Positive convolution and weighted Cauchy--Schwarz give
$\phi'^2/\phi\le Z^{-1}\kappa*(\phi_0'^2/\phi_0)$; Young's inequality
then adds only the normalization multiplier at most $1.000001$.
No actual conditional law is asserted uniform in this argument.

\begin{lemma}[Deterministic-time CDF interface]\label{lem:cdf}
For every measurable radial cut $b(S)$,
\[
 \int w(S)|K_\psi(S,b(S))-K_G(S,b(S))|\dd S\le2\norm{\psi-G}.
\]
Under the already transported actual-law cap $C$, the corresponding
binary-test mismatch is at most $2C\norm{\psi-G}$ per cut.
\end{lemma}
\begin{proof}
Multiply the CDF difference by $w=\int|\psi|^2$. The exact identity is
\[
 w(K_\psi-K_G)=\int(\one_{r\le b}-K_G)(|\psi|^2-|G|^2)\dd r.
\]
The coefficient has modulus at most one and the density difference has
$L^1$ norm at most $2\norm{\psi-G}$. Uniformity of the exact rank
under its own conditional \emph{wave} measure identifies this integral
with the test mismatch. Domination then supplies the actual-law bound.
Zero comparison fibres are charged by their exact mass. The same proof
works with live conditioning pair $(S,Z)$. It is a fixed-time result;
it is not substituted at a random clock hitting time.
\end{proof}

\section{How the exact writer controls the new rank terms}
Let $G_*=V(S)(\psi_{\rm old}g_0)$ and $\delta=\Psi-G_*$ as above.
Throughout this prefix section $T=t_a-t_h=0.0386\,\mathrm s$,
with elapsed time measured from the handoff $t_h=-0.002\,\mathrm s$.
On $S<S_{\rm on}$ the displacement and its derivatives are the identity.
Every hybrid functional integral in this section is restricted to
$Q_{\rm pre}=\{S<S_{\rm on}\}$ before integration. The separately
charged good-clock event keeps actual prefix trajectories in that set.
Global error and moment norms may upper-bound these restricted integrals,
but $G_*=\psi_{\rm old}g_0$ and the corresponding small pointer-jet
substitution are asserted only there. No derivative of the region's
indicator is taken, and no full-space equality of the two references is
being asserted.

For the clock component use the original characteristic helper and
its already proved positive current majorant. Inserting the sum of its
old error and $\delta$ into \eqref{eq:rankstable} gives the increment
\[
 \Delta I_S\le\frac\hbar M T
 [4(g_{1,o}+e_{1,o})d_S+2d_S^2+14dQ_{G,o}+9dg_{2,o}+d_{SS}].
\]
For the pointer component use the \emph{different} reference
$\psi_{\rm old}g_0$, whose radial direction is independent of $Z$.
Its pointer rank functional is exactly zero. Gaussian integration gives
$g_{1,Z}=1/\sqrt2$ and $g_{2,Z}=Q_Z=\sqrt3/2$. Therefore
\[
 I_Z\le\frac{T}{\mu_ZT_0}
 [2\sqrt2d_Z+2d_Z^2+(23\sqrt3/2)d+d_{ZZ}].
\]
The references are chosen separately because each has the appropriate
proved coefficients; no wave or current is transferred between them
without the displayed identity.

Write $A=\partial_Z+Z-\Pi_1c(s)$, $c=b+i\mu_Zb'$.
On the quiet prefix $A=B=\partial_Z+Z$, and
$B\delta=A\Psi$, $B^2\delta=A^2\Psi$. Oscillator integration by parts gives
\[
 d_Z\le\sqrt{A_1^2+d^2},\qquad
 d_{ZZ}\le\sqrt{A_2^2+4A_1^2+3d^2},
 \quad A_j=\norm{A^j\Psi}.
\]
Indeed $\norm{(p_Z^2+Z^2)f}^2
=\norm{B^2f}^2+4\norm{Bf}^2+\norm f^2$, while
$\Re\langle p_Z^2f,Z^2f\rangle=\norm{Zf'}^2-\norm f^2$.
All identities hold on the complete vector-valued fibres before
integration. A clock indicator is not differentiated.

\subsection{Derivative closure and its proof}
Appendix~\ref{prefix:main} supplies the Gaussian polynomial recursion,
the four simultaneous positive moment inequalities, and the strict
supersolution tests for the same new wave. It retains the additional
$-\kappa\Pi_1(c')^2\Psi/(vT_0)^2$ source in the $A^2\Psi$ equation.
The localized calculation then improves the pointer-weighted first
clock-error estimate to $Q(\delta_S)<616735\,\mathrm m^{-1}$ and proves
\[
 \sup\|\delta_S\|<0.1\,\mathrm m^{-1},\qquad
 \sup\|\delta_{SS}\|<1.117215\times10^{21}\,\mathrm m^{-2}.
\]
The old third-clock input and all differentiated residual coefficients
are included there. Appendix~\ref{prep:main} derives the original-stock
handoff data, including the fourth-spatial estimate and its Gaussian-tail
terms. Thus these derivative bounds are established for the declared
entrance wave, rather than imposed on a replacement handoff wave.

\section{Outward arithmetic and the complete event sum}\label{sec:events}
The preceding arguments identify each error as an event or a nonnegative
current functional under the same initial coupling. The preparation,
old-wave, new-prefix and copy/hold appendices specify their coefficient
formulae and evaluation bounds. No numerical certificate is used in
place of an analytic inequality.

For reproducible square-root enclosures, given a nonnegative rational
$x$ and $n=10^{60}$, choose the least integer $k$ with $k^2\ge n^2x$;
then $k/n$ is an upper bound. Polynomial integrations are exact rational
operations, and positive Taylor remainders or explicit derivative bounds
control the nonpolynomial terms as detailed in the appendices. The
following current ceilings have been rounded outward before combining:
\[
\begin{split}
 I_{\rm old}&<6.655453611235769\times10^{-16},\\
 \Delta I_S&<4.547786946512146880\times10^{-16},\\
 I_Z&<2.574814383253091344\times10^{-15}.
\end{split}
\]
The rank row below is evaluated as $2C\sqrt{2N(I_{\rm old}+\Delta I_S+I_Z)}$,
with $C=270000/67499$ and $N=256001$. Every displayed table entry is
itself a rational upper bound; unlike a shortened diagnostic display,
these entries may be added directly. The harmless rounding of the
pointer-inner event is deliberately included in the sum.
\begin{center}\small
\begin{tabular}{@{}lr@{}}\toprule
Event or contribution & Outward upper bound\\\midrule
Conditional reference endpoint & $0.002588532236512050$\\
Preparation rank & $0.000003084949592861$\\
Initial exact CDF & $0.000061578013525596$\\
Conditional age union & $0.000263432327209778$\\
Conditional exterior union & $0.000000015476669581$\\
Terminal exact CDF & $0.000062145214303409$\\
Aggregate own-rank history & $0.000347974650032331$\\
Copy baseline & $0.000537032962395582$\\
Finite-clock copy correction & $0.000010915589215099$\\
Positive pointer-inner event & $0.000000000000000001$\\
Holding current & $0.000000003271569844$\\
Shared clock core & $0.000004378860028174$\\
Shared whole-interval clock deviation & $0.000000000000106690$\\
Shared actual radial exit & $0.000000091272176629$\\
\\\bottomrule
\end{tabular}
\end{center}
The first seven rows belong to earlier-label transfer, the next four to
copying and retention, and the last three to both arguments. In both
proofs the last three are the same clock-core, entire-clock-deviation
and actual-radial-exit events of the new wave. Their sum is therefore
charged once in the union. Direct addition of the displayed rational
ceilings gives
\[
\begin{aligned}
 p_{\rm label}&\le 0.003331233000157099<.003331233001,\\
 p_{\rm copy/hold}&\le 0.000552421955492019<.000552421956,\\
 p_{\rm union}&\le 0.003879184823337625,\\
 p_{\rm union}+.002207585289&\le 0.006086770112337625<.006086770113.
\end{aligned}
\]
The reference-label ceiling $.002207585289$ already includes both
original $10^{-4}$ allowances. Applying Lemma~\ref{lem:joint} proves
the advertised joint-record bound without counting either allowance
again. The explicit outward sums also show that the last displayed
headline digits do not depend on substituting rounded intermediate
values into an unreported computation.

\section{Comparison and physical boundary}
The two theorems share a statistical and current-based proof strategy,
but their constants cannot be interchanged. The periodic theorem allows
its explicitly admitted inaccessible quantum reference and singular
passive-coordinate laws. The radial theorem is stated for a normalized
two-component qubit with no arbitrary inaccessible quantum-reference
extension, and narrows the reference law to a full-cap class before
adding its single joint auxiliary conditional. Its small entrance,
preparation, finite radial cutoff and added writer require new proofs.

The periodic apparatus does not generate its prescribed potential from
the retained elastic modes: they are decoupled spectators. The radial
apparatus does not acquire a finite material source merely because a
separate field candidate exists. Source preparation, actual-position
loading or proxy agreement after that enlargement, finite propagation,
clock/pointer support and whole-interval retention must be shown for one
combined Hamiltonian. The unspent distance below one percent is only a
possible allocation in such a future proof.

\section{Conclusion}
The achieved statements concern calibrated, historically faithful
finite records under explicit deterministic dynamics and restricted
nonequilibrium laws. Their mathematical content can be evaluated
without adopting a broader philosophical framework. What is not
derived is the physical provenance of every initial-law restriction,
the selection of guidance among alternative dynamics, or an
established-matter realization of every interaction. Independent
mathematical review should first test the precise current/domain and
event-composition claims under the stated assumptions. A physical
programme would additionally need a material model and its own
quantitative reduction to the retained-record observable.

\paragraph{Funding and collaboration.}
This research was conducted independently by Jeremy Rodgers without
external research funding. The author welcomes collaboration on
independent mathematical review, computational replication,
physical implementation and empirical testing, particularly where
specialist expertise, equipment or institutional resources are required.

\paragraph{AI assistance and verification.}
The author developed this research programme, including
the mathematical arguments. AI tools were used as an aid in that work and in preparing this
publication version, including manuscript consolidation, literature checking, mathematical checks,
and internal computational checks. The author retains responsibility for the manuscript.

\clearpage
\appendix

\section{Explicit estimates for the periodic apparatus}\label{app:periodic}
This appendix supplies the polynomial and Gaussian calculations used in
Part I. All parameter bounds refer to that part's finite rectangle.
They do not enlarge the law class of Part II.

\subsection{Polynomial derivatives and clock normalization}
Writing $z=1-2a$, direct differentiation gives
\[
 B_9'(a)=630a^4(1-a)^4,\quad
 B_9''(a)=\frac{315}{8}z(1-z^2)^3,\quad
 B_9'''(a)=-\frac{315}{4}(1-z^2)^2(1-7z^2).
\]
The first maximum is $315/128$; the second occurs at
$|z|=1/\sqrt7$ and is below $10$. For the third, the only interior
stationary squared arguments are $0$ and $3/7$, giving a bound $315/4$.
Consequently the writer displacement has derivative bounds
\[
 \|b'\|_\infty<1478,\qquad \|b''\|_\infty<150200,
 \qquad \|b'''\|_\infty<2.958\times10^7.
\]
For the four envelope derivatives the exact integrals
$J_k=\int_0^1(B_9^{(k)}(a))^2\dd a$ are
\[
 (J_1,J_2,J_3,J_4)
 =\left(\frac{4410}{2431},\frac{5040}{143},
              \frac{272160}{143},\frac{1814400}{11}\right).
\]
Also $\int A_c^2=176818/230945=Z_c$. The change of variable on
both transition intervals gives
$\|\varphi^{(k)}\|^2=2\,5^{2k-1}J_k/Z_c$.
Squaring the four claimed upper bounds $6,112,4000,190000$ verifies
all four inequalities using rational arithmetic alone.

\subsection{Driven Gaussian jets}
Let $d=v\lambda$, $r=y-q$, and let $u$ denote the normalized coherent
Gaussian in $G=\varphi(\xi)e^{-it/(2\mu)}u$. The start of every
characteristic in the envelope support precedes the drive. Thus the
characteristic solution depends on $t+\xi/d$, and differentiation in
$\xi$ is time differentiation divided by $d$ on $q,q',\vartheta$.
The forcing $\mathcal F=\Omega^2b+b''$ and the error $e=q-b$ obey
\[
 |e|\le\frac{|1-\lambda^2|\,2\|b'\|_\infty}{\lambda\Omega}<.001,
 \qquad
 |e'|\le\frac{|1-\lambda^2|\,2\|b'\|_\infty}{\lambda}<.07.
\]
Indeed the total variation of the single positive hump $b'$ is
$2\|b'\|_\infty$; integrate the sine and cosine representations
of $e,e'$ against $b''(\lambda t+\xi/v)$. It follows that
$|q|<25$, $|q'|<1630$, $|q''|<170000$, and
$|q'''|<36000000$, using
$q''=b''-\Omega^2e$ and $q'''=\lambda b'''-\Omega^2e'$.
Moreover $|\mathcal F|<421000$ and
$|\dot{\mathcal F}|<5.24\times10^7$.

Define the complex coefficient $a_1$ and real coefficient $\beta$ by
\[
 a_1=\frac{q'+i\mu q''}{d},\qquad
 \beta=\frac{\mu}{2d}(qq''-q'^2+\mathcal F e).
\]
Direct differentiation, including the scalar compensation phase, gives
\begin{align*}
 u_\xi&=(a_1r+i\beta)u,\\
 u_{\xi\xi}&=\left[a_1^2r^2+
 ((a_1)_\xi+2i\beta a_1)r+
 i\beta_\xi-a_1q_\xi-\beta^2\right]u,\\
 (a_1)_\xi&=(q''+i\mu q''')/d^2,\\
 \beta_\xi&=\frac{\mu}{2d^2}
 (qq'''-q'q''+\dot{\mathcal F}e+\mathcal F e').
\end{align*}
Since $d\ge126.99873$, substitution of the preceding bounds yields
\[
 |a_1|<19,\quad |(a_1)_\xi|<25,\quad |\beta|<273,
 \quad |\beta_\xi|<367,\quad |q_\xi|<13,
 \quad \mu|q'|<17.
\]
For example the numerator bounding $\beta$ is
$\mu_+(25\cdot170000+1630^2+421000\cdot.001)$;
for $\beta_\xi$ it is
$\mu_+(25\cdot36000000+1630\cdot170000
+5.24\,10^7\cdot.001+421000\cdot.07)$.
The denominators are respectively $2d_-$ and $2d_-^2$.

For the Gaussian measure $|u|^2\dd r$,
$\|r^k u\|^2=(2k-1)!!/2^k$. The quadratic polynomial in
$u_{\xi\xi}$ therefore has coefficient moduli bounded by
$361,10399,75143$. Applying the triangle inequality to these three
monomials and to their derivatives proves
\[
 \begin{array}{c|rr}
 &u_\xi&u_{\xi\xi}\\\hline
 \|\cdot\|&287&83000\\
 \|r\,\cdot\|&211&63000\\
 \|p_y\,\cdot\|&5110&1485000
 \end{array}
\]
as strict upper bounds. Here
$p_y(P(r)u)=[-iP'(r)+(ir+\mu q')P(r)]u$, which explains the
last row without omitting the phase momentum. Leibniz' rule now gives
\begin{align*}
 \|G_{\xi\xi}\|&<112+12(287)+83000=86556<110000,\\
 \|p_yG_{\xi\xi}\|&<112(18)+12(5110)+1485000<1.55\,10^6,\\
 \|yG_{\xi\xi}\|&<25(86556)+112/\sqrt2+12(211)+63000
                         <2.23\,10^6.
\end{align*}
These explicit calculations prove both weighted bounds used in
Lemma~\ref{pa:lem:aux}.

\subsection{Rotor enclosures and the quiet-prefix current}
For completeness, an outward square-root operation on a nonnegative
rational $x$ can be defined using integers alone. Set $n=10^{32}$,
$k=\lfloor\sqrt{\lfloor n^2x\rfloor}\rfloor$, and return $k/n$ if
$k^2\ge n^2x$, otherwise $(k+1)/n$. This is an upper enclosure whose
error is at most $10^{-32}$. Applying it to the four rational squares
in Lemma~\ref{ph:lem:rotor} gives
\[
\begin{aligned}
 d_0&<.006517177225775,&\quad \epsilon_R&<.000715486047289,\\
 d&<.007232663273064,&q&<.003982026605165.
\end{aligned}
\]
The strict comparisons with $D=.007233$ and $Q=.003983$ follow.
The trigonometric estimates used there also have elementary rational
checks: alternating Taylor sums through degrees $14,16$ bound
$\cos(1.3)$, and degrees $15,17$ bound $\sin(1.3)$.
They give $\cos(1.3)>1/4$, $\tan(1.3)<4$, and
$\sec(1.3)+\tan(1.3)<8$. Since the positive exponential Taylor
sum through degree $14$ at $2.1$ exceeds $8$,
\[
 \int_0^{1.3}\sec^3 t\dd t
 =\tfrac12[\sec(1.3)\tan(1.3)+
                 \log(\sec(1.3)+\tan(1.3))]<8.
\]
The same cosine upper sum at $.998(1.22)$ gives
$\cos(.998(1.22))/2+.01<1/5$, placing both compact waves
strictly inside their respective soft-classifier plateaux.

The Sobolev constants in the quiet-prefix estimate require no private
auxiliary result. For a normalized sector rotor $r_i$ with
$0\le V_i\le g_iK/8$, energy and graph-norm conservation from the
initial flat wave give
\[
 \|r_i'\|^2\le mg_iK/4,\qquad
 \|r_i''\|\le mg_iK/2.
\]
Thus $\|r_i'\|\le K/h$ and $\|r_i''\|\le K^2/h^2$.
On a circle of length two,
$\|f\|_\infty^2\le\|f\|^2/2+2\|f\|\|f'\|$.
Apply this first to $r_i$ and then to $r_i'$ to obtain
$\|\rho_r\|_\infty\le3K/h$ and $\|j_r\|_\infty\le8K$.
The factorized pretrigger wave errors then give
\[
 \|\delta j_X\|_1\le4h\epsilon_{\rm pre},\qquad
 \|\delta\rho\|_1\le3\epsilon_{\rm pre},
\]
and hence the lifted-rank integrand is at most
$36K\epsilon_{\rm pre}$. Integrating up to $1.22$ is bounded by
$120K\epsilon_{\rm pre}<6\,10^{-18}$ as used in the main text.

\subsection{Gaussian tails and the whole holding interval}
After capture, every comparison centre is within $.012$ of its intended
centre, including the offset $z_0$. For $a>0$, integration by parts gives
\[
 \|\one_{|r|\ge a}\gamma\|^2\le
 \frac{e^{-a^2}}{a\sqrt\pi},\qquad
 \|\one_{|r|\ge a}r\gamma\|^2\le
 \frac{e^{-a^2}}{\sqrt\pi}\left(a+\frac1{2a}\right).
\]
Take $a=7.98$. The positive Taylor sum of $e^{a^2}$ through degree
$200$ exceeds $4\,10^{27}$. With $\mu|q'|<.000701$, these inequalities
give comparison tail norms below $10^{-14}$ and momentum norms below
$10^{-13}$ on all four holding bands. They also imply the weaker
wrong-inner-region amplitude bound $10^{-12}$ at capture.

For each inner/outer pair choose a smooth cutoff changing from zero
to one across the intervening band of width two, with
$|\chi'|\le15/16$. Any path that starts in the inner region and
leaves its outer region has cutoff variation at least one. Equivariance
therefore bounds its wave probability by
$\int_{t_b}^3\int |\chi'(Z)j_Z|\dd q\dd t$.
This is a band integral of absolute current; no pointwise boundary
trace estimate is needed. Expanding the exact auxiliary wave as
$G+E$ bounds its integrand by the comparison current plus
$(e_A10^{-13}+10^{-14}p_A+e_Ap_A)/\mu_-$.
Summing the four cutoffs, multiplying by the actual-law cap $16$,
and using $3-1.28=43/25$ proves the displayed holding bound in Part I.


\section{Preparation estimates and transfer to the enlarged Hamiltonian}
\label{prep:main}
All momenta in this appendix have the constant Galilean carrier removed.
A superscript $g$ additionally removes the bounded preparation phase
$\theta$. Thus an ordinary clock momentum here means $p_S-Mv$ in the
laboratory frame. We give the finite coefficient construction as well as
its numerical enclosures; the latter are consequences of the construction,
not extra hypotheses.

\subsection{Coefficient conventions and the Gaussian residual}
Write $T=0.104$, $T_0=0.03$, $L_i=10^{-4}$, $L_f=10^{-2}$,
$M=10$, $v=1$, $K=100$, and $\lambda=2^{52}T_0/L_i^2$, so that
$m=\hbar\lambda$. Lengths and times in this appendix are in metres and
seconds. Let $b_p$ denote the preparation dilation, to distinguish it
from the writer displacement. Set $a=b_p'/b_p$, $f_*=0.201$,
$F_*=0.080802$, $g_1=0.0005$, $g_2=0.00025$ and $\gamma_*=141$.
Here $F(r)=2\int_0^r f$ is the bounded preparation profile; the subscript
on $F_*$ in this subsection is unrelated to the writer force.
Here $b_0(s)=1+99B_9(s/0.1)$ has its constant extensions, and
$b_p$ is its convolution with the fixed $10^{-5}$-scale mollifier.
The explicit clamped ninth-degree polynomial and its fixed positive
mollifier give the derivative bounds
\begin{align}
 B_1&=99(315/128)/0.1,& B_2&=99(2520/64)/0.1^2,\\
 B_3&=99(7560/16+5040/64)/0.1^3,&
 B_4&=99(1360800)/0.1^4. \label{prep:bjets}
\end{align}
The sharper relative bounds $b_p'\le140b_p$ and
$|b_p''|\le24000b_p$ follow from polynomial positivity before
convolution. One finite verification is to express each of
$140b_0-b_0'$, $24000b_0-b_0''$, $24000b_0+b_0''$ in Bernstein
form on dyadic subintervals of $[0,1]$; bisect a subinterval if any
coefficient is negative. All leaves are nonnegative by depth five.
For reproducibility, if $P(x)=\sum_{j=0}^n p_jx^j$, its Bernstein
coefficient of index $k$ on $[u,u+w]$ is
\[
 \sum_{j=0}^k\frac{\binom{k}{j}}{\binom{n}{j}}
 \sum_{l=j}^n p_l\binom{l}{j}u^{l-j}w^j.
\]
Positive convolution preserves these inequalities. Put
\begin{align}
 a_0&=140,\\
 a_1&=24000+a_0^2=43600,\\
 a_2&=B_3+3a_0(24000)+2a_0^3=70141750,\\
 a_3&=B_4+4a_0B_3+3(24000)^2+12a_0^2(24000)+6a_0^4
       =1387431060000. \label{prep:ajets}
\end{align}
These bound $|a^{(j)}|$. Let $B_5$ be the largest absolute Bernstein
coefficient on $[0,1]$ of $d^5b_0/ds^5$ and set
\[
 a_4=B_5+5a_0B_4+10(24000)B_3+20a_0^2B_3
       +30a_0(24000)^2+60a_0^3(24000)+24a_0^5.
\]
The normalized clock envelope has derivative norms bounded by
\begin{align*}
 p_0&=1,&p_1&=q\sqrt{16/7}\,\pi/(2\sigma),&
 p_2&=q\sqrt{512/35}\,[\pi/(2\sigma)]^2,\\
 p_3&=q\sqrt{1024/7}\,[\pi/(2\sigma)]^3,&
 p_4&=q(2\pi/\sigma)^4,
\end{align*}
where $\sigma=0.001$, $q=1.000001$. The convolution scale is
$10^{-10}$ and $q(1-10^{-10}p_1/q)>1$, which proves the normalization
allowance. These are weak-derivative estimates for the compact envelope,
not a band-limit assertion.

Let $A=2r\exp[-r^2/(2L^2)]/(\pi^{1/4}L^{3/2})$,
$L=L_ib_p(S)$, and $G=\phi(S-S_c-vt)A$. Set $z=(r/L)^2$,
$M_j=(2j+1)!!/2^j$ and $\nu_j=\sqrt{M_{2j}}$.
Thus $\|z^jA\|=\nu_j$. For a polynomial $P=\sum c_jz^j$ define
$\mathcal N(P)=\sum|c_j|\nu_j$. This coefficient norm is deliberately
subadditive; signs of independently bounded derivatives cannot cancel.
Use
\[
 H_1=z-\tfrac32,\quad H_2=z^2-5z+\tfrac94,\quad
 H_3=z^3-\tfrac{21}2z^2+\tfrac{79}4z-\tfrac{27}8.
\]
Put $k=B_1$, $k_1=B_2+B_1^2$, $k_2=B_3+3B_1B_2+2B_1^3$, and
\begin{align*}
 C&=\lambda L_i^2(100B_2+B_1^2)/2,\\
 C_k&=\lambda L_i^2(B_1B_2+B_1^3)/2,\\
 C_2&=\lambda L_i^2(100B_3+3B_1B_2+2B_1^3)/2,\\
 C_3&=\lambda L_i^2(100B_4+4B_1B_3+3B_2^2
                         +12B_1^2B_2+6B_1^4)/2.
\end{align*}
The three nonnegative coefficient bounds
\begin{align*}
 r_0&=p_2+3kp_1+\tfrac32 k_1+\tfrac94k^2,\\
 r_1&=2kp_1+k_1+5k^2+2Cp_1+3C_k+C_2,\\
 r_2&=k^2+2C_k
\end{align*}
give the residual norm coefficient $R_0=\sum_{j=0}^2r_j\nu_j$.
For its radial derivative set
\[
 R_r=\sum_{j=0}^2r_j\bigl[2j\sqrt{M_{2j-1}}+\sqrt{M_{2j+1}}\bigr],
\]
where the first summand is zero at $j=0$. This follows equally by the
unitary radial reduction of the three-dimensional Gaussian gradient;
its integration-by-parts identity includes the reduced amplitude's
factor $r$. For the clock derivative set
\begin{align*}
 U_1&=k\mathcal N(H_1),\\
 U_2&=k_1\mathcal N(H_1)+k^2\mathcal N(H_2),\\
 U_3&=k_2\mathcal N(H_1)+3kk_1\mathcal N(H_2)+k^3\mathcal N(H_3),\\
 R_S&=p_3+3p_2U_1+3p_1U_2+U_3
       +2p_2C\nu_1+4p_1kC\mathcal N(zH_1)+3p_1C_2\nu_1\\
 &\quad+2C[k_1\mathcal N(zH_1)+k^2\mathcal N(zH_2)]
            +3kC_2\mathcal N(zH_1)+C_3\nu_1.
\end{align*}
These are obtained by differentiating the complete gauged residual
\begin{equation}
 \frac{R^g}{\hbar}=\frac{\hbar}{2M}
 [G_{SS}+2i\theta_SG_S+i\theta_{SS}G]
 +ia\phi[(f-r)A_r+(f'-1)A/2]+(Q-Q_c)G/\hbar.
 \label{prep:residual}
\end{equation}
In particular, no derivative of a restored phase is added to this
\emph{gauged} residual.

Here are all coefficients needed to include its tails and the weak trap.
Write $\omega_i=(\lambda L_i^2)^{-1}$, $q_2=m\omega_i^2/2$,
$q_0=3\hbar\omega_i/2$, $q_{2S}=4a_0q_2$,
$q_{0S}=2a_0q_0$, $r_c=0.201$, $c_1=5000$, and
$\epsilon_{20}=10^{12}/7^{100}$. Set
\begin{align*}
 Q_0&=q_2r_c^2+q_0,&Q_r&=2q_2r_c+c_1Q_0,\\
 Q_S&=q_{2S}r_c^2+q_{0S}+c_1Q_0,&
 Q_t&=q_2L_f^2+q_0,& Q_{tS}&=q_{2S}L_f^2+q_{0S}.
\end{align*}
The complete residual rates are
\begin{align}
 d&=\hbar R_0/(2M)
 +(\lambda\omega_i^2L_f^2/2+3\omega_i/2+3a_0)\epsilon_{20},\\
 e_r&=\hbar^2R_r/(2ML_i)
 +[c_1Q_t+2q_2L_f+Q_t/L_i
       +\hbar a_0(3/0.2+7/L_i+2000)]\epsilon_{20},\\
 e_S&=\hbar^2R_S/(2M)
 +[Q_{tS}+Q_t(p_1+5a_0/2)
      +\hbar\{3(a_1+a_0p_1)+19a_0^2/2\}]\epsilon_{20}.
 \label{prep:source-rates}
\end{align}
The Gaussian tail recurrence
$\int_{20}^\infty y^n e^{-y^2}dy
=20^{n-1}e^{-400}/2+(n-1)\int_{20}^\infty y^{n-2}e^{-y^2}dy/2$
proves the common tail norm used here; $e^2>7$ makes its stated bound
rational.

\subsection{First moments and preparation current}\label{prep:first}
In the gauge the velocity coefficients are $u_r=af$ and
$u_S=v+\delta$, $\delta=ma'F/(2M)$. Put
\begin{align*}
 d_r&=ma_1f_*/M,&d_S&=ma_2F_*/(2M),& u_{rS}&=a_1f_*,\\
 z_r&=Q_r+\tfrac\hbar2(a_0\,4000+ma_2f_*/M),\\
 z_S&=Q_S+\tfrac\hbar2(a_1+ma_3F_*/(2M)).
\end{align*}
If $W_0=\hbar2^{35}/T_0$ and
$W_1=\hbar(2^{52}/256+5\cdot2^{35})/(T_0L_i)$, the remote
activation forces are bounded by
$F_r=W_1+Q_r$, $F_S=25000(W_0+Q_0)+Q_S$.
Let $P_j=\|p_j(\Psi^g-G)\|$ and $D=\|\Psi^g-G\|$.
Differentiating the symmetrized drift gives
\begin{align}
 D(t)&\le dt,\\
 P_r'&\le a_0P_r+d_rP_S+F_rn+e_r+z_rdt,\\
 P_S'&\le u_{rS}P_r+d_SP_S+F_Sn+e_S+z_Sdt,\\
 n'&\le g_0+P_S/(Mg_1),\\
 g_0&=\hbar[p_1+a_0\sqrt{3/2}]/(Mg_1). \label{prep:first-system}
\end{align}
The collar defining $n$ translates at $v+\|\delta\|_\infty$ on the
right and $v-\|\delta\|_\infty$ on the left; its convective contribution
is nonpositive. Both collars vanish on the initial stock. Their final
full-one edges are respectively below $0.1035000001262$ and above
$0.1004999998738$, so that they cover the remote activation and the
left phase support respectively. Thus the force term $F_jn$ is valid
throughout preparation, including the exact wave's tails.

For the sharper clock estimate take $\alpha=10^{-12}$ and
\[
 s=e_r+e_S/K+\alpha g_0,\qquad c=z_r+z_S/K,\qquad E_*=(11/4)^{15}.
\]
The three scaled row sums
$a_0+Kd_r+F_r/\alpha$, $u_{rS}/K+d_S+F_S/(K\alpha)$ and
$\alpha K/(Mg_1)$ are less than $141$. A positive supersolution gives
\begin{align}
 B(t)&=e^{141t}(st+cdt^2/2),&
 P_r(t)&\le B(t),&P_S(t)&\le KB(t),\\
 B_T&=E_*(sT+cdT^2/2),&
 B_I&=E_*(sT^2/2+cdT^3/6),\\
 n_{1T}&=g_0T+KB_I/(Mg_1). \label{prep:anisotropic}
\end{align}
Using $e^{141T}<E_*$ is legitimate since $141T<15$ and $e<11/4$.
These terminal majorants are increasing and apply at every earlier time.

For the smaller radial estimate and the preparation current retain also
the mass-weighted system. Put $\ell=10^{-6}$,
\begin{align*}
 e_0&=e_r/\sqrt m+e_S/\sqrt M,&
 c_0&=z_r/\sqrt m+z_S/\sqrt M,\\
 s_0&=e_0+\ell g_0,&
 Z_T&=E_*(s_0T+c_0dT^2/2),\\
 Z_I&=E_*(s_0T^2/2+c_0dT^3/6),&
 n_w&=g_0T+Z_I/(\sqrt M g_1).
\end{align*}
Indeed the weighted momentum plus $\ell n$ has growth at most $141$:
the drift row bound is
$a_0+\sqrt{m/M}a_1f_*+d_S$, increased by
$\ell/(\sqrt M g_1)$, and
$(F_r/\sqrt m+F_S/\sqrt M)/\ell<141$.
Consequently the ordinary radial error at the handoff is bounded by
\begin{equation}
 P_{r,o}=\sqrt m\,Z_T+ma_0f_*n_w,
 \qquad P_{S,o}=KB_T+ma_1F_*n_{1T}/2,
 \qquad D_o=dT. \label{prep:old-handoff}
\end{equation}
The phase terms here are localized to the reflected collar; no phase is
set to zero on an exact wave tail.

For the Gaussian rank let
$k_r=6/(7L_i)$, $k_S=2a_0$,
$c_r=k_r\sqrt{3/2}/L_i$, $c_S=k_S(p_1+a_0\sqrt{3/2})$.
The exact material-rank identity and the bilinear current expansion give
\begin{align}
 I_o&=(c_r/\lambda+\hbar c_S/M)dT^2/2
       +(k_r/\sqrt m+k_S/\sqrt M)Z_I,\\
 \beta&=mL_i^2(B_1B_2+a_0B_1^2)T/M
                   +Ta_0\,32000/7^{200}. \label{prep:old-current}
\end{align}
The first term of $\beta$ bounds $K_S\delta$ pointwise; the second
bounds the nonquadratic radial tail. The remaining current has absolute
value before integration over any live coordinate. This proves
$I_o<1.974548968\cdot10^{-19}$ and $\beta<2.638660\cdot10^{-13}$.

\subsection{Second spatial moments and the second collar}\label{prep:hessian}
Here is an explicit construction of every second-order residual
coefficient. Continue $H_j$ by
$H_0=1$, $H_{j+1}=(z-3/2)H_j-2zH_j'$, and take absolute
coefficients before adding independently bounded terms. The positive
Gaussian clock polynomials are
\begin{align*}
 A_0&=1,& A_1&=a_0|H_1|,\\
 A_2&=a_1|H_1|+a_0^2|H_2|,&
 A_3&=a_2|H_1|+3a_0a_1|H_2|+a_0^3|H_3|,\\
 A_4&=a_3|H_1|+(4a_0a_2+3a_1^2)|H_2|
                       +6a_0^2a_1|H_3|+a_0^4|H_4|.
\end{align*}
Here $|P|$ means coefficientwise absolute value. To evaluate radial
norms, form
$P_{j,r}(y)=(\partial_y-y)^r y^{2j+1}$ and set
\[
 \mathcal R_r(j;s)=L_i^{-r}
       \sum_e |[y^e]P_{j,r}|\sqrt{M_{e-1+s}},\qquad M_{-1}=2.
\]
Only indices at least $-1$ occur. For a factor $F/L^2$ and $r\le2$
replace this by
\[
 \mathcal R_r^F(j)=\sum_{l=0}^r\binom rl c_lL_i^{-l}
                         \mathcal R_{r-l}(j;2-l),
 \qquad(c_0,c_1,c_2)=(1,2,2).
\]
This is the product rule using $F\le r^2$, $|F'|\le2r$,
$|F''|\le2$. Define
\[
 J_{n,r}=\sum_{j=0}^n\binom nj p_{n-j}
       \sum_l[z^l]A_j\,\mathcal R_r(l;0),
\]
and $J_{n,r}^F$ by replacing $\mathcal R_r$ with $\mathcal R_r^F$.
With $C_j^*=\lambda L_f^2a_j/2$ for $1\le j\le4$, the sources are
\begin{align*}
 e_{rr}&=\frac{\hbar^3}{2M}
       (J_{2,2}+2C_1^*J_{1,2}^F+C_2^*J_{0,2}^F)+10^{-100},\\
 e_{rS}&=\frac{\hbar^3}{2M}
       (J_{3,1}+2C_1^*J_{2,1}^F+3C_2^*J_{1,1}^F+C_3^*J_{0,1}^F)
       +10^{-100},\\
 e_{SS}&=\frac{\hbar^3}{2M}
       (J_{4,0}+2C_1^*J_{3,0}^F+5C_2^*J_{2,0}^F
                           +4C_3^*J_{1,0}^F+C_4^*J_{0,0}^F)+10^{-100}.
\end{align*}
For example the last formula differentiates
$G_{SS}+2i\theta_SG_S+i\theta_{SS}G$ twice and retains coefficients
$2,5,4,1$. The nonquadratic and trap tails after two derivatives have
Gaussian degree at most seven and total frequency coefficient below
$10^{30}$; their momentum contribution is less than
$\hbar^2 10^{30}\epsilon_{20}<10^{-100}$ in each row. This can also be
checked by the general product-rule construction below.

For clarity we spell out the second-force coefficients. Set
$c_2=1.28\cdot10^8$, $q_{t1}=q_{2S}r_c^2+q_{0S}$,
$q_{t2}=(4a_1+16a_0^2)q_2r_c^2+(2a_1+4a_0^2)q_0$. Then
\begin{align*}
 Q_{rr}&=2q_2+4c_1q_2r_c+c_2Q_0,\\
 Q_{rS}&=c_1Q_r+2q_{2S}r_c+c_1q_{t1},\\
 Q_{SS}&=c_2Q_0+2c_1q_{t1}+q_{t2}.
\end{align*}
Put $\gamma_1=25000$,
$\gamma_2=(2520/64)/(0.0001)^2$, and
\begin{align*}
 W_2&=\frac{\hbar}{T_0L_i^2}
 [2^{52}(1+(15/16)2^{17})+5\cdot2^{52}/128+10^5\cdot2^{35}],\\
 F_{rr}&=W_2+Q_{rr},\\
 F_{rS}&=\gamma_1(W_1+Q_r)+Q_{rS},\\
 F_{SS}&=\gamma_2(W_0+Q_0)+2\gamma_1Q_S+Q_{SS}.
\end{align*}
Let $v_{jr}$ and $v_{jS}$ denote the component derivatives of
$(af,ma'F/(2M))$. The second-derivative ceilings used below are
\[
\begin{array}{c|ccc}
 &rr&rS&SS\\\hline
 u_r&a_0\,4000&a_1&a_2f_*\\
 \delta&ma_1/M&ma_2f_*/M&ma_3F_*/(2M)\\
 \operatorname{div}u&a_0\,10^7+ma_2/M&
 a_1\,4000+ma_3f_*/M&a_2+ma_4F_*/(2M)
\end{array}
\]
Write $d^{\rm div}_r=a_0\,4000+ma_2f_*/M$ and
$d^{\rm div}_S=a_1+ma_3F_*/(2M)$ for the first divergence
ceilings. If $U_{jk},V_{jk},D_{jk}$ denote the three rows of this table,
all scalar coefficients of the Hessian comparison are
\begin{align*}
 e_2&=e_{rr}+e_{rS}/K+e_{SS}/K^2,\\
 c_B&=\hbar(U_{rr}+KV_{rr}+d^{\rm div}_r)+2Q_r\\
 &\quad+\{\hbar[U_{rS}+KV_{rS}+(Kd^{\rm div}_r+d^{\rm div}_S)/2]
                  +KQ_r+Q_S\}/K\\
 &\quad+\{\hbar[U_{SS}+KV_{SS}+Kd^{\rm div}_S]+2KQ_S\}/K^2,\\
 c_D&=\tfrac{\hbar^2}{2}(D_{rr}+D_{rS}/K+D_{SS}/K^2)
                  +\hbar(Q_{rr}+Q_{rS}/K+Q_{SS}/K^2),\\
 c_n&=\hbar(F_{rr}+F_{rS}/K+F_{SS}/K^2)
                  +2F_r\hbar/(g_2K)+4F_S\hbar/(g_2K^2),\\
 c_A&=2F_r+2F_S/K.
\end{align*}
These follow directly by commuting two momenta through
$\{u_j,p_j\}/2+Q_c+V_A$. For example the mixed commutator contributes
$\hbar\sum_l\|u_{l,jk}\|P_l$ and
$\hbar(\|\partial_j\operatorname{div}u\|P_k+
\|\partial_k\operatorname{div}u\|P_j)/2$; both occur above.
The three homogeneous scaled row sums are below $282$.

With $Z_2=\max(\|p_r^2E^g\|,\|p_rp_SE^g\|/K,
\|p_S^2E^g\|/K^2)$, nested clock cutoffs give
\[
 Z_2'\le282Z_2+c_A\sqrt{n_1}\sqrt{Z_2}
                   +e_2+c_BB+c_Ddt+c_nn_1.
\]
Indeed a cutoff supported where the first collar is one satisfies
$\|\chi p_rE^g\|\le\sqrt{n_1\|p_r^2E^g\|}$ and
$\|\chi p_SE^g\|\le\sqrt{n_1\|p_S^2E^g\|}+2\hbar n_1/g_2$.
These inequalities are integration by parts, not finite propagation.
Let $b_n=K/(Mg_1)$, $\gamma=141$, and define
\begin{align*}
 J_n&=\frac{\sqrt{g_0\pi}}{2\gamma^{3/2}}
       +\frac{\sqrt{b_ns/2}}{(\gamma/2)^2}
       +\frac{3\sqrt{\pi b_ncd/6}}{4(\gamma/2)^{5/2}},\\
 J_f&=\frac{e_2}{2\gamma}
       +c_B(s/\gamma^2+cd/\gamma^3)
       +\frac{c_Dd}{(2\gamma)^2}\\
 &\quad+c_n\{g_0/(2\gamma)^2+b_n(s/\gamma^3+cd/\gamma^4)\},\\
 Z_{2T}&=e^{2\gamma T}(c_AJ_n/2+\sqrt{J_f})^2,
 \qquad H_{SS}=K^2Z_{2T}. \label{prep:hessian-construction}
\end{align*}
For the last display, extend positive integrals to infinity after using
$n_1(t)\le g_0t+b_ne^{\gamma t}(st^2/2+cdt^3/6)$.
The comparison follows by setting $y=e^{-\gamma t}\sqrt{Z_2}$:
the sum of the separate positive supersolutions for $y'\le A+B/(2y)$
is a supersolution, with zero handled by regularization.
The factor $e^{\gamma T}$ is enclosed by its 90-term Taylor sum and
the remaining positive geometric tail with ratio $\gamma T/91$; its
upper enclosure is squared when evaluating $e^{2\gamma T}$.

The second collar amplitude and the physical phase restoration are
\begin{align}
 n_{2T}&=\frac{T}{Mg_2}
           [\sqrt{n_{1T}H_{SS}}+2\hbar D_o/g_1],\\
 L_S&=\sqrt{n_{2T}H_{SS}}+2\hbar n_{1T}/g_2,
 &L_r&=\sqrt{n_{2T}H_{SS}/K^2},\\
 H_{SS}^{o}&=H_{SS}+2q_1L_S+(q_1^2+\hbar q_2^*)n_{2T},
 &q_1&=ma_1F_*/2,\quad q_2^*=ma_2F_*/2.
 \label{prep:second-collar}
\end{align}
Both reflected second collars fit before the relevant remote support:
their final edges are below $0.103750000127$ and above
$0.100249999873$. These same increasing majorants hold for all preceding
times. In particular the ordinary derivatives on the writer support
are bounded by $L_S,L_r$, because both the helper and the phase vanish
there.

\subsection{Fourth spatial closure and its complete coefficient ledger}
\label{prep:fourth}
Only ordinary spatial derivatives are used here. There is no assumption
that the unmollified radial reference lies in the domain of a fourth
power of its Hamiltonian. Use the commuting scaled momenta
$D_r=p_r$, $D_S=p_S/K$, and
$Z_j=\max_{a+b=j}\|D_r^aD_S^bE^g\|$.

The higher coefficient construction is finite. Define $q_j=b_p^{(j)}/b_p$.
Starting with $P_0(q)=q_1$, form
\[
 P_{j+1}=\sum_{l=1}^6(q_{l+1}-q_1q_l)\partial_{q_l}P_j,
 \qquad 0\le j<6.
\]
Evaluate the absolute coefficients at
\[
 (\bar q_1,\ldots,\bar q_7)
 =(140,24000,B_3,B_4,B_5,4B_5/10^{-5},300B_5/10^{-10})
\]
to obtain $\bar a_j\ge|a^{(j)}|$. For $j\le4$ these agree with
the bounds above. The last two entries use convolution of $b_0^{(5)}$
with the first and second derivative of the mollifier. Similarly set
$p_5=4p_4/10^{-10}$, $p_6=300p_4/10^{-20}$.
The normalized unit bump has derivative polynomials generated by
\begin{equation}
 U_0(x,u)=1,\qquad
 U_{j+1}=\partial_xU_j+2xu^2\partial_uU_j-2xu^2U_j,
 \qquad u=(1-x^2)^{-1}. \label{prep:bump-recurrence}
\end{equation}
Its normalization is greater than $1/4$; hence its $j$th derivative
supremum is bounded by
$4\sum_{l,k}|[x^lu^k]U_j|\,k^k/(8/3)^k$, with the $k=0$
factor defined as one. This proves the quoted mollifier bounds (the
first derivative also uses unimodality), as well as all tent jets below.

Let the positive partial Bell coefficients be
\[
 \mathcal B_{0,0}=1,\qquad
 \mathcal B_{n,k}=\sum_{j=1}^{n-k+1}\binom{n-1}{j-1}
                   \bar a_{j-1}\mathcal B_{n-j,k-1}.
\]
All unspecified entries are zero. For $n=0$ put $\mathcal A_0=1$;
otherwise $\mathcal A_n=\sum_{k=1}^n\mathcal B_{n,k}|H_k|$.
For $\varepsilon=0,1$, define the computable Gaussian norm majorant
\[
 \mathcal J_{n,r}^{(\varepsilon)}=
 \sum_{j=0}^n\binom nj p_{n-j}
 \sum_l[z^l]\mathcal A_j\,L_i^{-r}
 \sum_e\left|[y^e](\partial_y-y)^r y^{2l+1+2\varepsilon}\right|
                                   \sqrt{M_{e-1}}.
\]
For the higher-core calculation one may use $\pi\le22/7$ in $p_j$
and $\hbar\le\bar h=1.055\cdot10^{-34}$. The fourth physical source
coefficient is bounded by
\begin{align}
 s_4&=\frac{\bar h^5}{2M}\sum_{n=0}^4K^{-n}\Bigl[
 \mathcal J_{n+2,4-n}^{(0)}
 +\sum_{j=1}^{n+2}c_{n,j}\frac{\lambda L_f^2\bar a_j}{2}
                              \mathcal J_{n+2-j,4-n}^{(1)}\Bigr],\\
 c_{n,j}&=2\binom n{j-1}+\binom n{j-2},
 \label{prep:fourth-core}
\end{align}
where a binomial coefficient outside its range is zero. Direct rational
arithmetic and outward square roots give $s_4<7.141\cdot10^{-127}$.
The physical phase-square cancellation in the gauge is essential to
this small coefficient.

We next give a complete way to check the nonquadratic tail bound, including
higher derivatives of $F-r^2$. Let $b_j$ denote the bump derivative
suprema just constructed and use
\[
\begin{gathered}
 (f_0,f_1,f_2,f_3,f_4,f_5)
 =(0.201,1,4000,10^7,2b_2/0.001^3,2b_3/0.001^4),\\
 F_0=0.080802,\qquad F_j=2f_{j-1}\quad(1\le j\le6).
\end{gathered}
\]
The disjoint jumps of the unsmoothed tent give these bounds. The scaled
velocity derivative sum at order $j$ is explicitly
\begin{equation}
 V_j=\max_{0\le k\le j}\frac{\bar a_k f_{j-k}}{K^k}
   +\max_{0\le k\le j}\frac{K\bar h\lambda\bar a_{k+1}F_{j-k}}
                                      {2MK^k}.
 \label{prep:velocity-jets}
\end{equation}
It gives $V_2<560001$, $V_3<1.400001\cdot10^9$,
$V_4<3.669751\cdot10^{13}$, $V_5<2.786850\cdot10^{18}$.

For an entirely algebraic tail check replace every Gaussian norm in
$\mathcal J_{n,r}^{(0)}$ by the absolute sum of its polynomial
coefficients; call the result $\mathcal C_{n,r}$. Keep powers of $y$
as formal factors. The derivatives of $F-r^2$ have coefficient ceilings
$\Delta F=(2,2,4,2f_2,2f_3)$, and those of $f-r$ have
$\Delta f=(2,2,f_2,f_3,f_4,f_5)$. The unit logistic cutoff has
order-$j$ derivative bounded by
\[
 t_j=8j!\sum_{k=1}^j\binom{j-1}{k-1}2^k
                 \frac{(j+k)^{j+k}}{(8/3)^{j+k}}\quad(1\le j\le4).
\]
To prove it on $x\le1/2$, put $u=1/x$: the exponent is at most
$-u+2$, its $j$th derivative at most $2j!u^{j+1}$, and apply the
Bell product rule and $\max u^ne^{-u}\le n^n/(8/3)^n$.
Symmetry handles the other half. Set $t_0=1$ and
$\tau_j=t_j/0.001^j$.
For $p=2,4$ define
$J_p(0)=1$, $J_p(n)=\sum_{k=1}^n\mathcal B_{n,k}p^k$.
With $A_2=4\cdot10^{-41}$ and $A_0=1.2\cdot10^{-48}$, put
\begin{align*}
 q(n,r)&=A_2J_4(n)(L_f^2,2L_f,2)_r
                      +\one_{r=0}A_0J_2(n),\\
 \widehat q(n,r)&=q(n,r)+
   \sum_{c=0}^n\sum_{e=0}^r\binom nc\binom re
                      \tau_c\tau_e q(n-c,r-e).
\end{align*}
The three-entry sequence is zero outside $r=0,1,2$.
A bound for the sum of all fourth-derivative tail frequency coefficients
is the following finite positive sum, with $r=4-n$:
\begin{align}
 \mathcal T={}&\sum_{n=0}^4\sum_{a=0}^n\sum_{b=0}^{r}
 \binom na\binom rb\Bigl[
 \frac{1.425\cdot10^{-12}}{2M}\bar a_{a+1}\Delta F_b
                                      \mathcal C_{n-a+1,r-b}\\
 &+\frac{1.425\cdot10^{-12}}{4M}\bar a_{a+2}\Delta F_b
                                      \mathcal C_{n-a,r-b}
 +\bar a_a\Delta f_b\mathcal C_{n-a,r-b+1}\\
 &+\tfrac12\bar a_a\Delta f_{b+1}\mathcal C_{n-a,r-b}
 +10^{34}\widehat q(a,b)\mathcal C_{n-a,r-b}\Bigr]<10^{70}.
 \label{prep:tail-coefficient}
\end{align}
This is simply the product rule applied respectively to the two phase
terms, transport tail, and $Q-Q_c$ in \eqref{prep:residual}.
Taking absolute coefficients before forming $\mathcal C$ prevents
cancellation between independently bounded clock jets. As a rough
independent size check, $p_j\le10^{9j}$,
$\bar a_j<10^3 10^{9j}$, and Gaussian mixed coefficients through order
six are below $10^{12}10^{9j}$. The transport term has at most five
derivative slots and fewer than $10^5$ product summands, giving
$10^{3+3+12+45+5}=10^{68}$. In the phase term $\hbar\theta$ replaces
the large phase coefficient by the physical mass: six derivative slots
give at most $10^{-12+3+3+12+54+5}=10^{65}$. The pure $G_{SS}$
term is smaller. The explicit positive sum above also checks the weak
potential term. No nonzero tail is discarded.

The degree is at most twelve. Integration by parts gives
\[
 \|\one_{y\ge20}y^{12}A\|^2
 \le\frac{2\cdot20^{25}}{1-25/800}e^{-400}
 <10^{40}/7^{200}.
\]
Thus the additional fourth physical source is
$s_{\rm tail}=\bar h^4 10^{90}/7^{100}<10^{-130}$.
The same finite product rule for the quiet potential gives, for total
order $j\le4$,
\begin{equation}
 Q_j\le10^{-40}10^{9j}. \label{prep:quiet-high}
\end{equation}
For an explicit verification use the preceding $J_p(n)$ and cutoffs,
replace $(L_f^2,2L_f,2)$ by $(1,1,2)$ on the cutoff support,
and evaluate
$\sum_{a=0}^n\sum_{b=0}^r\binom na\binom rb
\tau_a\tau_b[A_2J_4(n-a)(1,1,2)_{r-b}
+\one_{r=b}A_0J_2(n-a)]$ for $n+r=j$.

For the remote potential define $W_j=\bar h\,w_j/(T_0L_i^j)$, where
\begin{align*}
 w_0&=2^{35},\qquad w_1=2^{52}/256+5\cdot2^{35},\\
 w_2&=2^{52}(1+(15/16)2^{17})+5\cdot2^{52}/128+10^5\cdot2^{35},\\
 w_3&=2^{112},\qquad w_4=2^{160}.
\end{align*}
Set
\begin{align*}
 (\gamma_0,\ldots,\gamma_4)&=\left(1,25000,
 \frac{2520/64}{10^{-8}},
 \frac{7560/16+5040/64}{10^{-12}},
 \frac{1360800}{10^{-16}}\right),\\
 R_j&=\max_{0\le k\le j}K^{-k}
 \left[\gamma_kW_{j-k}+
       \sum_{l=0}^k\binom kl\gamma_lQ_{j-l}\right].
\end{align*}
These bounds follow by differentiating the displayed capped radial
potential and activation polynomial; the last two deliberately loose
bounds require only its bounded fourth weak derivative. All remote
supports lie at $S\ge0.104$.

Add two smooth analytic cutoffs of width $g=10^{-4}$, with final
transition intervals $(0.103751,0.103851)$ and
$(0.103851,0.103951)$. They fit between the second old collar and
activation. Their first two derivative bounds are $5/g$ and $128/g^2$.
Put $a_c=\bar h/(Kg)$. With $n$ the second-collar amplitude, integration
by parts and interpolation give local derivative bounds
\begin{align}
 N_1&\le\sqrt{nZ_2}+10a_cn,\\
 N_2&\le\sqrt n(\sqrt{Z_4}+18a_c\sqrt{Z_2}),\\
 N_3&\le\sqrt{N_2Z_4}+10a_cN_2. \label{prep:local-high}
\end{align}
For example the squared second-derivative norm is at most
$nZ_4+20a_cn\sqrt{Z_2Z_4}+306a_c^2nZ_2$; its square root is
bounded as displayed since $306<18^2$. The global interpolation
$Z_3\le\sqrt{Z_2Z_4}$ moves one commuting derivative between the
two factors in its squared norm. Odd extension removes the radial
boundary term.

The useful time envelopes, with $Z_{1T}=B_T$, are
\begin{align*}
 Z_1(t)&\le Z_{1T}e^{-\gamma(T-t)},&
 Z_2(t)&\le Z_{2T}e^{-2\gamma(T-t)},\\
 n(t)&\le A_nt^{3/2}e^{\gamma t}+B_nt^2e^{3\gamma t/2}+C_nt^2,\\
 A_n&=\sqrt{K^2Z_{2T}g_0}/(Mg_2e^{\gamma T}),\\
 B_n&=\sqrt{K^2Z_{2T}n_{1T}}/(Mg_2Te^{3\gamma T/2}),\\
 C_n&=2\bar h D_T/(Mg_2g_1T),\qquad D_T=3.006684\cdot10^{-11}.
\end{align*}
Here $D_T$ is used only in the fourth-order majorant, not to replace the
sharper handoff norm in the final probability accounting.
Writing $Z=Z_4$, four commutations yield
\begin{align}
 Z'&\le4\gamma Z+\alpha(t)Z^{3/4}+\beta(t)Z^{1/2}+f(t),\\
 \alpha&=4R_1n^{1/4},\\
 \beta&=4R_1[\sqrt{18a_c}\,n^{1/4}Z_2^{1/4}+10a_cn^{1/2}]
       +6\bar hR_2n^{1/2}+(8\bar hV_2+4Q_1)\sqrt{Z_2},\\
 f&=(720R_1a_c^2+108\bar hR_2a_c)\sqrt{nZ_2}
       +4\bar h^2R_3(\sqrt{nZ_2}+10a_cn)+\bar h^3R_4n\\
 &\quad+(7\bar h^2V_3+6\bar hQ_2)Z_2
       +(3\bar h^3V_4+4\bar h^2Q_3)Z_1
       +(\bar h^4V_5/2+\bar h^3Q_4)dt+s_4+s_{\rm tail}.
 \label{prep:fourth-system}
\end{align}
The factors $8,7,3,1/2$ in the drift terms are the sum of the
first-order drift and its divergence product-rule coefficients.
Putting $y=e^{-\gamma t}Z^{1/4}$ and adding the separate positive
supersolutions gives
\begin{equation}
 Z_4(T)^{1/4}\le e^{\gamma T}
 \left[\tfrac14\int_0^T\alpha e^{-\gamma t}dt
 +\sqrt{\tfrac12\int_0^T\beta e^{-2\gamma t}dt}
 +\left(\int_0^Tf e^{-4\gamma t}dt\right)^{1/4}\right].
 \label{prep:fourth-solution}
\end{equation}
This includes the nonlinear third-derivative coupling; it is not an
exponential multiplying a terminal forcing value.

For completeness all integrals in this expression can be evaluated by
one finite rule. For $0<p\le1$ use subadditivity on the three terms of
$n(t)$; combine with $Z_2(t)^q$ and the integrating factor, and extend
each positive integral to infinity. If its time power is $u\ge0$ and
its decay is $\rho>0$, with $j=\lfloor u\rfloor$, use
\[
 \int_0^\infty t^ue^{-\rho t}dt
 \le j!(j+1)^{u-j}/\rho^{u+1}.
\]
This is H\"older interpolation of adjacent integer moments. All decay
rates in \eqref{prep:fourth-solution} are positive. Integer roots are
enclosed rationally; no numerical gamma function is needed.
For a final far cutoff of width $0.01$ before $S=0.14145$, put
\begin{equation}
 P_4=K^4Z_4(T),\qquad
 L_{SS}=K^2\sqrt{n_{2T}}
 [\sqrt{Z_4(T)}+18\bar h\sqrt{Z_{2T}}/(K\cdot0.01)].
 \label{prep:far-hessian}
\end{equation}
The unshortened positive constructions define the values used in later
interfaces; their shortened displays are $P_4<1.075567\cdot10^{-76}$ and
$L_{SS}<2.534075\cdot10^{-57}$ uniformly through preparation.

\subsection{Original-stock transfer to the writer}\label{prep:transfer}
Let $\Psi_o$ be the exact preparation without the writer, and $\Psi_n$
the enlarged exact wave, with their common initial state
$\Psi_n(0)=\Psi_o(0)g_0$. Here $g_0(Z)=\pi^{-1/4}e^{-Z^2/2}$.
Write $\Delta=\Psi_n-\Psi_og_0$. The ground oscillator energy is
subtracted and the additional potential is
$J=(\hbar/T_0)\Pi_1[C_w(s)-F_w(s)Z]$. Thus
\[
 i\hbar\partial_t\Delta=H_n\Delta+J\Psi_og_0,\qquad\Delta(0)=0.
\]
No comparison flow or initial law is substituted in this identity.
The writer is supported on $S\ge0.14145$. Throughout preparation,
the old helper and its bounded phase both vanish on this support. The source amplitudes and momentum
norms are therefore bounded by $n_{2T},L_S,L_r,L_{SS}$ established above.

The following constants are exact rational majorants. Put $\mu=0.01$,
$w=0.003$ and
\[
 (b_0,b_1,b_2,b_3,b_4)=24\left(1,\frac{315}{128w},
 \frac{2520}{64w^2},\frac{7560/16+5040/64}{w^3},
 \frac{1360800}{w^4}\right).
\]
In this subsection $b_j$ denotes a writer derivative bound, not the
preparation dilation or bump derivative. Define
\begin{align*}
 F_j&=b_j/\mu+\mu b_{j+2}\quad(0\le j\le2),\\
 C_0&=b_0^2/(2\mu)+\mu b_0b_2/2,\\
 C_1&=b_0b_1/\mu+\mu(b_1b_2+b_0b_3)/2,\\
 C_2&=(b_1^2+b_0b_2)/\mu
                  +\mu(b_2^2+2b_1b_3+b_0b_4)/2,\\
 c_0&=b_0+\mu b_1,& c_1&=b_1+\mu b_2.
\end{align*}
These bound $F_w,C_w$ and their dimensionless derivatives. Hereafter
$F_{\rm old}=8\cdot10^{-6}+10^{-17}$ and $F_{{\rm old},r}=10^{-6}$.
These force ceilings follow by differentiating the displayed potential:
the preparation clock force is bounded by
\[
 ma_2(0.080802)/2+ma_0a_1(0.201)^2
       +m^2a_1a_2(0.080802)^2/(4M)<8\cdot10^{-6},
\]
and its radial force by
$ma_1(0.201)+ma_0^2(0.201)+m^2a_1^2(0.080802)(0.201)/(2M)$.
Adding the quiet and activation derivatives gives the stated ceilings;
their clock contribution is below $10^{-17}$.
A second differentiation gives the conservative bound
$F_{{\rm old},SS}\le1$: its preparation part is
\[
 ma_3(0.080802)/2+m(a_1^2+a_0a_2)(0.201)^2
 +m^2(a_2^2+a_1a_3)(0.080802)^2/(4M),
\]
with $Q_{SS}+\gamma_2(W_0+Q_0)+2\gamma_1q_{t1}$ added.

For concise exact formulas set $n_o=n_{2T}$ and
\begin{align}
 d_\Delta&=(C_0+F_0/\sqrt2)n_o/T_0,&
 q_\Delta&=(C_0+\sqrt{3/2}F_0)n_o/T_0,\\
 r_S&=(C_0+F_0/\sqrt2)L_S/T_0
      +\hbar(C_1+F_1/\sqrt2)n_o/T_0^2,\\
 r_r&=(C_0+F_0/\sqrt2)L_r/T_0. \label{prep:transfer-source}
\end{align}
The exact oscillator identities
$\|Zg_0\|=\|p_Zg_0\|=1/\sqrt2$ and
$\|Z^2g_0\|^2=\|p_ZZg_0\|^2=3/4$ prove these coefficients.
Unitary Duhamel and oscillator rotation in the direct sum of its two
quadratures then give increasing bounds
\begin{align}
 D_\Delta(t)&=d_\Delta t,\\
 Q_\Delta(t)&=q_\Delta t+F_0d_\Delta t^2/(2T_0),\\
 P_{\Delta,S}(t)&=r_St+F_{\rm old}d_\Delta t^2/2\\
 &\quad+\frac\hbar{T_0^2}
  [F_1(q_\Delta t^2/2+F_0d_\Delta t^3/(6T_0))
                     +C_1d_\Delta t^2/2],\\
 P_{\Delta,r}(t)&=r_rt+F_{{\rm old},r}d_\Delta t^2/2.
 \label{prep:transfer-bounds}
\end{align}
Here $Q_\Delta$ bounds
$(\|Z\Delta\|^2+\|p_Z\Delta\|^2)^{1/2}$, not an unbounded
coordinate times a bare norm estimate. In the clock inequality the
term $F_1Q_\Delta$ is retained explicitly.

Let $B=\partial_Z+Z$. Its exact equations have sources
$-\Pi_1F_w\Psi_n$ and $-2\Pi_1F_wB\Psi_n$ for $B\Psi_n$ and
$B^2\Psi_n$, respectively. Because $B$ annihilates the old product
wave, integration gives
\begin{align}
 B_1(t)&=F_0(n_ot+d_\Delta t^2/2)/T_0,\\
 B_2(t)&=F_0^2(n_ot^2+d_\Delta t^3/3)/T_0^2,\\
 A_h&=B_1(T)+c_0[n_o+D_\Delta(T)],\\
 A_{2h}&=B_2(T)+2c_0B_1(T)+c_0^2[n_o+D_\Delta(T)].
 \label{prep:annihilator-handoff}
\end{align}
The last quantity bounds the norm of the squared moving annihilator;
it is not the square of its norm. These are operator identities on the
full oscillator Gaussian domain, with no truncation of $Z$.
For $f=B\Psi_n$, oscillator algebra gives
$\|Zf\|^2+\|p_Zf\|^2=\|Bf\|^2+\|f\|^2$.
Consequently, with integrals over $[0,T]$,
\begin{align*}
 P_B={}&F_{\rm old}\int B_1
 +\frac\hbar{T_0^2}\left[F_1\int(B_2+B_1)+C_1\int B_1\right]\\
 &+\frac{F_0}{T_0}[L_ST+TP_{\Delta,S}(T)]
 +\frac{\hbar F_1}{T_0^2}(n_oT+d_\Delta T^2/2),\\
 U_h={}&[P_B+c_0(L_S+P_{\Delta,S}(T))
           +\hbar c_1(n_o+D_\Delta(T))/T_0]/\hbar.
\end{align*}
This proves $\|\partial_S(A\Psi_n)(T)\|\le U_h$.
All displayed integrals are polynomials; for example
$\int B_1=F_0(n_oT^2/2+d_\Delta T^3/6)/T_0$ and
$\int B_2=F_0^2(n_oT^3/3+d_\Delta T^4/12)/T_0^2$.

The exact symbolic entrance values used subsequently are
\begin{align}
 D_h&=3.006683\cdot10^{-11}+D_\Delta(T),\\
 n_h&=n_{1T}+D_\Delta(T),\\
 P_h&=\hbar p_1+P_{S,o}+P_{\Delta,S}(T),\\
 P_{r,h}&=P_{r,o}+P_{\Delta,r}(T)+\epsilon_q,
 \qquad 0\le\epsilon_q<10^{-100},\\
 R_h&=1+Q_\Delta(T)<2. \label{prep:new-handoff}
\end{align}
Here $\epsilon_q$ is the quiet-reference comparison error in the radial
momentum, as bounded in the activation appendix. In particular the
larger norm $3.006684\cdot10^{-11}$ used in the fourth-moment
estimate is \emph{not} substituted for $D_h$.

\subsection{Localized repair of the transfer Hessian}
\label{prep:transfer-hessian}
A coarse second-momentum bound is useful only to generate source-free
left collars. We give it explicitly to avoid circular reuse of the
improved result. Let $P=P_{\Delta,S}(T)$, $D=D_\Delta(T)$,
$Q=Q_\Delta(T)$, and
\begin{align*}
 r_{SQ}&=(C_0+\sqrt{3/2}F_0)L_S/T_0
             +\hbar(C_1+\sqrt{3/2}F_1)n_o/T_0^2,\\
 Q_S&=F_0PT/T_0+F_{\rm old}\int Q_\Delta
       +\frac\hbar{T_0^2}\left[C_1\int Q_\Delta
                  +F_1\int(B_2+2B_1+2D_\Delta)\right]+r_{SQ}T,\\
 r_{SS}(L)&=(C_0+F_0/\sqrt2)L/T_0
       +2\hbar(C_1+F_1/\sqrt2)L_S/T_0^2
       +\hbar^2(C_2+F_2/\sqrt2)n_o/T_0^3,\\
 W_\Delta&=2\hbar T(C_1P+F_1Q_S)/T_0^2
                    +\hbar^2T(C_2D+F_2Q)/T_0^3,\\
 H_c&=r_{SS}(H_{SS})T+2F_{\rm old}PT+\hbar DT+W_\Delta.
\end{align*}
The quadrature estimate $\|Z^2\Delta\|$ and its companion are
controlled by $B_2+2B_1+2D_\Delta$, so that $Q_S$ bounds the
mixed weighted momentum without replacing $Z$ by a bounded operator.
Differentiating the inhomogeneous equation twice proves
$\|p_S^2\Delta\|\le H_c$; the coarse value is below
$\hbar^2(2.486519\cdot10^{27})$.

Choose decreasing left cutoffs with transitions $[0.126,0.136]$ and
$[0.116,0.126]$, width $g=0.01$. The writer forcing vanishes on
both supports and every old clock force is contained in their full-one
regions, because it vanishes for $S\ge0.106$. Positive drift favours
exit from each left region. Continuity and one integration by parts give
\begin{equation}
 n_{\Delta1}=TP/(Mg),\qquad
 n_{\Delta2}=\frac{T}{Mg}
       [\sqrt{n_{\Delta1}H_c}+2\hbar D/g].
 \label{prep:delta-collars}
\end{equation}
For the new Hessian $H(t)=\|p_S^2\Delta(t)\|$ the old-force
momentum is instead localized by
$\|\one_{\rm oldforce}p_S\Delta\|
\le\sqrt{n_{\Delta2}H(t)}+2\hbar n_{\Delta1}/g$.
Use the proved far old-wave Hessian $L_{SS}$ in the source.
The differential inequality has the form $H\prime\le2\alpha\sqrt H+f$,
where $\alpha=F_{\rm old}\sqrt{n_{\Delta2}}$. The function
$[\alpha t+\sqrt{\int_0^t f}]^2$ is a supersolution, because its
derivative is at least $2\alpha\sqrt H+f$ when evaluated at that
function. Bounding the nonnegative source integral gives
\begin{equation}
 H_\Delta=\left[F_{\rm old}T\sqrt{n_{\Delta2}}+
 \sqrt{r_{SS}(L_{SS})T+4F_{\rm old}\hbar n_{\Delta1}T/g
                         +\hbar n_{\Delta2}T+W_\Delta}\right]^2.
 \label{prep:improved-hessian}
\end{equation}
It proves $\|p_S^2\Delta(T)\|\le H_\Delta$ and
$H_\Delta/\hbar^2<1.828510\cdot10^{19}$.
The coarse Hessian only generated the amplitude collar; it was not
relabelled as the improved source or inserted into this nonlinear
old-force term.

\subsection{Handoff enclosures and the preparation event}
\label{prep:handoff-table}
The constructions above give the following outward decimal enclosures.
Every endpoint in the table is a rational number. The unrounded
expressions, not shortened displays, define the quantities when used
in the final event sum.
\begin{center}\small
\begin{tabular}{ll}
\toprule
Quantity&Strict upper bound\\\midrule
$D_o$&$3.00668265254473406785\cdot10^{-11}$\\
$P_{S,o}$&$6.98319079275307899984\cdot10^{-32}$\\
$n_{1T}$&$6.24924697825783043607\cdot10^{-30}$\\
$H_{SS}$&$3.23800849281954087759\cdot10^{-49}$\\
$H_{SS}^{o}$&$3.24499145042629721024\cdot10^{-49}$\\
$n_{2T}$&$5.97036914776425754723\cdot10^{-38}$\\
$P_4$&$1.07556686771573853653\cdot10^{-76}$\\
$L_{SS}$&$2.53407404062459981525\cdot10^{-57}$\\
$D_\Delta(T)$&$2.767839\cdot10^{-30}$\\
$P_{\Delta,S}(T)$&$7.597256\cdot10^{-36}$\\
$P_{\Delta,r}(T)$&$2.08385966319223143790\cdot10^{-37}$\\
$A_h$&$5.04959232864420059458\cdot10^{-24}$\\
$A_{2h}$&$1.22824296520291159974\cdot10^{-17}$\\
$U_h$&$276122.969821611569892$\\
$P_h$&$3.20281666862106903351\cdot10^{-31}$\\
$H_\Delta/\hbar^2$&$1.82850947100693947781\cdot10^{19}$\\
\bottomrule
\end{tabular}
\end{center}

\subsection{Preparation history under the enlarged exact flow}
\label{prep:history}
The preparation rank $K(S,r)$ is independent of $Z$; hence its material
derivative under the enlarged flow has no pointer-current term. Its
explicit drift $\beta$ in \eqref{prep:old-current} is unchanged. The
error in its current must nevertheless be recomputed for that flow.
In the bounded gauge the additional momentum ceilings are
\[
 d_r^g=P_{\Delta,r}(T)+ma_0(0.201)D_\Delta(T),\qquad
 d_S^g=P_{\Delta,S}(T)+ma_1(0.080802)D_\Delta(T)/2.
\]
The exact normalized wave stays in the derivative-error factor of the
bilinear expansion, retaining the quadratic error terms. Consequently
\begin{equation}
 \Delta I=T\left[(\hbar c_r/m+\hbar c_S/M)D_\Delta(T)
                           +(k_r/m)d_r^g+(k_S/M)d_S^g\right]
       <1.304454833\cdot10^{-22}. \label{prep:extra-current}
\end{equation}
Under the single original joint domination constant
$C=270000/67499$, $N$ fixed radial cuts therefore have preparation
failure probability bounded by
\begin{equation}
 E_p(N)=2CN\beta+2C\sqrt{2N(I_o+\Delta I)}
                        +C[D_o+D_\Delta(T)]^2.
 \label{prep:event}
\end{equation}
This follows by taking initial rank collars of width
$\beta+\epsilon$, using Markov only on the unknown absolute current,
and minimizing in $\epsilon$. The last term covers the terminal clock
suffix where the helper vanishes. It is an estimate on the enlarged
flow with its original initial law; closeness of two waves alone has
not been used to assert history agreement.
The initial conditional-CDF interface costs, separately,
\begin{equation}
 E_{\rm cdf}(N)=2CN[D_o+D_\Delta(T)]. \label{prep:cdf}
\end{equation}
For the final cut family, including the exterior ray, use $N=256001$
in both formulas. Evaluating the finite expressions gives
\[
 \begin{aligned}
 E_p(256001)&<0.000003084949592860634224432805,\\
 E_{\rm cdf}(256001)&<0.0000615780135255954031515710.
 \end{aligned}
\]
The pointer remains part of the one jointly dominated auxiliary stock;
no independent pointer cap, sector sample, equilibrium reset, or new law
at handoff enters these estimates.


\section{Radial conditional comparison and the old clock current}
\label{old:activation}\label{old:main}
This appendix supplies the radial comparison estimates used before the
writer is active. Detector coordinates and ages are used until physical
clock units are explicitly restored. Put
\[
 \begin{gathered}
 m=2^{52},\quad h=1/128,\quad L=100,\quad V=mh^2/8=2^{35},\\
 \eta=2^{-17},\quad g=1/300,\quad s_-=355/300,\quad s_+=377/300.
 \end{gathered}
\]
Write $W_i=\Theta(r-1000)V_i(r)$, $F=\chi_q U_f$, $\chi_q(r)=\Theta((r-2000)/10)$,
$U_f=(r^2/L^2-3)/(2mL^2)$, and
\[
 H_i(s)=-(2m)^{-1}\partial_r^2+\gamma(s)W_i+(1-\gamma(s))F,
 \qquad \gamma(s)=B_9(s/g).
\]
The decreasing cutoffs have their specified constant continuations.
The exact conditional solution starts from
$\chi(r)=2r e^{-r^2/(2L^2)}/(\pi^{1/4}L^{3/2})$ on $r>0$.
All estimates hold separately in each conserved sector and therefore for
their coherent spinor, without sampling a sector. Norms are half-line
norms or, equivalently, norms of the normalized odd extension.

\subsection{Gaussian folds and transported cell comparisons}\label{old:lens}
Gaussian integration gives
\[
 d_j^2=\|\chi^{(j)}\|^2=\frac{(2j+1)!!}{2^jL^{2j}},\qquad
 B_j=(d_j^2+2h d_jd_{j+1})^{1/2}.
\]
Indeed $h\sum_k f(y+kh)\leq\int f+h\TV(f)$ for nonnegative
integrable $f$ of bounded variation: compare the value in each cell to its
cell average and sum the cell variations. Apply this to the square of
the $j$th derivative of the odd extension, whose variation is at most
$2d_jd_{j+1}$. Multiplication by an $h$-periodic taper of normalized cell
norm $A_k$ consequently costs at most $B_jA_kh^{-k}$.
For a taper polynomial $P$ rising from zero to one on a unit transition,
with two transitions of width $\eta$ and two omitted edge strips of width
$\eta$, define
\begin{align*}
 e_0^2&=2\eta+2\eta\int_0^1(1-P)^2, & A_0&=1,\\
 A_k^2&=2\eta^{1-2k}\int_0^1(P^{(k)})^2,&
 C_k&=\sum_{j=0}^k\binom kj B_j A_{k-j}h^{-(k-j)},\\
 D_0&=B_0e_0,& c&=C_2/(2m).
\end{align*}
All taper integrals here are polynomial integrals with rational
coefficients. Two different analytical tapers will be used for the same
exact conditional wave. Their estimates are not interchanged.

The lens $b''+\gamma b=0$, $b(0)=1$, $b'(0)=0$ satisfies, for
$0\leq s<\pi/2$,
\[
 \cos s\leq b\leq1,\quad b'\leq0,\quad |b'|\leq\sin s,
 \quad b'^2+\gamma b^2\leq1.
\]
For the first inequality, $b-\cos s$ is the sine convolution of
$(1-\gamma)b$ until a hypothetical first zero; it prevents that zero.
Also $(b'^2+b^2)'=2b'b(1-\gamma)\leq0$, and differentiating $b-\cos s$
or integrating the equation yields the derivative bound.
In each cell centered at $c_i$ set $x=r-c_i$, $y=x/b$, and
\[
 \Phi_i(s,r)=b^{-1/2}e^{im b'x^2/(2b)}
             \chi(c_i+y)a(y/h).
\]
The core residual has norm $\leq cb^{-2}$. Its radial derivative,
divided by $mh$, has norm at most
$ c|b'|/(2b^2)+C_3/(2m^2hb^3)$, by differentiating its amplitude and
quadratic phase. At an origin well the helper is odd; at an origin
barrier it vanishes in a collar. Thus the even extension of each potential
and the odd extension of the helper justify integration by parts without
an origin force or a self-adjoint half-line momentum assertion.

Here and below use the explicit weak-potential bounds
\begin{align*}
 F_0&=\frac{2010^2}{2mL^4}+\frac{3}{2mL^2},\\
 F_1&=F_0/2+2010/(mL^4),\\
 F_2&=1000F_0+2010/(mL^4)+1/(mL^4).
\end{align*}
For the outer mismatch the harmonic backflow of $r\geq1000$ is
larger than $990$. The recurrence
\[
 I_k(z)=\tfrac12z^{k-1}e^{-z^2}+\tfrac{k-1}{2}I_{k-2}(z),
 \qquad I_0(z)\leq e^{-z^2}/(2z)
\]
for Gaussian tails, together with $e^2>7$, gives
\[
 T_0^2=21/7^{49},\quad T_1^2=2040/(L^2 7^{49}),\quad
 T_2^2=205100/(L^4 7^{49})
\]
as strict upper bounds for the first three Gaussian tail norms.
Put $\epsilon=VT_0+F_0$.

For the quintic taper $P=10z^3-15z^4+6z^5$, direct integration gives
\[
 e_0^2=2\eta+181\eta/231,\quad A_1^2=20/(7\eta),\quad
 A_2^2=240/(7\eta^3),\quad A_3^2=1440/\eta^5.
\]
Define $M_1=15/(8\eta h)$,
$f=1/2+5h/8+F_1/(mh)$,
$q_0=(B_1e_0+B_0A_1/h)/(mh)$, and
$J_3=(\sec s\tan s+\log(\sec s+\tan s))/2$.
Duhamel's inequality and its differentiated equation give the explicit
increasing majorants
\begin{align}
 D(s)&=D_0+c\tan s+\epsilon s,\label{old:D}\\
 q(s)&=q_0+f\{D_0s+c\log\sec s+\epsilon s^2/2\}
       +c(\sec s-1)/2+C_3J_3/(2m^2h)\notag\\
 &\quad +(1/2+5h/8)T_0s+VT_0(1-\cos s)/2
       +\frac{V(T_1+M_1T_0)}{mh}\log(\sec s+\tan s)\notag\\
 &\quad +\frac{F_1s}{mh}
       +F_0\left\{(1-\cos s)/2+\frac{C_1}{mh}
                                      \log(\sec s+\tan s)\right\}.
 \label{old:q}
\end{align}
They bound $\|u_i-\Phi_i\|$ and
$\|\partial_r(u_i-\Phi_i)\|/(mh)$. In detail, the force term in the
differentiated error equation is $\leq f\|u_i-\Phi_i\|$;
the outer-mismatch derivative contributes
$(1/2+5h/8)T_0+V|b'|T_0/2+V(T_1+M_1T_0)/(mhb)$;
and the weak-trap derivative contributes
$F_1/(mh)+F_0\{|b'|/2+C_1/(mhb)\}$.
Integrating these terms and the core residual proves
\eqref{old:D}--\eqref{old:q}, rather than assuming an abrupt switch.

\subsection{The fixed-age sorting bound and the complete age union}
\label{old:endpoint}
Let $G_1(s)=\int_0^s\gamma(x)\dd x$ and
$G_2(s)=\int_0^s(s-x)\gamma(x)\dd x$ during the gate.
For $s\leq g$, $b\leq1-\cos(g)G_2(s)$ and $|b'|\leq G_1(s)$.
The identities $G_1(g)=g/2$, $G_2(g)=3g^2/22$ imply, for $s\geq g$,
\begin{align*}
 b_U(s)&=(1-3g^2\cos g/22)\cos(s-g)
                          -(g\cos g/2)\sin(s-g),\\
 p_U(s)&=\sin(s-g)+(g/2)\cos(s-g).
\end{align*}
Use the preceding bounds as $b_U,p_U$ before $g$ and put
$G_-(s)=1/2-(1/2-\eta)b_U(s)$,
$G_+'(s)=(1/2-\eta)p_U(s)$.
These functions are nonnegative and increasing on the needed interval;
$G_-$ is a lower gap and $G_+'$ an upper gap speed, not derivatives
of a common artificial lens.

For completeness the moving-gap estimate is obtained with linear CDF
cutoffs centered on each positive barrier, with half-width $w(s)$.
The helper vanishes in these gaps, so the mass and integrated absolute
current there are at most $D(s)^2$ and $hD(s)q(s)$ after disjoint gaps
are summed. Differentiating the expectation of a moving linear cutoff
costs at most $w'D^2/(2w)+hDq/(2w)$. Initial symmetric smoothing costs
$\epsilon_0$ below, and replacing a soft terminal cut by each of the two
associated sharp classifier cuts costs their intervening error mass.
Summing over both cuts at every positive barrier gives
\begin{equation}
 S\leq2\epsilon_0+2D(s_-)^2+
     \int_0^{s_-}\frac{G_+'(s)D(s)^2+D(s)q(s)}{G_-(s)}\dd s.
 \label{old:sorting}
\end{equation}
There is no collar across zero: its cumulative mass is exactly zero, and
the single first-cell classifier is paid by its share of $D^2$.
At $s_-$ the helper gap contains the classifiers because $G_-(s_-)>1/4$.
The initial Gaussian density $\rho_0=\chi^2$ obeys
\[
 \TV(\rho_0)<12/(7L),\qquad \TV(|\rho_0'|)\leq20/L^2.
\]
The first follows from its single maximum
$4/(e\sqrt\pi L)<6/(7L)$; the second from integrating the
Gaussian derivative polynomial. Symmetric linear smoothing and linear
interpolation in half-cells of length $\ell=h/2$ give
\begin{align*}
 \epsilon_0&=\frac{(h\eta)^2}{6}
       \left\{\frac{12}{7Lh}+\frac{20}{L^2}\right\},\\
 J&=\frac\ell8\frac{12}{7L}+\frac{\ell^2}8\frac{20}{L^2},\qquad
 R=0.000052316289\left(1+\ell\frac{12}{7L}\right).
\end{align*}
For the interpolation estimate, the integrated error of a linear
interpolant on a cell is at most one eighth of the cell length times
the variation of the function there; apply this first to $\rho_0$ and
then to its derivative. The reference preimage in each half-cell is
shifted from the ideal sector cut by at most
$0.000052316289\ell$, as follows from the periodic rotor error
$(0.007233)^2$. Summing the corresponding initial masses with the
Gaussian fold inequality gives $R$.
Thus $R$ refers to the original reference age $1.22$, not $s_-$.

Here is also an explicit finite-range bound used in this comparison.
For the periodic conditional solution $z_i$ starting flat on a circle
of length two, energy differentiation gives
$\|z_i'\|\leq mh/2$, since
$\frac d{ds}\langle T+\gamma V_i\rangle=\gamma'\langle V_i\rangle
\leq\gamma'V$. Each periodic cell has mass $h/2$.
The auxiliary wave $\sqrt2\chi z_i$ has residual norm at most
\begin{equation}
 \beta=B_2/(2m)+(h/2)B_1+VT_G+F_0B_0,
 \qquad T_G=(22/7^{49})^{1/2}.
 \label{old:outerbeta}
\end{equation}
The two kinetic products use the Gaussian folds; the last two terms
pay the finite outer and weak cutoffs. Its outer norm is at most $T_G$,
because $h\sum_{k\geq0}\rho_0(1000+kh)\leq
\int_{1000}^\infty\rho_0+h\rho_0(1000)<22/7^{49}$.
Consequently the fixed-age mismatch with the original periodic label is
\begin{equation}
 0.000002144+C\{S+J+R+(s_-\beta+T_G)^2+201/(10\,7^{50})\},
 \qquad C=270000/67499.
 \label{old:endpointexpression}
\end{equation}
This is less than $0.002588532236512050$. The quiet-prefix allowance
is included once.

For the whole age interval use the fixed transition bands between
relative distances $h/5$ and $3h/10$ from neighboring wells.
The quintic helper and its derivative vanish on these bands because
$b(s)\leq\cos(s-g)$ and $\cos(s_--g)<0.4$.
A Lipschitz soft binary label with slope at most $10/h$ must vary by
at least $1/2$ along any trajectory that starts outside the bands and
subsequently changes its label. Hence, directly under the reference
wave law,
\begin{equation}
 q_{\rm age}=D(s_-)^2+20\int_{s_-}^{s_+}D(s)q(s)\dd s,
 \qquad Cq_{\rm age}<0.000263432327209778.
 \label{old:ageunion}
\end{equation}
This is a pathwise union estimate and permits the final age to depend on
all other coordinates. One does not replace an age union by its largest
single-age probability.

For an outer threshold the situation is different: reference ranks that
escape at age $s$ form the upper ray $[F_s(1000),1]$. Their union is one
ray. Therefore, on $[-1/8,13/10]$, its mass is at most
$q_{\rm out}=((13/10)\beta+T_G)^2$, and
\begin{equation}
 Cq_{\rm out}<1.547666958047712\times10^{-8}.\label{old:outer}
\end{equation}
The negative-age extension is the exact finite-trap evolution treated
below and has a still smaller tail. If the actual rank differs from the
original Gaussian rank by at most $\delta$, expanding this ray by
$\delta$ also covers the actual terminal escape. Adding it to the binary
moving cuts uses $N=256001$ cuts in the conservative $2N\delta$ allowance.
With known drift $\beta_*$ and remaining wave-law mean absolute drift
$I$, Markov's inequality and optimization in $\delta-\beta_*$ give
\begin{equation}
 C(q_{\rm age}+q_{\rm out})+2CN\beta_*+2C\sqrt{2NI}
       +\Prob(\hbox{clock-age failure}).\label{old:rankbridge}
\end{equation}
Only domination of the original law is used. The same expanded outer
ray pays whole-history escape if clock confinement and the absolute
rank integral hold throughout the whole interval.

\subsection{Energy graphs and the reference clock functional}
\label{old:graphs}
The following coefficients also supply the energy inputs for holding.
Put $A=W-F$, $A_0=V+1$, and use
\[
 A_1=2^{46},\quad A_2=2^{80},\quad A_3=2^{112},\quad A_4=2^{160}
\]
for derivative ceilings of $A$, with $\|U'\|\leq A_1$.
They are deliberately loose bounds, not additional hypotheses. For the
cap its first derivative is $-S_5(z)(1-2\eta z)$; the sum of absolute
coefficients is at most $32$, and the next three derivatives are bounded
by $32\,6^{k-1}\eta^{-(k-1)}$ with the cell powers of $h$ restored.
Leibniz's rule gives the displayed bounds after the outer cutoff.
The smooth logistic cutoff has derivatives through order four bounded
by $10^{10}$: on its central half the logit derivatives are bounded by
$32,256,3072,49152$; at an edge put $x=1/t\geq4$, use sigmoid derivatives
$\leq75e^{-A}$, $e^{-A}\leq4e^{-x}$, and the Bell-polynomial bound
$50x^8$. Its maximum is at most $15000\,8^8/7^4<10^{10}$.
The weak cutoff rescales these derivatives by $10^{-k}$.
The capped potential is $C^3$ with bounded weak fourth derivative,
which suffices for these third energy graphs; no fourth radial
Hamiltonian graph is invoked.

Define, entirely algebraically,
\begin{align*}
 N_1&=A_0, & P&=\{2m(N_1+1)\}^{1/2},\\
 B_*&=A_2/(2m)+A_1P/m,& N_2&=(1001/1000)A_0^2,\\
 P_H&=\{2m(N_1N_2+N_1^2)\}^{1/2},&
 U_{rr}&=2m(N_1+A_0),\qquad N_3=(1001/1000)A_0^3.
\end{align*}
Then $\|H^j u\|\leq N_j$, $j=1,2,3$, throughout the gate and static
plateau. To check this directly, $\|H_0^2\chi\|,\|H_0^3\chi\|<1$
by the Gaussian products and the finite weak cutoff. The first graph obeys
$\|H u\|\leq\|H_0\chi\|+V+F_0<V+1=N_1$ by graph Duhamel
and $\int\gamma'=1$. The identities
\[
 [H,A]=-(A''+2A'\partial_r)/(2m),\quad
 [H,[H,A]]=\frac{A''''+4A'''\partial_r+4A''\partial_r^2}{4m^2}
                                      +A'U'/m
\]
and $\int\gamma'=1$ give
\begin{align*}
 \|H^2u\|&\leq1+2A_0+A_0^2+B_*<N_2,\\
 \|H^3u\|&\leq1+3A_0(1+A_0+A_0^2/3+B_*/2)\\
 &\quad+3\{A_2N_1/(2m)+A_1P_H/m\}
       +(A_4+4A_3P+4A_2U_{rr})/(4m^2)+A_1^2/m<N_3.
\end{align*}
For the second line retain the intermediate bound
$\|H^2u\|\leq(1+A_0\gamma)^2+B_*\gamma$ and integrate
$(H^3)'=\gamma'(3AH^2+3[H,A]H+[H,[H,A]])$.
Smooth approximation on the common graph domain justifies the identities.

To evaluate the smaller reference-current bound use the transported
cell construction with $P=B_9$ and $C_k$ through order four. Fix a
horizon $a<\pi/2$ (here $a=s_+$; the holding calculation uses $4/3$).
With $M_1=(315/128)/(\eta h)$, $M_2=(2520/64)/(\eta^2h^2)$ put
\begin{align*}
 T_{A1}&=T_1+M_1T_0,\\
 T_{A2}&=T_2+2M_1T_1+M_2T_0,\\
 P_t&=mhT_0/2+\sec(a)T_{A1},\\
 P_{t2}&=m^2h^2T_0/4+m\tan(a)T_0
                 +mh\sec(a)T_{A1}+\sec(a)^2T_{A2},\\
 W_1&=mh/2+5V,\\
 W_2&=m(1+15/(16\eta))+5mh+10^5V,\\
 B_c&=(W_2T_0+2W_1P_t+VP_{t2})/(2m)+(V+F_0)VT_0,\\
 P_\Phi&=mh/2+\sec(a)C_1,\\
 H_\Phi&=V+\tan(a)/2+\tan(a)hC_1/2+c\sec(a)^2+VT_0+F_0,\\
 B_f&=F_0H_\Phi+F_2/(2m)+F_1P_\Phi/m,\\
 B_m&=\{VT_{A2}/(2m)+F_0c\}\sec(a)^2.
\end{align*}
Acting with the harmonic core on its residual gives
\[
 \|H R_{\rm core}\|\leq Vcb^{-2}
   +(-b'/b)(hC_3+C_2)b^{-2}/(4m)+C_4b^{-4}/(4m^2).
\]
The preceding $B_c,B_f,B_m$ bound respectively the differentiated outer
mismatch, weak mismatch, and their action on the core residual. Thus
\begin{align}
 D_9(s)&=D_0+c\tan s+\epsilon s,\notag\\
 Q_9(s)&=c+F_0D_0+(V+F_0)D_9(g)+Vc\tan s
       +\frac{hC_3+C_2}{8m}(\sec^2s-1)\notag\\
 &\quad+\frac{C_4}{4m^2}(\tan s+\tan^3s/3)
                         +(B_c+B_f+B_m)s.\label{old:DQ9}
\end{align}
These bound $\|u-\Phi\|$ and $\|H(u-\Phi)\|$. The gate derivative
source is at most $(V+F_0)D_9(g)$; it is not omitted.
Set $P_E=\{2mD_9(Q_9+F_0D_9)\}^{1/2}$,
$P_{E2}=2m\{Q_9+(V+F_0)D_9\}$ at $s=a$.
The form inequality and half-line Sobolev give
$\|E\|_\infty\leq(2D_9P_E)^{1/2}$,
$\|E'\|_\infty\leq(2P_EP_{E2})^{1/2}$.
With $\rho_*=6/(7L)$ use
\begin{align*}
 P_\infty&=(\rho_*\sec a)^{1/2},\\
 P_\infty'&=mhP_\infty/2+
   \sec(a)^{3/2}\{(2d_1d_2)^{1/2}+M_1\sqrt{\rho_*}\},\\
 j_*&=(h/2)\tan(a)\rho_*+
 \{P_\infty\sqrt{2P_EP_{E2}}+
   P_\infty'\sqrt{2D_9P_E}+
   \sqrt{2D_9P_E}\sqrt{2P_EP_{E2}}\}/m.
\end{align*}
This follows by expanding all three error terms in the bilinear current,
including $E^*E'$. For $a=s_+$, $D_9<0.006877$, $Q_9<2.364\times10^8$,
and $j_*<9.194\times10^6$.

Restore $M=10$, $v=1$, $T_0=0.03$, $\sigma=0.001$ in SI units and
$\hbar=6.62607015\times10^{-34}/(2\pi)$.
For the next two current identities, use physical radius $r_{\rm phys}=10^{-4}r$
and its normalized conditional wave, of density $\rho=\sum_i|u_i|^2$;
sector sums are implicit in products and $r$ in their integrals denotes
$r_{\rm phys}$. The radial physical mass is
$m_{\rm phys}=2^{52}\hbar T_0/(10^{-4})^2$.
For $G=\phi(S-vt-S_c)u(S,r_{\rm phys})$ the radial energy density and clock correction
are $e=\Re(u^*H_{\rm phys}u)$ and $k=-e/(Mv)$.
If $K$ is the conditional CDF, write $B_k=\int_0^r k$, $b_k=\int k$.
Substituting into the exact rank numerator gives
\[
 \mathcal F(G)=-\phi^2kj/v-
 \rho\{2\phi\phi'(B_k-Kb_k)+\phi^2(\partial_SB_k-K\partial_Sb_k)\}.
\]
Energy continuity gives
\[
 \partial_SB_k-K\partial_Sb_k=\frac{j_E}{Mv^2}
 -\frac{\dot\gamma}{Mv^2}
       \left\{\int_0^r(W_{\rm phys}-F_{\rm phys})\rho
                        -K\langle W_{\rm phys}-F_{\rm phys}\rangle\right\},
\]
where $j_E=(\hbar/(2m_{\rm phys}))\Im(u^*\partial_{r_{\rm phys}}(H_{\rm phys}u)
 +(H_{\rm phys}u)^*\partial_{r_{\rm phys}}u)$.
The intrinsic term is bounded by
$\hbar a j_*N_1/(Mv^2T_0)$.
Returning to detector units, for the energy term use the $H^2$ current
$J_2=m^{-1}\Im(u^*(Hu)'+(Hu)^*u')$,
$\|J_2\|_1\leq(P_H+N_1P)/m$.
For the conservative duration $T_s=0.0377$ let
$\varepsilon=\hbar T_s/(2Mv^2T_0^2)$ and bound its density by
\[
 \rho_e=\left[P_\infty+
 \sqrt{2\widetilde D\sqrt{2m\widetilde D
                 (\widetilde Q+5\times10^{-18}\widetilde D)}}\right]^2,
 \quad \widetilde D=D_9+\varepsilon N_2,\quad
 \widetilde Q=Q_9+\varepsilon N_3.
\]
This enlarged density also bounds the unchanged reference.
The energy contribution is at most
$\varepsilon\rho_e(P_H+N_1P)/m$.
Finally $|B_k-Kb_k|\leq2\|H_{\rm phys}u\|/(Mv)$,
$\int2|\phi\phi'|\leq2p_1$, and $\int\dot\gamma\dd t=1$
for each clock offset. With $T=0.04$, the reference functional is bounded by
\begin{equation}
 I_G\leq\frac{\hbar a j_*N_1}{Mv^2T_0}
 +\frac{\varepsilon\rho_e(P_H+N_1P)}m
 +\frac{4p_1T\hbar N_1}{MvT_0}
 +\frac{2\hbar A_0}{Mv^2T_0}+I_q
 <2.382267143194879\times10^{-16}.
 \label{old:referencecurrent}
\end{equation}
The quiet correction $I_q$ is explicitly controlled below.
This is a functional of the comparison wave, not an assertion that
$G$ solves the autonomous Schr\"odinger equation.

\subsection{The finite-trap quiet boundary}\label{old:quiet}
The Gaussian is not stationary for the finite trap. Put
$\alpha=(2m)^{-1}$, $w=(\chi_q-1)U_f$, $u=w\chi$.
Then
\begin{align*}
 H_F\chi&=u,&H_F^2\chi&=-\alpha u''+Fu,\\
 H_F^3\chi&=\alpha^2u''''-\alpha F''u-2\alpha F'u'
                               -2\alpha Fu''+F^2u.
\end{align*}
All terms are supported at $r\geq2000$. Here is an explicit polynomial
recipe for their norms. Set $y=r/L$, $q_f=1/(2mL^2)$,
$U_0(y)=y^2+3$, $U_1(y)=2y/L$, $U_2(y)=2/L^2$,
and $(c_0,c_1,c_2,c_3,c_4)=(1,1/2,32/25,1000,20000)$.
These bounds follow by rescaling the cutoff derivatives
$5,128,10^6,2\times10^8$ by its width ten.
Define
\begin{align*}
 W_j(y)&=\sum_{k=0}^{\min(2,j)}\binom jk c_{j-k}U_k(y),\\
 (R_0,R_1,R_2,R_3,R_4)&=\left(1,\frac{1+y}L,
 \frac{3+y^2}{L^2},\frac{3+6y+y^3}{L^3},
 \frac{15+10y^2+y^4}{L^4}\right),\\
 Z_j(y)&=\sum_{k=0}^j\binom jk W_{j-k}(y)R_k(y),\qquad
 U_j^*=q_f Z_j(1)10^{12}/7^{100}.
\end{align*}
Each $Z_j$ has nonnegative coefficients and degree at most six.
The Gaussian recurrence gives
$\|\one_{y\geq20}y^k\chi\|\leq10^{12}/7^{100}$ for $k\leq6$,
so $\|u^{(j)}\|\leq U_j^*$. Use
\begin{align*}
 Q_1&=2600/(2mL^2 7^{100}),\\
 Q_2&=\alpha U_2^*+F_0U_0^*,\\
 Q_3&=\alpha^2U_4^*+\alpha F_2U_0^*
                +2\alpha F_1U_1^*+2\alpha F_0U_2^*+F_0^2U_0^*.
\end{align*}
The sharper first bound is the direct recurrence for
$(y^2+3)\chi$ on $y\geq20$.
For $u(s)=e^{-isH_F}\chi$, $-1/8\leq s\leq0$, the boundary
errors after multiplication by the real clock packet are
\begin{equation}
 \Delta_k=p_kQ_1/8+
 \sum_{j=1}^k\binom kj p_{k-j}Q_j/(vT_0)^j,
 \qquad 0\leq k\leq3.\label{old:quietjets}
\end{equation}
In particular
$(\Delta_0,\Delta_1,\Delta_2,\Delta_3)<
(1.116\times10^{-102},2.947\times10^{-99},
1.195\times10^{-95},6.072\times10^{-92})$ in the corresponding
clock derivative units. The constants $p_k$ are given next.
Substituting these in the rank functional pays the negligible $I_q$;
$10^{-40}$ is a convenient outward ceiling for that term in
\eqref{old:referencecurrent}.

\subsection{Clock coefficients and the five-error decomposition}
\label{old:clockcoeff}
The smoothed cosine-fourth packet has the following derivative ceilings,
with $q_s=1.000001$ and $p_0=1$:
\[
 (p_1,p_2,p_3,p_4)=q_s\left(
 \sqrt{16/7}\frac\pi{2\sigma},\quad
 \sqrt{512/35}\left(\frac\pi{2\sigma}\right)^2,\quad
 \sqrt{1024/7}\left(\frac\pi{2\sigma}\right)^3,\quad
 \sqrt{69632/35}\left(\frac\pi{2\sigma}\right)^4\right).
\]
They follow by expanding the powers of sine and cosine and integrating
on the packet support. Convolution with the nonnegative mollifier
contracts every derivative norm; normalization costs at most $q_s$.
Let $d=0.0001$ seconds and
\[
 b_1=(5/2)/d,\quad b_2=(2520/64)/d^2,\quad
 b_3=(7560/16+5040/64)/d^3,\quad b_4=1360800/d^4.
\]
These are bounds for the first four physical-time derivatives of $B_9$;
the fourth bound is the absolute coefficient sum. Define
\begin{align*}
 c_1&=N_1/(vT_0),\\
 c_2&=N_2/(v^2T_0^2)+b_1A_0/(v^2T_0),\\
 c_3&=N_3/(v^3T_0^3)+b_1(3A_0N_1+B_*)/(v^3T_0^2)
                                      +b_2A_0/(v^3T_0),\\
 G_1&=(p_1^2+c_1^2)^{1/2},\qquad
 G_2=p_2+2p_1c_1+c_2,\\
 G_3&=p_3+3p_2c_1+3p_1c_2+c_3,\qquad
 Q_G=q_s(4\pi^2/\sigma^2)\sqrt{3/35}+c_1^2.
\end{align*}
Differentiating the conditional equation gives these bounds for
$\|G_S\|,\|G_{SS}\|,\|G_{SSS}\|$ and its logarithmic-envelope
coefficient. In the third derivative use
$2H_SH+HH_S=3H_SH+[H,H_S]$; the commutator costs $B_*$.
The first derivative improves to a sum of squares because the real
packet derivative and the unitary conditional derivative have zero real
cross term.

The denominator-free rank stability estimate of
Section~\ref{sec:rankproof}, for an error with norms $D,e_1,e_2$, is
\begin{equation}
 \mathcal R[D,e_1,e_2]=\frac\hbar M\int
 (4G_1e_1+2e_1^2+14DQ_G+9DG_2+e_2)\dd t.
 \label{old:Rfunctional}
\end{equation}
The derivatives here are clock derivatives. Decompose the exact error
from the same prepared stock as follows. If $E_h$ is the prepared error,
$U_a$ the actual propagator, $\psi_c$ the evolution of the product handoff
under the clipped gate, and $\psi_p$ agrees with $\psi_c$ until
$T_g=0.0033$ and evolves statically thereafter, put
\[
 E_p=U_aE_h,\quad E_c=U_a(\phi\chi)-\psi_c,\quad
 E_l=\psi_c-\psi_p.
\]
Before $T_g$ put $E_r=\psi_c-G$, $E_s=0$; afterward put
\[
 E_r=U_W(t-T_g)(\psi_c(T_g)-G(T_g)),\quad
 E_s=U_W(t-T_g)G(T_g)-G(t).
\]
Then $\Psi-G=E_p+E_c+E_l+E_r+E_s$ identically.
The carried ramp error is counted once, the new static error starts
from zero, and the quiet boundary difference is included in $E_r$.
The free clock and its Galilean transport belong to every propagator.

\subsection{Propagation of the prepared error under the common operator}
\label{old:homogeneous}
Define the old-wave preparation inputs by
\[
 \begin{aligned}
 D_{\rm old}&:=D_o,& P_{10}&:=P_{S,o},& P_{20}&:=H_{SS}^{o},\\
 n_{10}&:=n_{1T},& n_{20}&:=n_{2T}.
 \end{aligned}
\]
Use the defining expressions in \eqref{prep:anisotropic},
\eqref{prep:old-handoff}, and \eqref{prep:second-collar}, not the rounded
enclosures in Section~\ref{prep:handoff-table}. In particular,
$H_{SS}^{o}$ is the ordinary, phase-restored Hessian bound, not the
gauged bound $H_{SS}$. The old-wave norm $D_{\rm old}$ is distinct from
the enlarged-wave norm $D_h$ in Appendix~\ref{prep:main}.
The remote force coefficients are explicitly
\begin{align*}
 F_c&=\mu a_2F_*/2+\mu a_0a_1f_*^2
                     +\mu^2a_1a_2F_*^2/(4M)+10^{-20},\\
 F_{c2}&=\mu a_3F_*/2+\mu(a_1^2+a_0a_2)f_*^2
             +\mu^2(a_2^2+a_1a_3)F_*^2/(4M)+10^{-18},
\end{align*}
where
$(a_0,a_1,a_2,a_3)=(140,43600,70141750,1387431060000)$,
$F_*=0.080802$, $f_*=0.201$, and
$\mu=2^{52}\hbar T_0/(10^{-4})^2$.
These are the product-rule bounds for the retained preparation phase
and its scalar selfpotential; $F_c<5\times10^{-6}$ and $F_{c2}<0.09$.
Set $A_{\rm phys}=\hbar A_0/T_0$, $F_a=b_1A_{\rm phys}$,
$F_{a2}=b_2A_{\rm phys}$ and collar widths $g_1=0.0005$, $g_2=0.00025$.
Initialize
\[
 L_{10}=\sqrt{n_{10}P_{20}}+2\hbar D_{\rm old}/g_1,\qquad
 L_{20}=\sqrt{n_{20}P_{20}}+2\hbar n_{10}/g_2.
\]
The six nonnegative variables $(P_1,P_2,n_1,L_1,n_2,L_2)$ are bounded
by the solution, starting at these values, of
\begin{align*}
 \dot P_1&=F_aD_{\rm old}+F_cn_2,&
 \dot P_2&=2F_aP_1+\hbar F_{a2}D_{\rm old}+2F_cL_2+\hbar F_{c2}n_2,\\
 \dot n_1&=P_1/(Mg_1),&
 \dot L_1&=P_2/(Mg_1)+(F_c+F_a)n_1,\\
 \dot n_2&=L_1/(Mg_2),&
 \dot L_2&=P_2/(Mg_2)+(F_c+F_a)n_2.
\end{align*}
Here the collars move at speed $v$, cancelling Galilean transport;
the inner transition lies where the outer collar equals one.
Clock commutators give the first two equations. Weighted continuity for
the error and for its first clock momentum gives the other four.
The radial kinetic term commutes with every collar. The two initial
local-momentum inequalities follow by integration by parts and
Cauchy--Schwarz. Thus this is a closed positive comparison system,
including the remote preparation force.
With $e_j=P_j/\hbar^j$, use its exact time integrals in
\eqref{old:Rfunctional}; bound $\int P_1^2\leq P_1(T)\int P_1$.
For $T=0.04$ the resulting contribution is
$<1.954625\times10^{-17}$.

\subsection{Short ramp and static correction}
\label{old:rampstatic}
Put $a_j=b_jA_0/(v^jT_0)$, $R_0=\hbar G_2/(2M)$ and
$R_1=\hbar G_3/(2M)$. On $0\leq t\leq T_g$, commutation with the
bounded clock derivatives of the gate gives
\begin{align*}
 D_r(t)&=\Delta_0+R_0t,\\
 e_{1r}(t)&=\Delta_1+R_1t+a_1\Delta_0t+a_1R_0t^2/2,\\
 e_{3r}(t)&=G_3+p_3+3a_1p_2t+3a_1^2p_1t^2+a_1^3t^3
                           +3a_2p_1t+3a_1a_2t^2+a_3t,\\
 e_{2r}(t)&=\sqrt{e_{1r}(t)e_{3r}(t)}.
\end{align*}
The last line is Fourier interpolation
$\|E_{SS}\|^2\leq\|E_S\|\|E_{SSS}\|$.
These are clock derivatives and require only the third radial energy
graph. After $T_g$ the ramp error's clock norms are conserved by the
static propagator. Endpoint ceilings for duration $T=0.04$ in
\eqref{old:Rfunctional} give $<4.073932\times10^{-16}$.
The time $T_g$ is counted from handoff: the trailing support point
$0.1009999999+T_g$ exceeds the gate end $0.1041$.

For the static component the energy gauge
$U(S)=\exp[-iH_{\rm phys}(S-S_0)/(\hbar v)]$ transforms the generator to
$vp+(p-H_{\rm phys}/v)^2/(2M)$.
Its three commuting corrections are generated by
$p^2/(2M)$, $-pH_{\rm phys}/(Mv)$ and $H_{\rm phys}^2/(2Mv^2)$.
The first two ordinary clock derivatives become powers of
$(p-H_{\rm phys}/v)/\hbar$, which commute with all three corrections.
For $T_s=0.0377$, $\varepsilon=\hbar T_s/(2Mv^2T_0^2)$, the triples
$(D,e_1,e_2)$ for these three factors are respectively
\begin{align*}
 D_f&=\hbar T_sp_2/(2M),\\
 e_{1f}&=\frac{\hbar T_s}{2M}(p_3+p_2N_1/(vT_0)),\\
 e_{2f}&=\frac{\hbar T_s}{2M}
          (p_4+2p_3N_1/(vT_0)+p_2N_2/(v^2T_0^2)),\\[2pt]
 D_x&=\hbar T_sp_1N_1/(MvT_0),\\
 e_{1x}&=\frac{\hbar T_s}{Mv}
                    (p_2N_1/T_0+p_1N_2/(vT_0^2)),\\
 e_{2x}&=\frac{\hbar T_s}{Mv}
       (p_3N_1/T_0+2p_2N_2/(vT_0^2)+p_1N_3/(v^2T_0^3)),\\[2pt]
 D_E&=\varepsilon N_2,\qquad
 e_{1E}=p_1\varepsilon N_2+\varepsilon N_3/(vT_0),\\
 e_{2E}&=p_2\varepsilon N_2+2p_1\varepsilon N_3/(vT_0)
                                   +\sqrt{2\varepsilon}N_3/(v^2T_0^2).
\end{align*}
The last term follows from
$|e^{-i\varepsilon\lambda^2}-1|\leq\sqrt{2\varepsilon}|\lambda|$.
Thus $\|H^2(e^{-i\varepsilon H^2}-1)u\|
\leq\sqrt{2\varepsilon}\|H^3u\|$; no fourth radial graph is required.
Sum the triples, and use the static versions of $G_2$
(with its gate term removed) and $G_1,Q_G$ in
\eqref{old:Rfunctional}. The static contribution is
$<3.770360\times10^{-19}$.

\subsection{Both remote clock comparisons}
\label{old:tails}
For the clipped evolution let $\mathcal A_j=A_{\rm phys}b_j/v^j$,
$j=1,\ldots,4$, and $\mathcal A_0=A_{\rm phys}$.
The global physical clock moments obey
\[
 \dot P_k\leq\sum_{j=1}^k\binom kj\hbar^{j-1}\mathcal A_jP_{k-j},
 \qquad P_k(0)=\hbar^kp_k,\quad P_0=1.
\]
Define their polynomial majorants recursively by equality. This is an
explicit finite triangular recursion to degree $k$, using no radial
derivative. Four static left collars of width $g_c=0.0002$ are supported
below $0.10081$, below the initial clock support and before activation.
Their transitions are nested; where they act all clipped clock forces
vanish. Weighted continuity therefore yields the local polynomials
\[
 l_k(t)=(Mg_c)^{-(4-k)}\mathcal I^{4-k}P_4(t),
 \qquad k=0,1,2,\qquad \mathcal I f(t)=\int_0^t f(s)\dd s.
\]
They bound $\|\chi_0p^k\psi_c\|$. Set
\[
 V_c=\mu a_1F_*/2+\mu a_0^2f_*^2/2
                +\mu^2a_1^2F_*^2/(8M)+10^{-20},
 \quad C_1=F_c+\mathcal A_1,\quad C_2=F_{c2}+\mathcal A_2.
\]
The common-minus-clipped source is supported in this collar.
Its norm and momentum rates are
\[
 f_0=V_cl_0/\hbar,\quad f_1=V_cl_1/\hbar+F_cl_0,
 \quad f_2=V_cl_2/\hbar+2F_cl_1+\hbar F_{c2}l_0.
\]
Starting from zero, its bounds are the explicit polynomials
\[
 D_c=\mathcal I f_0,\quad P_{1c}=\mathcal I(f_1+C_1D_c),\quad
 P_{2c}=\mathcal I(f_2+2C_1P_{1c}+\hbar C_2D_c).
\]
This retains the common propagator's force; differentiating a Duhamel
integral as though the propagator commuted with $p$ would not suffice.
Their exact polynomial time integrals in \eqref{old:Rfunctional}, with
$\int P_{1c}^2\leq P_{1c}(T)\int P_{1c}$, give
$<2.240929\times10^{-21}$.

For the clipped-to-static comparison use the moving weight
$\chi(t,S)=\min(1,\exp[-10^6(S-0.1009-vt)])$.
Its derivative satisfies $\chi_t+v\chi_S=0$ and
$|\chi_S|\leq10^6\chi$. Let $L_k=\|\chi p^k\psi_c\|$,
$0\leq k\leq3$, and bound them with the positive system
\[
 \dot L_k=(10^6/M)L_{k+1}
   +\sum_{j=1}^k\binom kj\hbar^{j-1}\mathcal A_jL_{k-j},
 \quad L_4=P_4,\quad L_k(0)=(2/7^{50})\hbar^kp_k.
\]
The initial bound follows from $e^{-99.9999}<2/7^{50}$.
For $t\geq T_g$, the weight equals one on the entire remaining gate,
so the late error obeys
\begin{align*}
 D_l&\leq(\mathcal A_0/\hbar)\int_0^T L_0,\\
 P_{1l}&\leq(\mathcal A_0/\hbar)\int_0^T L_1
                                  +\mathcal A_1\int_0^T L_0,\\
 P_{2l}&\leq(\mathcal A_0/\hbar)\int_0^T L_2
                    +2\mathcal A_1\int_0^T L_1
                    +\hbar\mathcal A_2\int_0^T L_0.
\end{align*}
Extending these nonnegative integrals back from $T_g$ to zero is
conservative. The static propagator commutes with $p$.
Using the endpoint triples over the duration $T$ in
\eqref{old:Rfunctional} gives $<2.817683\times10^{-29}$.
Thus neither remote exact clock tail has been set to zero.

\subsection{Final old current, terminal CDF, and reproducible arithmetic}
\label{old:current}
Each of the five component expressions includes its own quadratic
$2e_1^2$ term. For their sum use $(\sum_{j=1}^5e_{1j})^2
\leq5\sum_{j=1}^5e_{1j}^2$. Add four further copies of the five
quadratic terms to their sum. This extra allowance is
$<2.117822196256773\times10^{-27}$.
Together with \eqref{old:referencecurrent} this gives
\begin{equation}
 I_{\rm old}<6.655453611235769\times10^{-16}.
 \label{old:Ifinal}
\end{equation}
All absolute values precede the live joint-coordinate integral.
The components are wave-law integrals; the original-law cap is inserted
only in the event conversion.

The scalar terminal comparison, on $T=0.04$, is independently bounded by
\begin{equation}
 D_{\rm end}=D_{\rm old}+\frac{\hbar Tp_2}{2M}
 +\frac{Tp_1\hbar N_1}{MvT_0}
 +\frac{\hbar A_0}{2Mv^2T_0}\,
       \frac52\left(1+\frac{2(\sigma+10^{-10})}{vd}\right)
 +\frac{T\hbar N_2}{2Mv^2T_0^2}+\Delta_0.
 \label{old:terminalD}
\end{equation}
The gate coefficient is a supremum over the whole clock packet, not the
single-offset integral. It is below $55$ and gives
$D_{\rm end}<3.034377484488931\times10^{-11}$.
For each radial cut the exact conditional CDF $K_\Psi$ and helper CDF
$K_u$ obey
\[
 \int w(S)|K_\Psi(S,b)-K_u(S,b)|\dd S
 \leq\|\,|\Psi|^2-|G|^2\,\|_1\leq2D_{\rm end}.
\]
Indeed subtract the two densities against the centered indicator
$\one_{r<b}-K_u(S,b)$, of modulus at most one. Under the exact terminal
wave law $K_\Psi$ is uniform conditional on $S$; equivariance transports
the original cap, giving $2CN D_{\rm end}$ for all $N=256001$ cuts.
This is terminal domination by the exact wave, never a newly imposed
cap relative to the helper.

All displayed decimal bounds can be checked by the following finite
rational procedure; the formulae above, rather than abbreviated decimals,
are the inputs. For $0\leq t\leq4/3$ use the 24-term alternating Taylor
sums for sine and cosine with the next-term one-sided remainder. Bound
square roots by rational bisection, and logarithms by
\[
 \log x=2\sum_{k=0}^{n-1}\frac{z^{2k+1}}{2k+1}+R_n,
 \quad z=(x-1)/(x+1),\quad
 |R_n|\leq\frac{2|z|^{2n+1}}{(2n+1)(1-z^2)},
\]
after taking reciprocal arguments where useful. Machin's identity
$\pi=16\arctan(1/5)-4\arctan(1/239)$ with alternating remainders
bounds $\pi$. Rational arithmetic permits arbitrary refinement.
For \eqref{old:sorting} use the breakpoints $0,s_-/2^{20},g$ and
64 equally spaced subintervals in every
$[s_-/2^j,s_-/2^{j-1}]$, $j=20,\ldots,1$.
Take every nonnegative numerator factor at the right endpoint and the
increasing lower gap at the left. No monotonicity of their ratio is
assumed. For \eqref{old:ageunion} use 128 equal subintervals and right
endpoint rectangles for the increasing product $Dq$.
These give
\begin{align*}
 \int_0^{s_-}\frac{G_+'D^2+Dq}{G_-}
       &<0.000517506143180416981,\\
 \int_{s_-}^{s_+}Dq&<0.000001583255601884495.
\end{align*}

For a positive system $y'=Ay$, rescale by a positive diagonal matrix
$S$, put $B=S^{-1}AS$, $z_0=S^{-1}y_0$, and sum
$z_k=(TB)^kz_0/k!$.
If $\alpha=T\|B\|_\infty<K+2$, the omitted norm after degree $K$
is at most
\[
 \|z_K\|_\infty\frac{\alpha}{K+1}
                     \frac1{1-\alpha/(K+2)}.
\]
For the time integral sum $Tz_k/(k+1)$ and multiply this tail by
$T/(K+2)$. For the homogeneous system append the constant state one,
use scales
$(10^{-29},10^{-47},10^{-28},10^{-35},10^{-33},10^{-40},1)$
and $K=80$. For the combined global/local tail system, ordered
$(P_0,\ldots,P_4,L_0,\ldots,L_3)$, use
$(1,10^{-19},10^{-38},10^{-57},10^{-76},10^{-42},10^{-60},10^{-68},10^{-72})$
and $K=100$. In both cases $\alpha<10$ for $T=0.04$.
This states every matrix entry, initial value, quadrature partition and
remainder rule needed to evaluate the bound without a separate
coefficient ledger or numerical propagation of the Schr\"odinger equation.


\section{Complete prefix estimates for the compensated writer}
\label{prefix:appendix}\label{prefix:main}

This appendix closes the clock and pointer derivative estimates used before
$t_a=0.0366\,\mathrm s$.  All norms are full spinor Hilbert norms, including
both internal sectors.  Write $t=0$ at $t_h=-0.002\,\mathrm s$, and set
\[
 \begin{aligned}
 T&=0.0386\,\mathrm s,&T_0&=0.03\,\mathrm s,&
 a&=vT_0=0.03\,\mathrm m,\\
 v&=1\,\mathrm{m/s},&M&=10\,\mathrm{kg},&
 \mu&=\mu_Z=0.01,\qquad\kappa=\hbar/M.
 \end{aligned}
\]
Clock derivatives are physical $S$ derivatives; pointer derivatives are
with respect to the dimensionless $Z$.  All scalar numbers below use SI
units for the clock.  Removing the same constant Galilean carrier from
both exact waves leaves the clock generator $vp_S+p_S^2/(2M)$.
The preparation estimates are those of
Sections~\ref{prep:transfer} and~\ref{prep:transfer-hessian}; the old-wave
bounds are obtained from Sections~\ref{old:homogeneous}
and~\ref{old:clockcoeff}.  Thus neither wave nor actual law is replaced
at the handoff.

\subsection{Finite Gaussian coefficient calculation}
\label{prefix:gaussian}
Let $w=0.003$, $b(s)=24B_9(s/w)$ with its constant extensions, and
\[
 \gamma(s,Z)=\pi^{-1/4}e^{-(Z-b)^2/2}
                   e^{\mathrm i\mu b'(Z-b/2)},\qquad y=Z-b.
\]
Its first logarithmic derivative is
\[
 \partial_s\log\gamma=c y+d,
 \quad c=b'+\mathrm i\mu b'',\qquad
 d=\tfrac{\mathrm i\mu}{2}(bb''-b'^2).
\]
The scalar phase in $d$ is retained.  Successive logarithmic derivatives
are $cy+d$, $c'y-cb'+d'$, $c''y-2c'b'-cb''+d''$, and
$c'''y-3c''b'-3c'b''-cb'''+d'''$.
The following rational derivative ceilings specify every coefficient
used below:
\[
 \begin{aligned}
 (B_0,B_1,B_2)&=(24,19687.5,105000000),\\
 (B_3,B_4,B_5)&=(4.9\times10^{11},4.032\times10^{17},
                                      4.7936\times10^{20}),
 \end{aligned}
 \qquad |b^{(j)}|\le B_j.
\]
Here $B_1=24(315/128)/w$ follows from
$B_9'=630x^4(1-x)^4$.  Differentiating this identity, using
$x(1-x)\le1/4$ and $|1-2x|\le1$, gives
$B_2=24(2520/64)/w^2$ and
$B_3=24(7560/16+5040/64)/w^3$.
For $j=4,5$ the sufficient coefficient bound is
\[
 B_j=\frac{24}{w^j}
       \sum_{k=5}^9 |\beta_k|\frac{k!}{(k-j)!},
 \qquad (\beta_5,\ldots,\beta_9)=(126,-420,540,-315,70).
\]
The polynomial and constant extensions are globally $C^4$; their fifth
weak derivative is bounded.  No endpoint distribution enters these
calculations.

Define $C_j=B_j+\mu B_{j+1}$ for $0\le j\le4$ and $q_0=C_0+1$.
The $C_j$ bound the derivatives of $b+\mathrm i\mu b'$.
Define four positive affine polynomials in a formal variable $Y$:
\begin{align*}
 d_0&=\mu(B_0B_2+B_1^2)/2,&
 d_1&=\mu(B_0B_3+B_1B_2)/2,\\
 d_2&=\mu(B_0B_4+B_2^2)/2,&
 d_3&=\mu(B_1B_4+B_0B_5+2B_2B_3)/2,\\
 L_1&=d_0+C_1Y,&L_2&=C_1B_1+d_1+C_2Y,\\
 L_3&=2C_2B_1+C_1B_2+d_2+C_3Y,&&\\[-2mm]
 L_4&=3C_3B_1+3C_2B_2+C_1B_3+d_3+C_4Y.&&
\end{align*}
For a positive polynomial $P(Y)=\sum p_jY^j$, put
\[
 \mathcal N(P)=\sum_jp_j\sqrt{(2j-1)!!/2^j},\qquad(-1)!!=1.
\]
This uses the exact Gaussian moment $\|y^j\gamma\|^2=(2j-1)!!/2^j$.
Complete, evaluable formulas for the four pointer jets are
\begin{align}
 g_0&=1,&g_1&=\mathcal N(L_1),&
 g_2&=\mathcal N(L_1^2+L_2),\nonumber\\
 g_3&=\mathcal N(L_1^3+3L_1L_2+L_3),&&&\nonumber\\[-1mm]
 g_4&=\mathcal N(L_1^4+6L_1^2L_2+3L_2^2+4L_1L_3+L_4),
 &&&\label{prefix:jets}
\end{align}
so that $\|\partial_s^j\gamma\|\le g_j$.  In particular,
\[
 \begin{aligned}
 g_1&<1.529437156625049\times10^7,&
 g_2&<2.344303538271762\times10^{14},\\
 g_3&<3.601493635795900\times10^{21},&
 g_4&<5.545932284532096\times10^{28}.
 \end{aligned}
\]
The defining expressions, rather than these shortened displays, are
used in the numerical bounds below.

For $A=\partial_Z+Z-\Pi_1(b+\mathrm i\mu b')$ and
$\mathsf Q(f)=(\|Zf\|^2+\|p_Zf\|^2)^{1/2}$, the oscillator form identity
and $[A,Z]=1$ give
\begin{align}
 \mathsf Q(f)&\le\|Af\|+q_0\|f\|,\nonumber\\
 \mathsf Q(Zf)&\le\|A^2f\|+2q_0\|Af\|+(q_0^2+1)\|f\|.
 \label{prefix:oscform}
\end{align}
They hold fibrewise and then for the Hilbert direct sum over all other
coordinates.  Differentiating $A\gamma=0$ gives the weighted jet bounds
\[
 q_j=q_0g_j+\sum_{k=1}^j\binom jk C_kg_{j-k},\qquad 1\le j\le3,
 \quad \mathsf Q(\partial_s^j\gamma)\le q_j.
\]
Finally, for the writer coefficients $F=b/\mu+\mu b''$ and
$C_b=b^2/(2\mu)+\mu bb''/2$, define
\begin{align}
 F_*&=B_0/\mu+\mu B_2=1052400,\nonumber\\
 F_1&=B_1/\mu+\mu B_3=4901968750,\nonumber\\
 F_2&=B_2/\mu+\mu B_4=4032010500000000,\nonumber\\
 K_1&=B_0B_1/\mu+\mu(B_1B_2+B_0B_3)/2=69183187500,\nonumber\\
 K_2&=(B_1^2+B_0B_2)/\mu
       +\mu(B_2^2+2B_1B_3+B_0B_4)/2
       =48535884509765625.
 \label{prefix:forces}
\end{align}
These bound $|F|,|F'|,|F''|,|C_b'|,|C_b''|$, respectively.

\subsection{Exact intertwiner and pulse integrals}
\label{prefix:intertwiner}
Let $G_*=V(S)(\psi_o g_0)$, where $\psi_o$ is the full exact old wave,
and $V$ applies the above Weyl displacement in sector one and the
identity in sector zero.  Since $\mathrm i\hbar v\gamma_S=H_Z\gamma$,
product differentiation gives, up to an irrelevant overall sign in the
forced error equation,
\[
 R/\hbar=\kappa(\gamma_S\psi_{o,S}+\gamma_{SS}\psi_o/2).
\]
This identity holds separately in the orthogonal sectors and hence for
any normalized qubit input.  Its support is exactly the open writer
pulse; the stationary displaced plateau contributes no source.
Write $\delta=\Psi-G_*$ and
$I_j=\int_0^T\|\mathbf1_{\rm pulse}\partial_S^j\psi_o\|\,dt$.
The old characteristic helper has clock centre $t_{\rm lab}+0.104$ and
half-width $0.0010000001$.  Its earliest pulse contact is
$0.0364499999\,\mathrm s$, leaving $\Delta=0.0001500001\,\mathrm s$.
Throughout those late slices the initial offset lies in an endpoint strip
of relative width $d=0.1500002$.  With $\sigma=0.001$ and
$m_\phi=1.000001$, define the rational bounds
\[
 \beta=m_\phi^2\frac{64}{35}(11/7)^8\frac{d^9}{9},\qquad
 \eta^2=m_\phi^2\frac{1024}{35}(11/7)^8
                         \frac{d^7}{7\sigma^2}.
\]
The inequality $\cos(\pi x/(2\sigma))\le
\pi(\sigma-x)/(2\sigma)$, positive convolution and Jensen's inequality
prove that $\beta$ bounds the strip mass and $\eta$ its amplitude
first-derivative norm.  The mollification radius is included in $d$.
On these slices $\|u_S\|\le N_1/a$, $N_1=34359738369$.
The old-wave estimates, evaluated through elapsed time $0.0404$,
give $\|\psi_o-G_o\|<3.1\times10^{-11}$ and
$\|\partial_S(\psi_o-G_o)\|<36000$.  Consequently
\begin{align}
 I_0&\le T(3.1\times10^{-11})+\Delta\sqrt\beta=:J_0,\nonumber\\
 I_1&\le T(36000)+\Delta(\eta+N_1\sqrt\beta/a)=:J_1,\nonumber\\
 I_2&\le T(1.315\times10^{24})=:J_2.
 \label{prefix:pulse}
\end{align}
The last bound uses the full old Hessian, including its error.
The preparation transfer and old right collar imply
$\|\delta(0)\|\le3\times10^{-30}+1.2\times10^{-37}$; the latter
term bounds $(V-I)\psi_og_0$ without discarding the tail.
Duhamel therefore proves the uniform ceiling
\begin{equation}
 \begin{aligned}
 D_*&:=3\times10^{-30}+1.2\times10^{-37}
  +\widehat\kappa(g_1J_1/a+g_2J_0/(2a^2)),\\
 \|\delta(t)\|&\le D_*<5.055628819933504\times10^{-22}.
 \end{aligned}
 \label{prefix:D}
\end{equation}
where $\widehat\kappa=1.055\times10^{-35}>\kappa$.
All subsequent sharp bounds use the defining expression for $D_*$.

\subsection{Closing the simultaneous pointer hierarchy}
\label{prefix:hierarchy}
For this first, deliberately coarse closure take
$h_-=1.054\times10^{-34}<\hbar$, and use $\widehat\kappa$ in positive
numerators and $h_-$ in denominators.  The old physical force ceilings
are $F_o=8.00000000001\times10^{-6}$ and $F_{o,2}=1$.
The handoff bounds proved in the preparation appendix are
\begin{center}
\begin{tabular}{lc}
\toprule
Quantity & Outward upper bound\\
\midrule
$\mathsf Q(\delta(0))$ & $6\times10^{-24}$\\
$\|\delta_S(0)\|$ & $0.073$\\
$\mathsf Q(\delta_S(0))$ & $300000$\\
$\|\Psi_{SS}(0)\|$ & $2.5\times10^{27}$\\
$A_h=\|A\Psi(0)\|$ & $5.049592328644200595\times10^{-24}$\\
$A_{2,h}=\|A^2\Psi(0)\|$ & $1.228242965202911600\times10^{-17}$\\
$U_h=\|\partial_S(A\Psi)(0)\|$ & $276122.9698216115699$\\
\bottomrule
\end{tabular}
\end{center}
These already include the remote Weyl change.
Differentiating the residual once gives the coefficients $1,3/2,1/2$.
Define its integrated bounds
\begin{align*}
 R_Q&=\widehat\kappa(q_1J_1/a+q_2J_0/(2a^2)),\\
 R_S&=\widehat\kappa(g_1J_2/a+3g_2J_1/(2a^2)+g_3J_0/(2a^3)),\\
 R_{SQ}&=\widehat\kappa(q_1J_2/a+3q_2J_1/(2a^2)+q_3J_0/(2a^3)).
\end{align*}
Oscillator rotation and the differentiated error equation give the
successive bounds
\begin{align*}
 Q_d&=6\times10^{-24}+F_*TD_*/T_0+R_Q,\\
 e_S&=0.073+F_oTD_*/h_-+T(F_1Q_d+K_1D_*)/(vT_0^2)+R_S,\\
 P_1&=1.146\times10^{12}+g_1/a+e_S,\\
 A_1&=A_h+\widehat\kappa\left[
 \frac{C_1}{a}(J_1+g_1J_0/a+Te_S)
                +\frac{C_2}{2a^2}(J_0+TD_*)\right].
\end{align*}
In particular $Q_d<6.847\times10^{-16}$, $e_S<1.482\times10^6$,
and $A_1<6.192373945249072\times10^{-23}$.
The last row follows from the exact identity
\begin{equation}
 (\mathrm i\partial_t-H/\hbar-1/(\mu T_0))A\Psi
 =-\kappa\Pi_1(c'\Psi_S/a+c''\Psi/(2a^2)),
 \qquad c=b+\mathrm i\mu b'.
 \label{prefix:Aeq}
\end{equation}
For the second power the exact equation is
\[
 (\mathrm i\partial_t-H/\hbar-2/(\mu T_0))A^2\Psi
 =-2\kappa\Pi_1(c'\partial_S(A\Psi)/a+c''A\Psi/(2a^2))
                     -\kappa\Pi_1(c')^2\Psi/a^2.
\]
The last source is essential and comes from $[A,\partial_S]=\Pi_1c'/a$.

For candidate constants $(X,P,U,B)$ bounding
$(\mathsf Q(\delta_S),\|\Psi_{SS}\|,\|\partial_S(A\Psi)\|,
\|A^2\Psi\|)$, put
\[
 Q_2=(q_0^2+1)D_*+2q_0A_1+B,\qquad
 P_{1Q}=q_0(1.146\times10^{12})+q_1/a+X.
\]
A simultaneous constant supersolution consists of the following four
strict inequalities:
\begin{align}
 X&>300000+F_*Te_S/T_0+F_oTQ_d/h_-
             +T(K_1Q_d+F_1Q_2)/(vT_0^2)+R_{SQ},\nonumber\\
 P&>2.5\times10^{27}+T\left[
 \frac{2F_oP_1+F_{o,2}}{h_-}
 +\frac{2(F_1P_{1Q}+K_1P_1)}{vT_0^2}
 +\frac{F_2(q_0+Q_d)+K_2}{v^2T_0^3}\right],\nonumber\\
 U&>U_h+T\left[\frac{F_oA_1}{h_-}
  +\frac{F_1(B+q_0A_1)+K_1A_1}{vT_0^2}
  +\widehat\kappa\left(\frac{C_1P}{a}
      +\frac{3C_2P_1}{2a^2}+\frac{C_3}{2a^3}\right)\right],\nonumber\\
 B&>A_{2,h}+T\widehat\kappa
         (2C_1U/a+C_2A_1/a^2+C_1^2/a^2).
 \label{prefix:fourrows}
\end{align}
These follow from the twice differentiated Schr\"odinger equation,
\eqref{prefix:Aeq}, and \eqref{prefix:oscform}.  In particular the force
on $A\Psi$ costs $B+q_0A_1$, not merely $A_1$.
For $(X,P,U,B)=(10^{14},10^{40},2\times10^{11},2\times10^{-17})$
the respective right sides are strictly below
\[
 \begin{aligned}
 \mathcal R_X&<4.011662441023\times10^{12},&
 \mathcal R_P&<6.718077927095\times10^{39},\\
 \mathcal R_U&<1.452034044217\times10^{11},&
 \mathcal R_B&<1.809106526629\times10^{-17}.
 \end{aligned}
\]
Every dependence on the other three unknowns has been retained.
The positive integral inequalities and a simultaneous first-crossing
argument prove the four bounds.  The large $10^{40}$ bound is used
only inside this hierarchy; the next subsection supplies the much
sharper error Hessian needed for the clock current.

\subsection{Old third clock derivative with complete coefficients}
\label{prefix:oldthird}
Only an ordinary third spatial clock derivative of the old wave is
needed; no fourth radial energy graph is invoked.
Sections~\ref{prep:hessian} and~\ref{prep:fourth} give the following
outward bounds for the gauged old preparation error $E_g$:
\[
 \begin{aligned}
 P_2&=3.238008492819540878\times10^{-49},\\
 P_4&=1.075566867715738537\times10^{-76},\\
 n_2&=5.970369147764257548\times10^{-38},
 \end{aligned}
\]
where $\|p_S^jE_g\|\le P_j$ and $n_2$ bounds the second left collar.
For a cutoff equal to one through $0.10005$, zero from $0.10015$, and
$g=0.0001$, the localized interpolation identities give
\[
 L_1=\sqrt{n_2P_2}+10\hbar_+ n_2/g,\qquad
 L_2=\sqrt{n_2}\bigl(\sqrt{P_4}+18\hbar_+\sqrt{P_2}/g\bigr),
 \quad \hbar_+=1.055\times10^{-34}.
\]
This cutoff covers every preparation-phase derivative, whose support
ends at $0.10001$, and lies inside the full-one region of the old
second collar, which ends at $0.100249999873$.
For the preparation-rate jets use
\[
 \begin{aligned}
 (a_0,a_1,a_2)&=(140,43600,70141750),\\
 (a_3,a_4)&=(1387431060000,986084475600000),\\
 \mu_r^+&=1.425\times10^{-12},\qquad F_p=0.080802,
 \end{aligned}
\]
and put $r_j=\mu_r^+F_pa_j/2$ for $j=1,2,3$.
Expanding $p_S^3(e^{\mathrm i\theta}E_g)$, with
$\hbar|\partial_S^j\theta|\le r_j$, proves the handoff bound
\begin{align}
 E_{3,h}:={}&h_-^{-3}\left[
 \sqrt{P_2P_4}+3r_1L_2+3(r_1^2+\hbar_+r_2)L_1\right.\nonumber\\[-1mm]
 &\left.\hspace{22mm}
 +(r_1^3+3\hbar_+r_1r_2+\hbar_+^2r_3)n_2\right]
                         +46933757821.201.
 \label{prefix:oldinitial}
\end{align}
The last term bounds the original packet's third derivative.  The
helper's radial state is $S$ independent at handoff and its preparation
phase is zero there.  Thus
$\|\psi_{o,SSS}(0)\|\le E_{3,h}<2.250395507380083\times10^{42}$.

Here are the old-wave input ceilings used on the enlarged horizon
$T_o=0.0404$, computed with the formulas of
Sections~\ref{old:homogeneous} and~\ref{old:clockcoeff}:
\begin{center}
\begin{tabular}{lc}
\toprule
Input & Outward upper bound\\
\midrule
$D_o=\sup\|\psi_o-G_o\|$ & $3.034654348976846863\times10^{-11}$\\
$e_{1,o}=\sup\|\partial_S(\psi_o-G_o)\|$ & $35444.22634847288176$\\
$e_{2,o}=\sup\|\partial_S^2(\psi_o-G_o)\|$ & $1.047307159416726327\times10^{21}$\\
$G_1=\sup\|G_{o,S}\|$ & $1145324612300.000003$\\
$G_2=\sup\|G_{o,SS}\|$ & $1.313080270080687692\times10^{24}$\\
$F_{c,1}$ & $4.388993584524007892\times10^{-6}$\\
$F_{c,2}$ & $0.08054047024909008647$\\
\bottomrule
\end{tabular}
\end{center}
For clarity the enlargement keeps the same ramp interval $0.0033$ and
uses static duration $T_o-0.0033=0.0371<0.0377$.
The helper age remains below $4/3<\pi/2$.
Its packet ceilings are
\[
 p_1=1.000001\sqrt{16/7}\,\pi/(2\sigma),\quad
 p_2=1.000001\sqrt{512/35}\,(\pi/(2\sigma))^2,
\]
so that, with $N_2=1.001(2^{35}+1)^2$,
\[
 G_1\le\sqrt{p_1^2+(N_1/a)^2},\qquad
 G_2\le p_2+2p_1N_1/a+N_2/a^2+25000N_1/a.
\]
The three error entries are the sums of the transported preparation,
short ramp, static, common-to-clipped and clipped-to-static component
bounds evaluated at these durations; no new propagation assumption is
introduced by extending the time argument.

The remote third force, including the preparation phase-square term,
has the complete upper expression
\begin{align}
 F_{c,3}={}&\mu_r^+a_4F_p/2
  +\mu_r^+(3a_1a_2+a_0a_3)(0.201)^2\nonumber\\
 &+\frac{(\mu_r^+)^2}{4M}(3a_2a_3+a_1a_4)F_p^2
                      +2.393692459677192\times10^{-27}.
 \label{prefix:Fc3}
\end{align}
The last summand is the third quiet-force derivative obtained from the
finite product-rule calculation in Section~\ref{old:quiet}.
Thus $F_{c,3}<68.48117358490407$.
Put
\[
 \begin{aligned}
 W_0&=\hbar_+(2^{35}+1)/T_0,&F_{a,1}&=25000W_0,\\
 F_{a,2}&=\frac{2520W_0}{64(0.0001)^2},&
 F_{a,3}&=\frac{(7560/16+5040/64)W_0}{(0.0001)^3}.
 \end{aligned}
\]
The remote common-minus-clipped force is supported on $S\le0.10001$.
The helper and all its $S$ derivatives vanish there for the entire
prefix, since its left support starts at $0.1009999999$ and moves right.
The third differentiated unitary equation therefore gives
\begin{align*}
 \frac{d}{dt}\|\psi_{o,SSS}\|\le R_3:={}&
 h_-^{-1}\bigl[3F_{c,1}e_{2,o}+3F_{c,2}e_{1,o}+F_{c,3}D_o\\
 &\qquad+3F_{a,1}(G_2+e_{2,o})
       +3F_{a,2}(G_1+e_{1,o})+F_{a,3}\bigr].
\end{align*}
In particular the large helper Hessian is not charged on the remote
force support.  Integrating the linear envelope once more yields
\begin{equation}
 J_3:=TE_{3,h}+\tfrac12T^2R_3
 <9.746858547719537661\times10^{46},\qquad I_3\le J_3.
 \label{prefix:J3}
\end{equation}

\subsection{Localized first derivative and refined mixed moment}
\label{prefix:localized}
Take a fixed decreasing Lipschitz cutoff $\chi=1$ on $S\le0.126$,
$\chi=0$ on $S\ge0.136$, with $|\chi'|\le1/g_L$, $g_L=0.01$.
Every old clock-force derivative is supported in its full-one region
($S\le0.106$), while the residual starts at $S=0.14145$.
The positive constant drift favors departure from this left region.
The forced continuity identity, with no source on the support of $\chi$,
gives
\begin{equation}
 n_L(t):=\|\chi\delta(t)\|
 \le n_{L,h}+\frac\kappa{g_L}\int_0^t\|\delta_S(u)\|\,du,
 \qquad n_{L,h}=7.901146211666425714\times10^{-36}.
 \label{prefix:left}
\end{equation}
Since $AG_*=0$, improve the oscillator bound to
$\widetilde Q=A_1+q_0D_*$.
Let $\widetilde R_S=(\kappa/\widehat\kappa)R_S$ and similarly
$\widetilde R_{SQ}=(\kappa/\widehat\kappa)R_{SQ}$.
For the candidate $e=0.1$, define
\begin{align}
 n_*&=n_{L,h}+\kappa Te/g_L,\nonumber\\
 \mathcal E&=0.073+F_oTn_*/\hbar
                  +T(F_1\widetilde Q+K_1D_*)/(vT_0^2)
                  +\widetilde R_S.
 \label{prefix:localfirst}
\end{align}
The positive two-variable comparison closes because
$\mathcal E<0.07327295767972046<e$.
Consequently $\|\delta_S\|\le0.1$ and $n_L\le n_*<1.197180\times10^{-35}$
throughout the prefix.  The improvement retains the old force but
multiplies it by its actual localized error norm.

On that force support $A=\partial_Z+Z$ and $AG_*=0$.  Applying the
oscillator identity on each $S$ slice, without differentiating a sharp
indicator, gives
$\mathsf Q(\mathbf1_{\rm old\ force}\delta)
 \le\sqrt{A_1^2+n_*^2}$.
Set $B=2\times10^{-17}$ in $Q_2$ above.  The differentiated oscillator
rotation equation now proves the much sharper uniform bound
\begin{align}
 X_*:={}&300000+F_*Te/T_0
 +F_oT\sqrt{A_1^2+n_*^2}/\hbar\nonumber\\
 &+T(K_1\widetilde Q+F_1Q_2)/(vT_0^2)
                    +\widetilde R_{SQ},
 \qquad\mathsf Q(\delta_S)\le X_*<616735.
 \label{prefix:refinedQ}
\end{align}
This is the mixed moment required in the clock estimate; the coarse
$10^{14}$ bound from \eqref{prefix:fourrows} is not substituted here.

\subsection{Second clock derivative and the two absolute currents}
\label{prefix:currents}
Two residual derivatives give exactly
\[
 R_{SS}/\hbar=\kappa\left(
 \gamma_S\psi_{o,SSS}+\tfrac52\gamma_{SS}\psi_{o,SS}
       +2\gamma_{SSS}\psi_{o,S}+\tfrac12\gamma_{SSSS}\psi_o\right).
\]
Its integrated norm is bounded by the explicit expression
\[
 R_{SS}^*=\kappa(g_1J_3/a+5g_2J_2/(2a^2)
                      +2g_3J_1/a^3+g_4J_0/(2a^4)).
\]
The refined new-minus-old handoff Hessian from
Section~\ref{prep:transfer-hessian}, plus both derivatives of the
remote Weyl change, gives the following expression.
Here $n_{2,R}=5.970369147764257548\times10^{-38}$ is the
\emph{right} second-collar ceiling from Section~\ref{prep:transfer}.
It equals the numerical left-collar ceiling $n_2$ because reflecting
the moving-cutoff construction gives the same positive momentum
inequality: the right boundary translates at $v$ plus the drift ceiling
and the left boundary at $v$ minus that ceiling.  The right collar is
one throughout the writer support; its handoff full-one edge is below
$0.103750000127$, while $S_{\rm on}=0.14145$.
Thus the required right-tail amplitude is bounded independently of the
left-tail localization used in \eqref{prefix:oldinitial}.  Explicitly,
\begin{align}
 H_h={}&1.828509471006939478\times10^{19}
       +\frac{2(2.534074040624599816\times10^{-57})}{\hbar^2}
       \nonumber\\
 &+\frac{2g_1(1.390399439216243060\times10^{-43})}{a\hbar}
                       +\frac{g_2n_{2,R}}{a^2}.
 \label{prefix:Hh}
\end{align}
For $H(t)=\|\delta_{SS}(t)\|$, integration by parts with $\chi$ gives
\[
 \|\mathbf1_{\rm old\ force}\delta_S\|
             \le\sqrt{n_*H(t)}+2D_*/g_L.
\]
The cutoff cost is retained.  Twice differentiating the error equation
therefore yields the differential comparison
$H^\prime\le 2\alpha\sqrt H+f(t)$, where
$\alpha=F_o\sqrt{n_*}/\hbar$, $f\ge0$, and
$H(0)+\int_0^T f(t)\,dt\le\mathcal R$ with
\begin{align}
 \mathcal R={}&H_h+R_{SS}^*
       +\frac{4F_oTD_*}{\hbar g_L}
       +\frac{F_{o,2}Tn_*}{\hbar}\nonumber\\
 &+\frac{2T(F_1X_*+K_1e)}{vT_0^2}
       +\frac{T(F_2\widetilde Q+K_2D_*)}{v^2T_0^3}.
 \label{prefix:Hrest}
\end{align}
Indeed, writing $F(t)=H_h+\int_0^t f(u)\,du$,
the function $[\alpha t+\sqrt{F(t)}]^2$ has derivative at least
$2\alpha[\alpha t+\sqrt{F(t)}]+f(t)$ and dominates the initial
value.  Scalar differential comparison therefore proves
\begin{equation}
 H(t)\le H_*:=\left(F_oT\sqrt{n_*}/\hbar+\sqrt{\mathcal R}\right)^2
 <1.117214708915152634\times10^{21}.
 \label{prefix:H}
\end{equation}
Regularization at zero justifies this comparison when a norm vanishes.
Smooth-domain approximation and
these uniform finite estimates justify the oscillator and weak clock
commutations without an oscillator cutoff.

On $S<S_{\rm on}$, $G_*=\psi_og_0$ and $A=\partial_Z+Z$.
The oscillator form and graph identities imply, after restriction to
that region,
\[
 \|\delta_Z\|\le e_Z:=\sqrt{A_1^2+D_*^2},\qquad
 \|\delta_{ZZ}\|\le e_{ZZ}:=\sqrt{B^2+4A_1^2+3D_*^2}.
\]
For example $\|p_Z^2f\|^2\le\|(p_Z^2+Z^2)f\|^2+2\|f\|^2$ and
$\|(\mathcal B^\dagger\mathcal B+1)f\|^2
=\|\mathcal B^2f\|^2+4\|\mathcal Bf\|^2+\|f\|^2$ with
$\mathcal B=\partial_Z+Z$ prove the second bound.
The conditional-rank stability functional then gives
\begin{equation}
 I_Z\le\frac{T}{\mu T_0}
 \left(2\sqrt2e_Z+2e_Z^2+\frac{23\sqrt3}{2}D_*+e_{ZZ}\right)
 <2.574814383253091344\times10^{-15}.
 \label{prefix:IZ}
\end{equation}
Here the comparison's radial direction is independent of $Z$;
$\|g_0'\|=1/\sqrt2$, $\|g_0''\|=\sqrt3/2$ and
$\|g_0'^2/g_0\|=\sqrt3/2$ supply the displayed coefficients.
For the clock functional use the old characteristic helper instead.
Expanding its positive stability majorant with the additional error
$\delta$ gives
\begin{align}
 \Delta I_S&\le\kappa T\left[
 4(G_1+e_{1,o})e+2e^2+14D_*Q_G+9D_*G_2+H_*\right],\nonumber\\
 Q_G&=1.000001\frac{4\pi^2}{\sigma^2}\sqrt{3/35}+(N_1/a)^2,
 \nonumber\\[-1mm]
 \Delta I_S&<4.547786946512146880\times10^{-16}.
 \label{prefix:IS}
\end{align}
This is an increment of an absolute-current upper bound, not a
subtraction of signed currents.  All hidden-coordinate absolute values
precede integration.  Both estimates apply on the common whole-prefix
good-clock event; its complement is charged separately in the event
composition.  The current components must be added before optimizing
the single rank collar.

\paragraph{Reproducible outward arithmetic.}
Equations~\eqref{prefix:jets}--\eqref{prefix:IS}, together with the
outward input tables, form a finite coefficient specification.
Every displayed decimal input is an exact rational ceiling.
One fully rational evaluation uses a bracket for $\pi$ of width
$2\times10^{-69}$ about
$3.141592653589793238462643383279502884197169399375105820974944592307817$
and $\hbar=6.62607015\times10^{-34}/(2\pi)$.
For a nonnegative rational $x=n/d$, replace each square root by
\[
 \operatorname{sqrt}_+(x)=
 \frac{1+\left\lfloor\sqrt{\left\lfloor10^{200}n/d\right\rfloor}
                         \right\rfloor}{10^{100}}.
\]
All these coefficient expressions have positive inputs; use upper
numerators and lower denominators.  This procedure gives
$X_*<616734.119066647397144$,
$H_*<1.117214708915152634\times10^{21}$, and the strict current bounds
\eqref{prefix:IZ}--\eqref{prefix:IS}.  It specifies the calculations
inside the manuscript and requires no external coefficient file.


\section{Copying and retention for the exact loaded wave}\label{hold:appendix}\label{hold:main}

All estimates in this appendix concern the same normalized enlarged wave
$\Psi$ and its original actual law.  Write $C=270000/67499$,
$m=2^{52}$, $h=1/128$, $M=10\,\mathrm{kg}$, $v=1\,\mathrm{m/s}$,
$T_0=0.03\,\mathrm s$, $\ell_0=10^{-4}\,\mathrm m$, and
$\mu=\mu_Z=0.01$.  Radial and pointer coordinates are dimensionless;
$S$ and the elapsed time $t$ below have physical units.  The origin
$t=0$ is the handoff at laboratory time $-0.002\,\mathrm s$.
The transferred entrance estimates are those proved in
Section~\ref{prep:transfer}; no law or wave is assigned anew there.

\subsection{Gaussian derivatives and the complete clock residual}
Put $b(s)=24B_9(s/w)$, $w=0.003$, with the constant extensions specified
in the model.  Direct differentiation gives the following rational
ceilings $b_j\geq\|b^{(j)}\|_\infty$:
\begin{equation}\label{hold:bjets}
 (b_0,b_1,b_2,b_3,b_4)
 =\left(24,\frac{39375}{2},105000000,
            490000000000,403200000000000000\right).
\end{equation}
Indeed $B_9'=630x^4(1-x)^4$; the first three derivative bounds follow
from $x(1-x)\leq1/4$ and $|1-2x|\leq1$, and the coefficient sum of
$B_9^{(4)}$ is $1360800$.  These bounds apply on both constant
extensions because $B_9$ has four matching endpoint derivatives.

For the exact sector-one pointer put $y=Z-b$ and
\[
 g_b=\pi^{-1/4}e^{-y^2/2}
       e^{\ii\mu b'(Z-b/2)},\qquad
 c=b'+\ii\mu b'',\qquad d=\frac{\ii\mu}{2}(bb''-b'^2).
\]
Then $g_b'=(cy+d)g_b$ and
\[
 g_b''=\bigl((cy+d)^2+c'y-cb'+d'\bigr)g_b,
 \qquad d'=\frac{\ii\mu}{2}(bb'''-b'b'').
\]
In particular the scalar phase is retained.  Define
\begin{align}
 c_0&=b_1+\mu b_2,&c_1&=b_2+\mu b_3,\notag\\
 d_0&=\mu(b_0b_2+b_1^2)/2,&
 d_1&=\mu(b_0b_3+b_1b_2)/2,\notag\\
 g_1&=d_0+c_0/\sqrt2,&
 g_2&=d_0^2+c_0b_1+d_1
       +(2d_0c_0+c_1)/\sqrt2+\sqrt3\,c_0^2/2.
 \label{hold:gjets}
\end{align}
The Gaussian moments $\|y^j g_b\|^2=(2j-1)!!/2^j$ prove
$\|g_b^{(j)}\|\leq g_j$ for $j=1,2$.
The sector-zero pointer is constant, so the same ceilings apply to
both sectors, uniformly in the normalized qubit input.

The characteristic comparison through readout is
\begin{equation}\label{hold:comparison}
 \begin{aligned}
 G(t,S,r,Z)=\phi(S-0.102-vt)
       \sum_{i=0}^1\alpha_i u_i(\sigma,r)g_i(s,Z),\\
 \sigma=\frac{S/v-0.104\,\mathrm s}{T_0},\qquad
 s=\frac{S-0.14145\,\mathrm m}{vT_0}.
 \end{aligned}
\end{equation}
Here $u_i$ is the exact smooth-activation radial family, continued
backwards by the quiet Hamiltonian at negative ages.  On the entire
compact support of $\phi$, the preparation phase potential has ended.
Substitution into the boosted Schr\"odinger equation leaves precisely
the residual $\pm\hbar^2G_{SS}/(2M)$.  The sign is immaterial for
the norm estimates, but both pointer and activation derivatives in
$G_{SS}$ must be included.

Let $N_0=1$, $N_1=V+1$, $N_2=1001(V+1)^2/1000$ and
$N_3=1001(V+1)^3/1000$, where $V=2^{35}$; these are the energy-graph
bounds proved in Section~\ref{old:graphs}.  Set
\begin{align*}
 F_0&=\frac{2010^2}{2m\,100^4}+\frac{3}{2m\,100^2},&
 F_1&=F_0/2+\frac{2010}{m\,100^4},& W_1&=2^{46},\\
 K_j&=\sqrt{2m(N_jN_{j+1}+F_0N_j^2)},&&&j&=0,1,2.
\end{align*}
The form inequality $H_\gamma\geq p_r^2/(2m)-F_0$, applied to
$H_\gamma^ju_i$, proves $\|\partial_rH_\gamma^ju_i\|\leq K_j$.
Only the third Hamiltonian graph is needed.  Since
\[
 \partial_\sigma^2u_i=-H_\gamma^2u_i
       -\ii\gamma_\sigma'(W-F)u_i,
 \qquad 0\leq\gamma_\sigma'\leq750,
\]
the product-rule bounds, including the gate derivative, are
\begin{align*}
 U_1&=N_1+g_1,\\
 U_2&=N_2+750(V+F_0)+2N_1g_1+g_2,\\
 L_0&=K_0,\\
 L_1&=K_1+g_1K_0,\\
 L_2&=K_2+750\{W_1+F_1+(V+F_0)K_0\}
                         +2g_1K_1+g_2K_0.
\end{align*}
For the normalized mollified clock packet, convolution contraction
and its normalization factor at most $1.000001$ give
\[
 p_1=1.000001\sqrt{16/7}\,\frac{\pi}{0.002},\qquad
 p_2=1.000001\sqrt{512/35}\left(\frac{\pi}{0.002}\right)^2
\]
as ceilings on $\|\phi'\|$ and $\|\phi''\|$.  Therefore
\begin{equation}\label{hold:residualjets}
 G_2=p_2+\frac{2p_1U_1}{vT_0}+\frac{U_2}{(vT_0)^2},\qquad
 G_{2r}=p_2L_0+\frac{2p_1L_1}{vT_0}+\frac{L_2}{(vT_0)^2}
\end{equation}
bound $\|G_{SS}\|$ and $\|\partial_rG_{SS}\|$.

Write $E=\Psi-G$.  Its entrance norm $e_h$ is $D_h$ from
Section~\ref{prep:transfer} plus the backward-quiet discrepancy.
For that discrepancy, if $e=(e^{-\ii\sigma H_F}-1)\chi$ and
$|\sigma|\leq1/8$, then
\[
 \|e\|\leq\|H_F\chi\|/8,\quad
 \|H_Fe\|\leq\|H_F^2\chi\|/8,\quad
 \|\partial_re\|^2\leq2m\|e\|(\|H_Fe\|+F_0\|e\|).
\]
The quiet tail estimates in Section~\ref{old:quiet} make its
physical momentum correction less than $10^{-100}\,\mathrm{kg,m/s}$.
Let $P_{r,h}$ be the old ordinary radial error momentum, the
loaded-minus-old momentum from Section~\ref{prep:transfer}, and this
quiet correction added together.

For completeness the global radial force used in propagating $E$ is
explicit.  Put $m_{\rm phys}=m\hbar T_0/\ell_0^2$,
$\omega=1/(mT_0)$, $R=0.201\,\mathrm m$, $a=140\,\mathrm{s}^{-1}$,
$a'=43600\,\mathrm{s}^{-2}$ and $F_{\max}=2R^2$.  Define
\begin{align*}
 Q_0&=m_{\rm phys}\omega^2R^2/2+3\hbar\omega/2,\\
 Q_r&=m_{\rm phys}\omega^2R+5000Q_0,&
 W_r&=\frac{\hbar(mh/2+5V)}{T_0\ell_0},\\
 V_{p,r}&=m_{\rm phys}(a'+a^2)R
          +\frac{m_{\rm phys}^2(a')^2F_{\max}R}{2M},&
 V_r&=V_{p,r}+Q_r+W_r.
\end{align*}
These follow by differentiating the displayed preparation, weak-trap
and capped detector potentials; all remote preparation terms remain.
The writer has zero radial derivative.  Unitary Duhamel, followed by
the ordinary radial momentum commutator, consequently gives
\begin{align}
 e(t)&\leq e_h+r_0t,&r_0&=\hbar G_2/(2M),\notag\\
 P_r(t)&\leq P_{r,h}+r_rt+V_r(e_ht+r_0t^2/2),&
 r_r&=\hbar^2G_{2r}/(2M\ell_0).
 \label{hold:radialerror}
\end{align}
The odd extension at the Dirichlet origin has no boundary term.
Approximation on the common smooth core and these uniform bounds
justify the differentiated evolution without a fourth radial graph.
At $T_c=0.0404\,\mathrm s$, denote the right sides by $e_*$ and
$P_{r,*}$ and put $p_*={\ell_0P_{r,*}}/({\hbar mh})$.
Evaluation gives
\[
 e_*<3.034680\times10^{-11},\qquad
 P_{r,*}<2.208802\times10^{-20}\,\mathrm{kg,m/s},\qquad
 p_*<0.000595293.
\]

\subsection{A closed clock and annihilator estimate through the hold}
\label{hold:clock}
The exact boosted generator contains $p_s+\epsilon p_s^2/2$,
$\epsilon=\hbar/(Mv^2T_0)$.  Define
$\mathcal A=\partial_Z+Z-\Pi_1(b+\ii\mu b')$.
Since the radial potentials commute with this operator, its complete
commutator is
\begin{equation}\label{hold:annihilator}
 (\ii\partial_\tau-H-1/\mu)\mathcal A\Psi
 =-\epsilon\Pi_1\bigl[(b'+\ii\mu b'')\partial_s
                          +(b''+\ii\mu b''')/2\bigr]\Psi.
\end{equation}
For example $[\mathcal A,H_{\rm ptr}]=(\mathcal A+\Pi_1c)/\mu
-\Pi_1F$ and $[\mathcal A,p_s]=-\ii\Pi_1c'$, where
$c=b+\ii\mu b'$; $c/\mu-F-\ii c'=0$ cancels the prescribed
driving terms.  The clock Laplacian supplies both terms on the right
of \eqref{hold:annihilator}.

Let $P(t)=\|-\ii\hbar\partial_S\Psi(t)\|$ be the boosted clock
momentum, $n(t)$ the moving left-collar norm of width
$g=0.0005\,\mathrm m$, and
$R_Z(t)=(\|Z\Psi\|^2+\|p_Z\Psi\|^2)^{1/2}$.
The entrance estimates allow the conservative values
\[
 P(0)\leq2\times10^{-27},\qquad n(0)\leq2\times10^{-26},
 \qquad R_Z(0)\leq2,
\]
in SI units where applicable; the much sharper actual entrance
$A_h=\|\mathcal A\Psi(0)\|$ is retained.
Set
\[
 F_*={b_0}/{\mu}+\mu b_2,\quad
 F_*'={b_1}/{\mu}+\mu b_3,\quad
 C_*'={b_0b_1}/{\mu}+\mu(b_1b_2+b_0b_3)/2.
\]
The Heisenberg equations for $(Z,p_Z)$ are a rotation with bounded
forcing $\Pi_1F$; its variation-of-constants formula gives
$R_Z(t)\leq2+F_*t/T_0$, without exponentiating the oscillator
frequency.  The translating collar and clock commutator give
\begin{equation}\label{hold:closedmoments}
 n'\leq\frac{P}{Mg},\qquad
 P'\leq F_p n+F_a+
       \frac{\hbar}{vT_0^2}\{F_*'R_Z+C_*'\},
 \qquad F_p=8\times10^{-6},\quad F_a=10^{-17}.
\end{equation}
Here $F_p$ and $F_a$ have units of force.  Explicit ceilings from
the original potentials, with $a''=70141750\,\mathrm{s}^{-3}$, are
\begin{align*}
 F_p^{\rm calc}&=
 \frac{m_{\rm phys}a''F_{\max}}{2v}
 +\frac{m_{\rm phys}aa'R^2}{v}
 +\frac{m_{\rm phys}^2a'a''F_{\max}^2}{4Mv^3}<F_p,\\
 Q_S&=2m_{\rm phys}\omega^2aR^2+3\hbar\omega a+5000Q_0,\\
 F_a^{\rm calc}&=Q_S+25000(\hbar V/T_0+Q_0)<F_a.
\end{align*}
The $C_*'$ term retains the clock force of the scalar compensation.
All inequalities apply to the full wave, including remote clock tails.
The oscillator estimates extend from finite-energy approximants by
the form-domain bounds just obtained.

For $q=0.0002\,\mathrm{kg,m/s}$ and $k=0.04\,\mathrm{s}^{-1}$,
$Y=P+qn$ satisfies $Y'\leq kY+f_0+f_1t$, where
\[
 Y_0=2\times10^{-27}+q(2\times10^{-26}),\qquad
 f_0=F_a+\frac{\hbar(F_*'\,2+C_*')}{vT_0^2},\qquad
 f_1=\frac{\hbar F_*'F_*}{vT_0^3}.
\]
Indeed $q/(Mg)=F_p/q=k$.  With $T=0.092\,\mathrm s$ and
$E_T=(1-kT)^{-1}\geq e^{kT}$, positive integration proves
\begin{align}
 P(t)&\leq E_T(Y_0+f_0T+f_1T^2/2),\notag\\
 \int_0^T P(t)\dd t&\leq
 \mathcal P:=E_T(Y_0T+f_0T^2/2+f_1T^3/6),\notag\\
 A_*&:=A_h+\frac{c_0\mathcal P}{MvT_0}
                +\frac{\hbar c_1T}{2Mv^2T_0^2}
       \geq\sup_{0\leq t\leq T}\|\mathcal A\Psi(t)\|.
 \label{hold:closedvalues}
\end{align}
The last line is unitary Duhamel applied to
\eqref{hold:annihilator}.  The resulting ceilings are
\[
 \mathcal P<2.667201\times10^{-18}\,\mathrm{kg,m},\qquad
 A_*<9.510237\times10^{-12}.
\]

Equivariance and Cauchy--Schwarz imply
$\E_{|\Psi|^2}\int_0^T|\dot S-v|\dd t\leq\mathcal P/M$.
Thus actual-law domination and Markov's inequality bound the failure
of the entire-path deviation $\int|\dot S-v|\leq10^{-5}\,\mathrm m$
by $C\mathcal P/(10^{-5}M)<1.067\times10^{-13}$.
The handoff core $|S_h-0.102|\leq0.000839\,\mathrm m$ has failure
at most
\begin{equation}\label{hold:core}
 e_{\rm core}=C(\sqrt{\beta_{\rm core}}+D_h)^2,
 \qquad
 \beta_{\rm core}=(1.000001)^2\frac{128}{35}
       \left(\frac{\pi}{2}\right)^8\frac{(0.1610001)^9}{9}.
\end{equation}
To see this, normalize the unsmoothed $\cos^4$ packet on its
$0.001\,\mathrm m$ half-width.  Its squared density has coefficient
$64/35$ after scaling to unit half-width; the two edges give the
factor $128/35$.  On either edge use
$\sin x\leq x$, integrate the eighth power, and enlarge the edge
width by the smoothing radius $10^{-10}\,\mathrm m$.
Positive convolution and Jensen preserve this two-sided bound,
with the displayed normalization factor.  The $D_h$ triangle then
uses the exact handoff wave, rather than replacing its marginal.

On the intersection of these two good events, the first possible
crossing of $S_{\rm on}$ is after $0.036601\,\mathrm s$, and the
last possible completion of the pulse is before
$0.038389\,\mathrm s$.  More directly, the pathwise bound
$|S_t-(0.104+vt_{\rm lab})|\leq0.000849\,\mathrm m$ implies
$S_t\geq0.14154\,\mathrm m$ for every
$t_{\rm lab}\in[0.0384,0.09]$.  This proves completion throughout
the whole hold, including the exclusion of later returns.  The same
two exceptional events are charged only once below.

\subsection{The actual earlier label and a positive capture event}
\label{hold:copy}
Let $H_r(r)\in[0,1]$ be the soft radial classifier, agreeing with the
sharp label $\ell(r)$ outside its bands $B$ and having
$|H_r'|\leq10/h$.  Put $I_0=[-8,8]$, $I_1=[16,32]$ and
\[
 M_r(r,Z)=(1-H_r(r))\one_{I_0^c}(Z)+H_r(r)\one_{I_1^c}(Z).
\]
For each absolutely continuous exact trajectory, with witness and
readout at laboratory times $t_a=0.0366$, $t_c=0.0384$,
\begin{equation}\label{hold:pathineq}
 \one_{\{Z_c\notin I_{\ell(r_a)}\}}
 \leq\one_B(r_a)+\Var_{[t_a,t_c]}H_r(r_t)+M_r(r_c,Z_c).
\end{equation}
Outside $B$, $H_r(r_a)$ is the binary earlier label; replacing this
coefficient by $H_r(r_c)$ changes the mismatch by at most its total
variation.  This proves the inequality without assigning a sampled
internal sector.

Throughout this copy interval the entire comparison packet has
conditional age in $[71/60,4/3]$.  Use the $B_9$-taper lens
expressions $D_9,Q_9$ of \eqref{old:DQ9}, evaluated at $4/3$,
and denote these by $D,Q$; set
\begin{equation}\label{hold:lensvalues}
 q_r=\frac{\sqrt{2mD(Q+F_0D)}}{mh}.
\end{equation}
They give $D<0.007640094$, $Q<262597634.037$ and
$q_r<0.003820673$.  The contraction scale is at most
$\cos(\sigma-1/300)$.  Since $\sigma-1/300\geq59/50$ and
$\cos(59/50)<2/5$, the compact comparison and its radial derivative
vanish on the bands $[h/5,3h/10]$ and their translates.  Hence
$\|\one_BG\|\leq D$ and $\|\one_B\partial_rG\|\leq mhq_r$.
Expanding the exact current of $G+E$, taking its modulus before
integrating any hidden variable, gives
\begin{equation}\label{hold:bandcurrent}
 \int_B|J_r[\Psi]|\leq
 h\{Dq_r+(D+e_*)p_*+e_*q_r\}.
\end{equation}
For example the four terms before division by $m$ are
$D(mhq_r)$, $D\|\partial_rE\|$,
$e_*(mhq_r)$ and $e_*\|\partial_rE\|$.

On the completed-clock region, let $P_g$ project in each internal
sector onto the normalized real Gaussian centered at $0$ or $24$.
The oscillator ladder spectrum proves the fibre inequality
\begin{equation}\label{hold:gap}
 \mathcal A^*\mathcal A\geq2(I-P_g).
\end{equation}
Both operators commute with completed-clock localization.  Thus the
localized excited component has norm at most $A_*/\sqrt2$.
The ground-state wrong-inner probability is $q_I=\operatorname{erfc}(8)$,
and contraction and the norm triangle give a pointer cost at most
$C(\sqrt{q_I}+A_*/\sqrt2)^2$.
To combine it with the radial mismatch, use the pointwise inequality
\[
 M_r(r,Z)\leq |H_r(r)-i|+\one_{I_i^c}(Z),\qquad i=0,1,
\]
and sum against the exact orthogonal component densities.
The square-root multiplier of the first term kills the compact
sector-$i$ lens comparison, so its total contribution is at most
$C(D+e_*)^2$.  Together with the initial band and
\eqref{hold:bandcurrent}, this proves the copy cost, apart from the
already identified exceptional events,
\begin{equation}\label{hold:copycost}
 C\left[2(D+e_*)^2+\frac35
       \{Dq_r+(D+e_*)p_*+e_*q_r\}\right]
       +C(\sqrt{q_I}+A_*/\sqrt2)^2.
\end{equation}
Here $10(t_c-t_a)/T_0=3/5$.  Each endpoint and variation refers to
the same exact trajectory and its own earlier radial label.

\subsection{Absolute current over the entire holding interval}
\label{hold:retention}
At each full configuration, including $S,r,Z$, the exact identity is
\begin{equation}\label{hold:exactcurrent}
 J_Z[\Psi]=\frac{1}{\mu}\Im(\Psi^\dagger\mathcal A\Psi)
                         +b'\Psi^\dagger\Pi_1\Psi.
\end{equation}
The second term has not been discarded: it vanishes on every visited
point of a good clock path because the pulse is completed there.
For $R_0=[-10,10]$, $R_1=[14,34]$, a path starting in $I_i$ and
leaving $R_i$ crosses one of the four width-two bands
$[-10,-8]$, $[8,10]$, $[14,16]$, $[32,34]$.
Linear cutoffs of slope $1/2$ on these disjoint bands require unit
variation for such an exit.  Restricted equivariance, actual-law
domination and then Cauchy--Schwarz therefore give
\begin{align}
 \Prob(\hbox{a hold exit, correct capture, good clock})
 &\leq\frac{C}{2\mu}\int_{1.28}^{3}
          \int|\Psi|\,|\mathcal A\Psi|\dd q\dd\tau\notag\\
 &\leq\frac{C}{2\mu}\frac{43}{25}A_*
  <3.271570\times10^{-9}.
 \label{hold:wholehold}
\end{align}
All hidden-coordinate moduli are taken before integration.  The
variation counts recrossings and every time of the interval; an
endpoint estimate is not being substituted for a path estimate.
The global annihilator norm bounds the restricted current, so no
unfinished part of the comparison wave has been removed.

\subsection{The loaded radial exit and the evaluated sum}
\label{hold:exit}
For the radial barrier increasing from zero at $0.099\,\mathrm m$
to one at $0.100\,\mathrm m$, the slope is
$1000\,\mathrm m^{-1}$.  Let
\[
 d_j^2=\frac{(2j+1)!!}{2^j100^{2j}},\quad
 B_j=(d_j^2+2hd_jd_{j+1})^{1/2},\quad
 T_G=\sqrt{22/7^{49}},\quad
 \beta=\frac{B_2}{2m}+\frac{hB_1}{2}+VT_G+F_0B_0.
\]
Here $d_j$ are the half-line Gaussian derivative norms, equivalently
the norms of its unitary odd extension.  For an integrable absolutely
continuous $f$ the grid sum obeys
\[
 \sup_{0\leq x<h}\left|h\sum_{k\in\mathbb Z}f(x+kh)
                      -\int_{\mathbb R}f\right|
       \leq h\int_{\mathbb R}|f'|.
\]
This follows by comparing the value at the chosen point of each
cell with its integral and then summing the fundamental theorem of
calculus bounds.  With $f=|\chi_{100}^{(j)}|^2$,
$\int|f'|\leq2d_jd_{j+1}$ proves the folded bound $B_j^2$.
The periodic reference has mass $h/2$ per cell and circle norm one.
Its kinetic energy starts at zero and increases by at most $V$ under
the nondecreasing gate, proving
$\|\partial_ru_{{\rm ref},i}^{\gamma}\|_{L^2(0,2)}\leq mh/2$.
Consequently the norms of
$\sqrt2\chi_{100}^{(j)}u_{{\rm ref},i}^{\gamma}$ and
$\sqrt2\chi_{100}'\partial_ru_{{\rm ref},i}^{\gamma}$
are bounded by $B_j$ and $B_1mh/2$, respectively.
For the auxiliary $\sqrt2\chi_{100}u_{{\rm ref},i}^{\gamma}$,
the product-rule residual consists respectively of the second
Gaussian derivative, the cross derivative, the outer cutoff tail
and the weak trap, and their sum is exactly $\beta$ above.
For the decreasing Gaussian density beyond $990$, periodic cell
masses give a folded tail at most
$\operatorname{Tail}(990)+h\rho(990)$.
Writing $z=9.9$, integration by parts gives
\[
 \int_z^\infty y^2e^{-y^2}\dd y
 \leq\left(\frac z2+\frac{1}{4z}\right)e^{-z^2}.
\]
The radial density is $4y^2e^{-y^2}/\sqrt\pi$ in this scaled
coordinate.  Thus $\operatorname{Tail}(990)<21/7^{49}$ and
$\rho(990)<4/7^{49}$ in detector coordinates, by $e^2>7$ and
$\sqrt\pi>1$.  Since $h<1/4$, their folded sum is below $T_G^2$.
Unitary integration for age at most $4/3$ yields the full-wave bound
\[
 \|\one_{r_{\rm phys}\geq0.099}\Psi\|
 \leq T_r:=\frac43\beta+T_G+e_*.
\]
Negative quiet ages obey the same larger bound by the quiet estimate
used in \eqref{hold:radialerror}; pointer tensor factors have norm one.

The same loaded clock estimate gives a uniform left-collar bound
$n_*=n_h+\mathcal P/(Mg)$.  Thus the exact radial momentum starts
below $P_{r,0}=\hbar\sqrt{3/2}/0.01+P_{r,h}$ and grows at most at
rate $F_r=W_r+Q_r+V_{p,r}n_*$.  There is no writer radial force.
The initial outer tail plus the absolute soft-barrier current prove
\begin{equation}\label{hold:radialexit}
 e_{\rm exit}=C\left[
 (\sqrt{21/7^{49}}+D_h)^2
 +\frac{1000T_r}{m_{\rm phys}}
       (P_{r,0}T_c+F_rT_c^2/2)\right]
 <9.127218\times10^{-8}.
\end{equation}
This bounds initial outside configurations as well as every later
exit through readout under the loaded wave.

For a directly evaluable decomposition, define
\begin{align*}
 e_{\rm base}&=C\{2D^2+(3/5)Dq_r\},\\
 e_{\rm corr}&=C\{2[(D+e_*)^2-D^2]
                   +(3/5)[(D+e_*)p_*+e_*q_r]\},\\
 e_{\rm inner}&=C\left(\sqrt{1/(14\,7^{32})}+A_*/\sqrt2\right)^2,
 &e_{\rm hold}&=43CA_*/(50\mu),\\
 e_{\rm dev}&=C\mathcal P/(10^{-5}M).
\end{align*}
The inner-tail replacement follows from
$\operatorname{erfc}(8)<e^{-64}/(8\sqrt\pi)
<1/(14\,7^{32})$.  All constants in these expressions have now
been specified by rational operations, square roots, $\pi$, and the
lens functions $D_9,Q_9$ in \eqref{old:DQ9}.  Outward evaluation gives
\begin{center}\small
\begin{tabular}{ll}\toprule
Term & Upper bound\\\midrule
$e_{\rm base}$ & $0.000537032962395581630703$\\
$e_{\rm corr}$ & $0.000010915589215098519870$\\
$e_{\rm inner}$ & $1.81324790661460\times10^{-22}$\\
$e_{\rm hold}$ & $3.271569843081749\times10^{-9}$\\
$e_{\rm core}$ & $0.000004378860028173163338$\\
$e_{\rm dev}$ & $1.06689598045912\times10^{-13}$\\
$e_{\rm exit}$ & $9.127217662854526\times10^{-8}$\\\bottomrule
\end{tabular}
\end{center}
Their sum is strictly below $0.000552421956$.
The event proof is \eqref{hold:pathineq}, the positive endpoint
bound, \eqref{hold:wholehold}, and the union with the three common
exceptional events (clock core, clock deviation, radial exit).
When combined with the earlier-label theorem, those same events
are counted once.  None of the estimates supplies an independent
internal-sector variable, a new actual-law marginal, or an additional
physical source beyond the stated effective Hamiltonian.

{\small\nocite{PortfolioControl,PortfolioReturn,PortfolioEquilibrium,PortfolioPilot,RodgersMonograph,RodgersPilot,RodgersMassive}
\begin{thebibliography}{10}

\bibitem{Bonds}
Liam Bonds, Brooke Burson, Kade Cicchella, Benjamin~H. Feintzeig, {Lynnx}, and
  Alia Yusaini.
\newblock Quantum probability via the method of arbitrary functions, 2024.
\newblock arXiv:2409.16457.

\bibitem{DGZ}
Detlef D{\"u}rr, Sheldon Goldstein, and Nino Zangh{\`i}.
\newblock Quantum equilibrium and the origin of absolute uncertainty.
\newblock {\em Journal of Statistical Physics}, 67:843--907, 1992.
\newblock DOI: \url{https://doi.org/10.1007/BF01049004}.

\bibitem{Operators}
Detlef D{\"u}rr, Sheldon Goldstein, and Nino Zangh{\`i}.
\newblock Quantum equilibrium and the role of operators as observables in
  quantum theory.
\newblock {\em Journal of Statistical Physics}, 116:959--1055, 2004.
\newblock DOI: \url{https://doi.org/10.1023/B:JOSS.0000037234.80916.d0}.

\bibitem{PortfolioEquilibrium}
Jeremy Rodgers.
\newblock Autonomous quantum measurement chains with faithful equilibrium
  records, 2026.
\newblock Author preprint. DOI: \url{https://doi.org/10.5281/zenodo.23131069}.

\bibitem{PortfolioControl}
Jeremy Rodgers.
\newblock Control consistency and born equilibrium: Uniqueness from local
  potentials and fixed interactions, 2026.
\newblock Author preprint. DOI: \url{https://doi.org/10.5281/zenodo.23131058}.

\bibitem{PortfolioPilot}
Jeremy Rodgers.
\newblock A deterministic hybrid medium for bell jump paths and autonomous
  records, 2026.
\newblock Author preprint. DOI: \url{https://doi.org/10.5281/zenodo.23131075}.

\bibitem{RodgersPilot}
Jeremy Rodgers.
\newblock A deterministic pilot medium: Bell path selection and autonomous
  material records, 2026.
\newblock Version 2; predecessor of the companion publication revision. DOI:
  \url{https://doi.org/10.5281/zenodo.22774634}.

\bibitem{RodgersMassive}
Jeremy Rodgers.
\newblock A massive configuration completion of the quantum measurement
  programme, 2026.
\newblock Version 2; predecessor of the companion publication revision. DOI:
  \url{https://doi.org/10.5281/zenodo.22774739}.

\bibitem{PortfolioReturn}
Jeremy Rodgers.
\newblock Preparation-return holonomy and equilibrium uniqueness in engineered
  interacting networks, 2026.
\newblock Author preprint. DOI: \url{https://doi.org/10.5281/zenodo.23131064}.

\bibitem{RodgersMonograph}
Jeremy Rodgers.
\newblock Shadow theory and quantum measurement: Source dynamics, event laws,
  and physical records. two constitutive completions, 2026.
\newblock Version 2, 15 September; public web edition. DOI:
  \url{https://doi.org/10.5281/zenodo.22774584}.

\bibitem{TT}
Stefan Teufel and Roderich Tumulka.
\newblock Simple proof for global existence of bohmian trajectories.
\newblock {\em Communications in Mathematical Physics}, 258:349--365, 2005.
\newblock DOI: \url{https://doi.org/10.1007/s00220-005-1302-0}.

\bibitem{Tor}
E.~Torrontegui, S.~Ib{\'a}{\~n}ez, Xi~Chen, A.~Ruschhaupt, D.~Gu{\'e}ry-Odelin,
  and J.~G. Muga.
\newblock Fast atomic transport without vibrational heating.
\newblock {\em Physical Review A}, 83:013415, 2011.
\newblock DOI: \url{https://doi.org/10.1103/PhysRevA.83.013415}.

\bibitem{ValentiniWestman}
Antony Valentini and Hans Westman.
\newblock Dynamical origin of quantum probabilities.
\newblock {\em Proceedings of the Royal Society A}, 461:253--272, 2005.
\newblock DOI: \url{https://doi.org/10.1098/rspa.2004.1394}.

\end{thebibliography}
}
\end{document}
